\documentclass[11pt]{article}
\usepackage{amsmath,amssymb,amsthm,mathtools}
\usepackage{enumitem}
\usepackage{graphicx}
\usepackage{longtable,array,booktabs}
\usepackage[margin=2.6cm]{geometry}
\usepackage[hidelinks]{hyperref}

\setlist{itemsep=1pt,topsep=3pt}
\allowdisplaybreaks
\emergencystretch=3em

% ---- the handbook's macros, verbatim (so the notebook speaks the handbook's language) ----
\newcommand{\cE}{\mathcal E}
\newcommand{\cB}{\mathcal B}
\newcommand{\cP}{\mathcal P}
\newcommand{\cQ}{\mathcal Q}
\newcommand{\cT}{\mathcal T}
\newcommand{\one}{\mathbf 1}
\newcommand{\mvec}{\mathbf m}
\newcommand{\bmu}{\boldsymbol\mu}
\newcommand{\btau}{\boldsymbol\tau}
\newcommand{\ip}[2]{\langle #1,#2\rangle}
\newcommand{\oE}{\mathbb E}
\newcommand{\sG}{\mathsf G}
\newcommand{\sZ}{\mathsf Z}
\newcommand{\sB}{\mathsf B}
\newcommand{\sT}{\mathsf T}
\newcommand{\sQ}{\mathsf Q}
\newcommand{\Gst}[4]{\Gamma_{\substack{#1\\#2}}\!\begin{bmatrix}#3\\#4\end{bmatrix}}
\newcommand{\IZ}{\mathbb I_{\mathbb Z}}
\newcommand{\IR}{\mathbb I_{\mathbb R}}
\newcommand{\Ilat}[1]{\mathbb I_{#1}}
\newcommand{\Tx}[3]{(\sT_{#1;(1)}[x]\,#2)((#3,\infty))}
\newcommand{\Th}[4]{(\sT_{#1;(1)}[x]\,#2)((#3,#4])}
\newcommand{\Ttr}[4]{(\sT_{#1;I_{#2}}[y]\,#3)(\{#4<y_1<\cdots<y_{#2}\})}
\DeclareMathOperator{\re}{Re}
\DeclareMathOperator{\im}{Im}
\DeclareMathOperator{\Reg}{Reg}
\DeclareMathOperator*{\CT}{CT}
\DeclareMathOperator{\cone}{cone}
\DeclareMathOperator{\interior}{int}
\DeclareMathOperator{\Dir}{Dir}
\DeclareMathOperator{\pat}{pat}
\DeclareMathOperator{\load}{load}
\newcommand{\reflint}[2]{\reflectbox{$#1\int$}}
\newcommand{\pseudoint}{\mathop{\mathpalette\reflint\relax}\nolimits}
\newcommand{\bsig}{\boldsymbol\sigma}
\newcommand{\bdel}{\boldsymbol\delta}
% ---- notebook-only ----
\newcommand{\status}[1]{\hfill{\small[\textit{#1}]}}
\newcommand{\hb}[2]{#1~(p.~#2)}   % printed number, page, of the handbook of 20 October 2026
\theoremstyle{plain}
\newtheorem{claim}{Claim}[section]
\theoremstyle{remark}
\newtheorem*{note}{Note}

\title{Notebook: the arithmetic of the output, read through the one Trailmoment}
\date{21 October 2026, entries 1--16}
\author{}

\begin{document}
\maketitle

\section*{What this is}

Running notes, not part of the handbook. They use the handbook's macros so that every object is written through its governing symbol, and they cite the handbook by its printed numbers and pages (the PDF compiled from the revised source of 20 October). Each claim carries a tag: \textit{derived} (inside the family's definitions, here or in the handbook), \textit{computed} (a finite computation done, with what was computed), \textit{recalled} (an outside value quoted from memory, to be verified), or \textit{not yet computed}.

\section{The one object and its slots}

The Trailmoment is
\[
(\sT_{\nu;t}[\mathbf c]\,\phi)(Y')=\int_{Y'}\sZ_{\nu;t}(\mathbf c(y))\,\phi(dy),
\]
with five slots: index $\nu$, configuration $t$, coefficient function $\mathbf c$, input $\phi$, cut region $Y'$ (Def.~\hb{5.1}{28}). A Trailmoment on a line is this object with the cut region a half-line, $\Tx{\nu}{\phi}{N}$; a head is it with the cut region $(N,N']$, $\Th{\nu}{\phi}{N}{N'}$; a trail is it with the configuration $I_d$ and the trail region, $\Ttr{\boldsymbol\nu}{d}{\phi}{N}$. These are corners of one symbol (Def.~\hb{5.6}{29}), never separate objects and never named constants. The arithmetic of the output is the study of which values the slots take along the Foreman's closings, and of what the Trailmoment's own relations---the cut relation (Prop.~\hb{5.3}{28}), transport (Thm.~\hb{5.9}{31}), the lattice relations (Thm.~\hb{5.12}{32}), the Struffle and the Shruffle (Thms.~\hb{5.25}{38}, \hb{5.26}{39}), the reflections (Thm.~\hb{5.18}{35})---impose among those values.

Two sides, one number. The native gauge of a configuration is read on the orthant side as the minimum of $\Phi_{\boldsymbol\xi}$ and on the share side as the maximum of $S$ on the real fibre, and the two agree (Thm.~\hb{4.3}{25}). Everything below is: \emph{decay is read on the orthant side, the row is read on the share side}.

\section{Chapters 31--37 stripped for parts}

Each device of the arithmetic chapters, the choice that privileges it, the verb of chapters 1--30 it instantiates, and what the choice is internally.

\begin{center}\small
\begin{tabular}{p{3.5cm}p{3.3cm}p{4.6cm}p{3.3cm}}
\toprule
device & privileged choice & the general verb & the choice, internally\\
\midrule
row, Def.~\hb{34.1}{300} & coordinates $=\Tx{k}{\IZ}{0}$ & the Laurent data of the share Gammas at the lattice points the pins allow (Props.~\hb{5.13}{32}, \hb{5.14}{33}) & unimodular pins, i.e.\ width one (\S5)\\
clearing, Lemma~\hb{35.20}{314}, Prop.~\hb{34.10}{302} & one free share, a paired word & the heads of Thm.~\hb{6.3}{43}, of a live input (Thm.~\hb{6.9}{45}); the power of the frame along a ray, Prop.~\hb{35.8}{309} (\S8) & how far the heads reach, times the denominator of the native value\\
gap, Thm.~\hb{34.3}{300} & none & keep & content-free\\
decay, Prop.~\hb{34.8}{302}, Thm.~\hb{34.19}{305} & a maximum over a cube, by hand & native value of the direction, Thm.~\hb{35.11}{310}; least Gamma rate of a chamber, Thm.~\hb{29.26}{264}; suppression by fractional cover, Thm.~\hb{32.7}{293} & the saddle\\
positivity, Thms.~\hb{32.1}{289}, \hb{32.3}{291}, \hb{32.12}{295}, \hb{33.6}{298} & particular insertions & the orthant; the completed source average, Cor.~\hb{29.27}{264} & sign of $d_i/L_i$ for linear letters $=$ width one\\
critical-line law, Def.~\hb{33.1}{297}, Thm.~\hb{33.2}{297} & two marks & the law of the free shares of an atom on its wall contour, Cor.~\hb{17.15}{111}, split by Cremona, Thms.~\hb{16.16}{99}, \hb{16.18}{99} & the fixed locus $a=1-a$ of Cremona\\
critical-line words, Def.~\hb{33.3}{298} & odd derivative of $p_{\mathrm{crit}}$ & the Cremona-odd part of the fibre law read against an insertion; Cor.~\hb{11.16}{69}; Thm.~\hb{5.17}{34} & odd part $\leftrightarrow$ odd indices\\
shape words, Def.~\hb{34.4}{301} & the cube and the seed $L$ & port words of a manufactured family, Def.~\hb{32.4}{291}, Thm.~\hb{32.6}{293} & a seed with two monomials\\
depth $n$ & one direction & rays, Def.~\hb{35.1}{307}; the centre ray, Def.~\hb{35.9}{309} & best saddle per unit clearing\\
Beta port, Def.~\hb{3.14}{22} & order $q$ on a pair shape & Thm.~\hb{3.15}{23}(ii): a $q$-th derivative moved onto the shape density by raising the shape index by $q$; moments are steps, Prop.~\hb{39.22}{345} & a jet is a step of the masses\\
$\sG_{\nu;(1,-1)}$, Cor.~\hb{17.20}{113}, Thm.~\hb{19.39}{153} & the configuration $(1,-1)$ & the one-scale kernel whose configuration is the window of levels, Thm.~\hb{7.11}{49}; Thm.~\hb{21.13}{178} at width two & width one\\
\bottomrule
\end{tabular}
\end{center}

\section{Where the collapse is built in}

Three definitions, not a theorem.

\begin{enumerate}[label=(\alph*)]
\item Def.~\hb{19.56}{158} sets $\mathcal Z$ to be the $\mathbb Q$-algebra generated by $\Tx{k}{\IZ}{0}$, $k\ge2$, before any pin has been looked at; the paragraph after Thm.~\hb{19.57}{159} then says the pins need not be unimodular and that closings land on lattice inputs with cosets and frequencies. The general content of a tree: the Trailmoments at the inputs its pins produce, with their trails, closed under the Trailmoment's own relations. The unimodular case is $\mathcal Z$.
\item Def.~\hb{34.1}{300} defines a row over fixed numbers ``such as Trailmoments of the counting input''. The general row is the one Trailmoment with the closing's rational word in its input slot, read through its cut relation (\S4).
\item Def.~\hb{33.1}{297} is the law of one share on one fibre. The general object is the law of the free shares of an atom on the wall contour of Cor.~\hb{17.15}{111}; two marks is the case where that fibre is a line.
\end{enumerate}
And $(1,-1)$ in Cor.~\hb{17.20}{113} and Thm.~\hb{19.39}{153} is width one wearing a general title (\S5).

\section{The row is a slot trade}

\subsection{The one-share fibre of $\binom{\cE}1$ on two marks}\status{derived}

Open both letters of $\binom{\cE}{1}$, $\cE=\{1,2\}$, $K=\Psi=u_1+u_2$, by Cor.~\hb{4.4}{26}: shares $a_1,a_2$ for $K$ (the share of $u_1^{-a_1}$ is the share of the basis $\{1\}$) and $b_1,b_2$ for $\Psi$, $b_1$ the share of the cobasis $u_2$ of $\{1\}$. The shape integral (Thm.~\hb{3.9}{21}) is then
\[
\int_{\Delta_{\cE}}u_1^{\,m_1-1-a_1-b_2}u_2^{\,m_2-1-a_2-b_1}\,d\omega=\Gst{2}{1}{m_1-a_1-b_2,\ m_2-a_2-b_1}{m_1+m_2-s-t},
\]
and on the balance wall the lower entry is $0$: the stacked Gamma becomes the pin $a_1+b_2=m_1$ (the mechanism after Thm.~\hb{17.12}{110}, ``the Gammas are the pins''). With $a:=a_1$, $a_2=s-a$, $b_2=t-b_1$, $b_1=a+t-m_1$:
\[
P_{\binom{\cE}1}(\mvec;s,t)=\frac1{2\pi i}\int_{(c)}\Gst21{a,\ s-a}{s}\,\Gst21{a+t-m_1,\ m_1-a}{t}\,da
=\Gst21{m_1,\ m_2}{m_1+m_2}=\sB^{\{\cE\}}_{\mvec,\,m_{\cE};\,I_{\cE}}(\one),
\]
the middle equality Cor.~\hb{4.7}{26}, the last Thm.~\hb{3.9}{21}. At $(\one;1,1)$ the integrand is $\Gst40{a,\ a,\ 1-a,\ 1-a}{}$ on $\re a=\frac12$, Lemma~\hb{5.16}{34}, and the period is one: Thm.~\hb{33.2}{297}.

\subsection{A ray on the share side}\status{derived}

Along an integral direction (Def.~\hb{35.1}{307}) each share kernel moves by the shift of Prop.~\hb{1.3}{11}(iii). On the centre ray at step $k$,
\[
\Gst21{a,\ k+1-a}{k+1}=\Gst21{a,\ 1-a}{1}\cdot\Gst11{k+1-a}{1-a}\,\Gst11{1}{k+1},
\]
the origin's kernel times a stacked Gamma that is a rational function of the share with rational coefficients: the frame of Prop.~\hb{35.8}{309}. Gammas \emph{above} the line in the step factor put zeros on the lattice, Gammas \emph{below} it put poles. An extra shape on a pair (Beukers' cube coordinates) integrates to
\[
\sB^{\{\{x,\bar x\}\}}_{(a+p,\,n+1),\,a+p+n+1;\,I}(\one)=\Gst21{a+p,\ n+1}{a+p+n+1},
\]
a factor with its Gammas below the line: that is where the share function of Def.~\hb{34.4}{301} gets its poles at $-p,\dots,-p-n$, and why the heads of a ray reach $h\sim n$. This ray itself has only zeros, and its values $\Gst21{k+1,\ k+1}{2k+2}$ are rational: no row. Rows need poles.

\subsection{The closing lands in the input slot}\status{derived}

Lemma~\hb{5.16}{34}(i), for a rational $g$ holomorphic on $\re a\ge c$ with $g=O(a^{-2})$:
\[
\frac1{2\pi i}\int_{(c)}\Gst40{a,\ a,\ 1-a,\ 1-a}{}\,g(a)\,da=\Tx{0}{(-g')\,\IZ}{0}.
\]
Index $0$, configuration $(1)$, coefficient function $x$, input the counting measure with the density $-g'$, cut region the half-line. A closing is a change of input; nothing has been evaluated. The double poles have no simple part because the word is balanced (Prop.~\hb{5.13}{32}): the index-one values that a single $\Gamma(a)^2$ would hand over cancel against those of $\Gamma(1-a)^2$.

\subsection{Chapter 6 trades the input slot for the index and cut slots}\status{derived}

With the rational word in its partial-fraction presentation (Def.~\hb{6.1}{42}), $-g'=\sum_{h,s}c_{h,s}(x+h)^{-s}$, Thm.~\hb{6.2}{42}:
\[
\Tx{0}{(-g')\,\IZ}{0}=\sum_{h,s}c_{h,s}\,\Tx{s}{\IZ}{h},
\]
each pole column becoming the index $s$ and the shift of the cut by $h$ (transport under $x\mapsto x+h$, Thm.~\hb{5.9}{31}). The cut relation Prop.~\hb{5.3}{28}, additivity in the cut slot, then brings the cuts together (Thm.~\hb{6.3}{43}):
\[
\Tx{s}{\IZ}{h}=\Tx{s}{\IZ}{0}-\Th{s}{\IZ}{0}{h},
\]
so that
\[
\Tx{0}{(-g')\,\IZ}{0}=\sum_s\Big(\sum_hc_{h,s}\Big)\Tx{s}{\IZ}{0}\;-\;\sum_{h,s}c_{h,s}\,\Th{s}{\IZ}{0}{h}.
\]
This display is Def.~\hb{34.1}{300} said by the object. The ``coordinates'' are the Trailmoment beyond the common cut at the indices $s$ whose column sum $C_s=\sum_hc_{h,s}$ (Thm.~\hb{6.4}{43}) is nonzero; the ``rational part'' is the same Trailmoment on finite cut regions, $\Th{s}{\IZ}{0}{h}=\sum_{i=1}^hi^{-s}$; the ``clearing'' is how far those cut regions reach, $d_h^{\,s}$ (Lemma~\hb{34.2}{300}, Prop.~\hb{34.10}{302}); and Lemma~\hb{35.20}{314} is ``the index slot never exceeds the order of the pole column''. Along a ray the heads reach $h\sim n$ by \S4.2, hence $d_{n+D}^{\,w}$ in Thm.~\hb{34.5}{301}.

\begin{note}
On the share side the Pochhammer zeros $(1-a)_n^2$ of the share function kill $g'(m)$ for $m\le n$: depth moves the cut of the Trailmoments out. The share side sees the coefficients grow at the sign-flipped native gauge (Thm.~\hb{35.12}{310}) and cannot see the cancellation that makes the value small; the orthant side sees the value small at the saddle (Thm.~\hb{35.11}{310}) and cannot see a lattice. That is \S1.
\end{note}

\section{Width one is one privilege}

\subsection{Pins of lists are unimodular}\status{computed; retracts an earlier claim}

For the atom formula of a list (Cor.~\hb{17.15}{111}) the pins are the share sums and the loads,
\[
\sum_Ta_T=s,\qquad\sum_Tb_T=t,\qquad\sum_{T\ni e}a_T+\sum_{T\not\ni e}b_T=m_e\quad(e\neq g),
\]
a $0/1$ matrix with $2|\cB|$ columns and $n+1$ rows. Its elementary divisors were computed for $\binom{\{1,2\}}1$, $\binom{\{1,2,3\}}1$, $\binom{\{1,2,3\}}2$, $\binom{\{1,\dots,4\}}2$, $\binom{\{1,\dots,4\}}3$, $\binom{\{1,\dots,5\}}2$: all equal to $1$. So the claim that $\binom{\{1,\dots,4\}}2$ has a pin with elementary divisor $2$ was wrong; at integer data the full fibre of the atom formula of these lists is an integral affine lattice, and closings of it land on the counting input. Not derived in general.

\subsection{The increments come from levels}\status{derived}

The scale of a bunch against its letters (Thm.~\hb{15.2}{87}, Prop.~\hb{15.1}{87}) is
\[
\int_0^\infty\rho^{\kappa_B-1}\prod_iL_i(\rho)^{-w_i}\,d\rho,\qquad L_i(\rho)=\sum_{l=0}^{c_i}\rho^{\,l}L^{(l)}_i,
\]
each letter a polynomial in the scale, its exponents the levels of its window (Thm.~\hb{7.11}{49}). By the letter presentation (Prop.~\hb{3.6}{20}) the levels $0,\dots,c_i$ are the configuration of the one-scale Gramma, and the share side of a Gramma is pinned by its configuration (Prop.~\hb{4.2}{25}) with the factor $|\det t|^{-1}$ of Cor.~\hb{2.6}{14}: the scale pin is
\[
\sum_i\sum_l l\,b_{i,l}=\kappa_B,
\]
with coefficients $0,1,\dots,c$. Thm.~\hb{21.13}{178} writes it at $c=2$ for one letter: pins $b_0+b_1+b_2=t$, $b_1+2b_2=\kappa_B$, kernel
\[
\tfrac12\,\Gst31{z,\ \tfrac{\kappa_B-z}2,\ \tfrac{2t-\kappa_B-z}2}{t},\qquad z=b_1,
\]
and a closing in $z$ at $z=-j$ puts the other two entries on $\frac12+\mathbb Z$ for half the $j$. Their Laurent data (Prop.~\hb{5.13}{32}) are Trailmoments of the lattice input $\Ilat{\frac12+\mathbb Z}$ (Def.~\hb{5.2}{28}), a frequency on the counting input by Thm.~\hb{5.12}{32}(ii). Hence:

\begin{claim}\label{cl:width}
The increment of the coordinates away from the integers is $1/c$ for a bunch of width $c$ read on its scale. At width one the scale pin has coefficients $0,1$, every closing lands on integers, the one-scale kernel is $\sG_{\nu;(1,-1)}$, and the coordinates are $\Tx{k}{\IZ}{N}$. The counting input and the configuration $(1,-1)$ are one privilege, width one: the first window, and nothing else. \status{derived}
\end{claim}

The $2^{2j+1}-1$ of Thm.~\hb{33.2}{297}(ii) is a different source: Cremona's fixed point $a=\frac12$ on the line, not a level.

\section{The centre ray of $\binom{\cE}1$ on $n$ marks}

\subsection{Closed form}\status{derived}

Cor.~\hb{17.20}{113} for the circuit $\binom{\cE}{n-1}$ with unit weights, masses on $m_{\cE}=(n-1)s+t$, $\nu_e=m_e-s$:
\[
P_{\binom{\cE}{n-1}}(\mvec;s,t)=\frac1{\Gamma(s)\Gamma(t)}\int_0^\infty y^{s-1}\prod_{e\in\cE}\sG_{\nu_e;(1,-1)}(1,y)\,dy .
\]
Prop.~\hb{8.14}{57}: the dual tables with the same exponents give the interest at the exponents swapped, so $P_{\binom{\cE}1}(\mvec;s,t)=P_{\binom{\cE}{n-1}}(\mvec;t,s)$. At the centre $s=t=\sigma$, $\mvec=\sigma\one$, every $\nu_e=0$:
\[
P_{\binom{\cE}1}(\sigma\one;\sigma,\sigma)=\frac1{\Gamma(\sigma)^2}\int_0^\infty y^{\sigma-1}\,\sG_{0;(1,-1)}(1,y)^n\,dy,\qquad \sG_{0;(1,-1)}(1,y)=\int_0^\infty u^{-1}e^{-u-y/u}\,du .
\]

\subsection{Check at two marks}\status{derived}

At $n=2$, $\sigma=1$, the integral over $y$ is read first, then scale and shape (Lemma~\hb{2.7}{14}) on $(u,v)$:
\[
\int_0^\infty\sG_{0;(1,-1)}(1,y)^2\,dy=\iint\frac{e^{-u-v}}{uv}\Big(\int_0^\infty e^{-y(\frac1u+\frac1v)}dy\Big)du\,dv=\iint\frac{e^{-u-v}}{u+v}\,du\,dv=\int_0^\infty e^{-\rho}\,d\rho=1,
\]
with $u+v=\rho$, $u=\rho w$, $du\,dv=\rho\,d\rho\,dw$: the period one of Thm.~\hb{33.2}{297}.

\subsection{What the kit says about the values}

\begin{itemize}
\item Native value of the direction $(\one;1,1)$: the stride is $f=\alpha^{\one}/(K\Psi)$ (Def.~\hb{35.9}{309}), the native gauge is $\one$, and $f(\one)=1/(K(\one)\Psi(\one))=1/n^2$; the values go like $n^{-2k}$ times the power $k^{-(n-1)/2}$ of Prop.~\hb{17.8}{109}. \status{computed}
\item The integration count of a period on $n$ marks is $n-1$ (Def.~\hb{19.56}{158}), so the Trailmoment part has index at most $n-1$; the lattice sum may leave the algebra (remark after Thm.~\hb{19.57}{159}). \status{derived}
\item Values at $\sigma=1$. At $n=3$ the value is $3\,\Tx{2}{(\Ilat{1+3\mathbb Z}-\Ilat{2+3\mathbb Z})}{0}$, a Trailmoment at index $2$ of a lattice input with the frequency $\frac13$ (\S8.5; the earlier recollection of $\Gamma(\frac13)^6/\pi$ belongs to another exponent). At $n=4$ the value is $8\,\Tx{3}{\Ilat{1+2\mathbb Z}}{0}=7\,\Tx3{\IZ}0$, the index bound attained, derived in \S8.6; along the whole centre ray the rows are $A_k\Tx3{\IZ}0-B_k$ with rational $A_k,B_k$ (\S8.7). At $n=5$ an $L$-value of a modular form. \status{computed at $n=3$; derived at $n=4$; recalled at $n=5$}
\item The clearing of the rows along the centre rays at $n=3$ and $n=4$: \S8.7. Both raw centre rays lose the gap, and not by the heads alone: the clearing carries the $k$-th power of the denominator of the native value, $9^k$ and $64^k$, before the heads $d_k^2$ and $d_k^3$ are counted. The cube insertion of Def.~\hb{34.4}{301} wins because its native value $(\sqrt2-1)^4$ is a unit (\S8.8). \status{computed}
\end{itemize}

\section{Purity, internally}

\subsection{What purity is}\status{derived}

By \S4.4 the coordinates of a closing are the Trailmoments at the indices whose column sums are nonzero. ``Pure at index $w$'' is $C_s=0$ for $s\neq w$. Which columns a word has is fixed before cancellation: the order of a pole at a lattice point is the number of Gammas above the line whose arguments fall into it, and the Laurent data at a pole of order $w$ (Prop.~\hb{5.14}{33}) carry every index up to $w$, with products---which are trails by the Struffle (Thm.~\hb{5.25}{38})---at coincident poles. Purity is $(\text{populated columns})-1$ linear conditions on the word.

\subsection{The two universal cancellations}\status{derived}

\begin{enumerate}[label=(\roman*)]
\item \emph{Balance.} Every share kernel is a balanced stacked Gamma because the Zamma divided out the Gamma of the index (Def.~\hb{2.12}{15}, Cor.~\hb{4.4}{26}), and at lattice data the regularized index-one values of a balanced word cancel (Prop.~\hb{5.13}{32}). On the input side: $C_1=0$ exactly when $R=O(x^{-2})$ (Thm.~\hb{6.4}{43}). Column $1$ is empty.
\item \emph{Cremona.} An exchange of the family (Def.~\hb{11.14}{69}) acts on the share side as the reflection of the lattice index. Thm.~\hb{6.6}{44}(iv): if $R(-L-x)=-R(x)$ then $C_s=0$ at every even $s$. The reason sits in the kernel slot: $\sZ_{n;(1)}(x)=x^{-n}$ has the parity of $n$ (after Thm.~\hb{5.17}{34}), so a word of one parity under the reflection carries one parity of index. Cor.~\hb{11.16}{69} is the same fact on the table.
\end{enumerate}

\subsection{The count}\status{derived}

For a word odd under Cremona with pole order $w$, after (i) and (ii) the populated columns are the odd indices $3,5,\dots$ up to $w$, together with the trails of odd total index with no part equal to $1$ at coincident poles:
\[
w=3:\ \text{index }3\text{ alone};\qquad w=5:\ \text{index }5\text{ and the trail }(2,3);\qquad w=7:\ \text{four}.
\]
Purity is free at pole order $3$ and nowhere else. The two-shape word of Thm.~\hb{34.5}{301} is easy for exactly this reason: one free share (no trails), pole order $3$ from the two shapes, and the family's two cancellations empty everything below. The alias $(1,-1)$ plays no part.

\begin{claim}
There is no third universal cancellation on a one-dimensional lattice: two reflections compose to a translation, and a rational word invariant under a translation is constant. Beyond pole order $3$ every further cancellation is a condition on the word---Prop.~\hb{34.18}{304}'s filters are elements of a relation block (Def.~\hb{36.7}{316}) chosen to vanish on the lower columns---and it must hold on the share side while positivity holds on the orthant side. \status{derived}
\end{claim}

\subsection{The family's way around purity}\status{derived from the handbook's statements}

A ray is a sequence of rows with a recurrence of bounded order (Thm.~\hb{35.3}{308}, Prop.~\hb{35.8}{309}), so consecutive rows come for free; Cor.~\hb{36.14}{318} turns $s+1$ consecutive rows with an invertible square matrix into linear independence over $\mathbb Q$ of all $s$ survivors at once, Thm.~\hb{36.12}{317} supplying the primitive separator. No column has to be emptied; the hypothesis moves to decay against the clearing of a determinant, still read on the two sides.

\subsection{Where symmetry can still do more}\status{not yet computed}

The bijections carrying or exchanging a table (Def.~\hb{11.14}{69}) act on the shares of its atom formula by permutation; the surviving coordinates of a multi-share closing are the invariants of that action on the multi-index columns. The lists $\binom{\cE}r$ are carried by every permutation of $\cE$, and an atom is the realization at which every mark reads alike (Def.~\hb{10.1}{64}). The one data point: the centre of $\binom{\cE}1$ on four marks is pure at the only weight-$3$ coordinate (\S6.3). Whether the action of the carrying bijections of $\binom{\cE}1$ on six marks kills the trail $(2,3)$ at weight $5$ is a finite computation on the lattice of one closing.

\section{The row of a closing with several free shares}

The question: Prop.~\hb{34.10}{302} and Lemma~\hb{35.20}{314} compute the row and the clearing when one free share is closed; what happens with several. The answer, in the slot language of \S1: closing a second share does not change the index or the cut slot of the Trailmoment that the closing produces, it changes the \emph{input} slot, from the counting input to a live input, the values of the smaller word that the closing leaves behind, read along the ray that the new pin traces. Everything after that is the trade of \S4.4 again, with the live input in place of $\IZ$ (Chapter~6, ``live inputs''). The new thing is not the row but the denominator along a ray: it is governed by the frame of the ray's recurrence (Prop.~\hb{35.8}{309}), which is a power of the product of the native values, and that power is what kills the raw centre rays.

\subsection{The word on the fibre}\status{derived}

At integer data with unit weights the atom formula (Cor.~\hb{17.15}{111}) is the fibre integral
\[
P_F(\mvec;\mathbf s,\mathbf t)=\frac1{\prod_j\Gamma(s_j)\Gamma(t_j)}\ \frac1{(2\pi i)^D}\int_{\mathbf c+i\mathbb R^D}\ \prod_{\alpha}\Gamma\big(\ell_\alpha(\mathbf z)\big)\,d\mathbf z,
\]
$\mathbf z$ integral coordinates on the fibre of the pins (\S5.1), each share $\ell_\alpha$ an affine function of $\mathbf z$ with integer coefficients, $D=\sum_j2(|\cB_j|-1)-(n-1)$ with both letters active. The word is balanced, $\sum_\alpha\ell_\alpha=\sum_j(s_j+t_j)$. Along a ray (Def.~\hb{35.1}{307}) each share moves by an integer $u_\alpha$ per step, and the shift (Prop.~\hb{1.3}{11}(iii)) writes the word at step $k$ as the word at the origin times a rational function of $\mathbf z$,
\[
\mathcal W_k(\mathbf z)=\mathcal W_0(\mathbf z)\,g_k(\mathbf z),\qquad g_k(\mathbf z)=\frac{\prod_\alpha(\ell_\alpha(\mathbf z))_{ku_\alpha}}{\prod_j(s_j)_{ks_{1j}}(t_j)_{kt_{1j}}},\qquad(\ell)_{-r}=\frac1{(\ell-r)_r},
\]
whose poles lie on the hyperplanes $\ell_\alpha=1,\dots,k|u_\alpha|$ of the shares that fall: their reach grows with $k$. This is \S4.2 in $D$ variables.

\subsection{One free share: the privilege is the pairing, not the number}\status{derived}

Close one share on the poles $\ell_1=-j$ of a Gamma whose argument falls (Thm.~\hb{4.5}{26}). At $D=1$ the residue is the rest of the word at the point: the Laurent data of the Gammas polar there (Prop.~\hb{5.14}{33}), the derivatives of $g_k$, and the constant
\[
c_j=\prod_{\alpha\notin P(j)}\Gamma\big(\ell_\alpha(-j)\big)\prod_{\alpha\in P(j)}\frac{(-1)^{j_\alpha}}{j_\alpha!},
\]
$P(j)$ the Gammas with nonpositive argument at the point: a balanced stacked Gamma at integers. It is a rational function of $j$ exactly when the Gammas above pair off by reflection, $\ell_\beta=1-\ell_\alpha$, for then $\Gamma(\ell_\alpha)\Gamma(\ell_\beta)=\pi/\sin\pi\ell_\alpha$ and the word is periodic along the lattice: its Laurent data are the same at every lattice point, the terms of the closing are rational in $j$, and Prop.~\hb{34.10}{302}, Lemma~\hb{35.20}{314} and Chapter~6 give the row over the counting input with the clearing $d^{\,w}_{\mathrm{reach}}$ (\S4.4). The two-mark word $\Gst40{a,\ a,\ 1-a,\ 1-a}{}$ is paired by its two load pins $a_e+b_{e'}=m_e=1$; the shape words of Thm.~\hb{34.5}{301} are that word times Beta corners rational in $a$. That is the whole privilege of ``one free share'': the one-share word of two marks is a product of reflection pairs. Without the pairing one share already gives a Gamma ratio, $c_j$ a sequence of tier one (Prop.~\hb{35.8}{309}), and the closing is a Trailmoment of a live input (Thms.~\hb{6.8}{45}, \hb{6.10}{46}); \S8.5 is such a case.

\subsection{Several free shares: the terms are the child along its ray}\status{derived}

With $D\ge2$ the residue at $\ell_1=-j$ is a fibre integral of dimension $D-1$ with one more pin (Thm.~\hb{4.5}{26}), the \emph{child}. On the child every share is affine in $j$: the pin moves the child's data along an integral direction, a ray of the child in the sense of Def.~\hb{35.1}{307}. So the terms of the closing are
\[
\rho(j)=f(j),\qquad f(j)=\text{the value of the child at the $j$-th point of its ray},
\]
and by Thm.~\hb{35.19}{313}(ii) $f$ satisfies a recurrence $\sum_{i=0}^{\tau}p_i(j)f(j+i)=0$, $p_i\in\mathbb Q[j]$, of order $\tau$ at most the dimension of the child's ray: a live input of finite tier. $D=1$ is tier one; the paired word is $f$ rational. Three things follow.

\begin{enumerate}[label=(\alph*)]
\item \emph{Convergence is the chamber.} The terms grow or fall like the root rate $r_G=\bar d^{\,\bar d}/\prod d_l^{\,d_l}$ of the child's Gamma scalar (Thm.~\hb{29.18}{260}, Cor.~\hb{29.19}{261}), and the closing exists only where the chamber rate is below one (Thm.~\hb{29.26}{264}, Cor.~\hb{29.24}{263}). Thm.~\hb{4.5}{26}'s hypothesis, that the shifted contour's contribution vanishes, is this. On the centre ray of $\binom{\cE}1$ on four marks the one-share closing exists at $k=1$ and at no later step (\S8.6); on three marks it exists at every step (\S8.5).
\item \emph{Coincident poles bring heads whose reach moves with the point.} Where several Gammas are polar on the crossed hyperplane (Thm.~\hb{29.29}{265}) the residue takes Laurent data. With $\varepsilon_\alpha=\ell_\alpha(\mathbf v)$ the deviations, $p_\alpha$ the positive and $-j_\alpha$ the nonpositive arguments at the point, Props.~\hb{5.13}{32}, \hb{5.14}{33} and the cut relation $\Tx{k}{\IZ}{N}=\Tx{k}{\IZ}{0}-\Th{k}{\IZ}{0}{N}$ give, for the Gammas above the line (sign reversed below), the degree-$k$ part of the logarithm of the word over its principal part:
\begin{multline*}
\frac1k\sum_{\alpha}(-\varepsilon_\alpha)^k\,\Tx{k}{\IZ}{0}\;-\;\frac1k\sum_{\alpha\notin P}(-\varepsilon_\alpha)^k\,\Th{k}{\IZ}{0}{p_\alpha-1}\\
+\;\frac1k\sum_{\alpha\in P}\varepsilon_\alpha^{\,k}\,\Th{k}{\IZ}{0}{j_\alpha}\qquad(k\ge2),
\end{multline*}
and at $k=1$ the regularized value drops out by balance, $\sum_\alpha\varepsilon_\alpha=0$, leaving the heads. The part at the cut $0$ is the same polynomial in $\mathbf v$ at every lattice point; only the heads move, with reach $p_\alpha-1$ or $j_\alpha$, affine in $j$. Summed against the live input they are trails with the live input outermost and the inner cut affine in the outer point: at three marks, $\sum_jf(j)\,r_k(j)\,(\sT_{1;(1,1)}[(j+k-1,\,i)]\,\IZ)((0,j+k])$ (\S8.5). Products of heads at a pole of higher order are products of Trailmoments on a line, hence trails by the Struffle (Thm.~\hb{5.25}{38}).
\item \emph{The row is over the Trailmoments of the live input}: \S8.4.
\end{enumerate}

\subsection{The row of a closing against a live input}\status{derived; the exponent of the clearing as in Thm.~6.9(iv)}

\begin{claim}\label{cl:liverow}
Let $x=\sum_{j>N}r(j)\,f(j)$ with $r(j)=\sum_{h,s}c_{h,s}(j+h)^{-s}$ a rational word (Def.~\hb{6.1}{42}) whose poles lie in finitely many cosets of $\mathbb Z$, of largest order $w$, and $f$ a live input with $\sum_{i=0}^\tau p_i(j)f(j+i)=0$, $p_i\in\mathbb Z[j]$ of degree at most $d$, $p_0$ and $p_\tau$ without zeros on the lattice beyond $N$.
\begin{enumerate}[label=(\roman*)]
\item $x=\sum_{h,s}c_{h,s}\,(\sT_{s;(1)}[x+h]\,f\,\IZ)((N,\infty))$: the row is over the Trailmoments of the live input, one coordinate per coset, index and shift (Thm.~\hb{6.2}{42} with the coefficient function $x+h$; Thm.~\hb{6.9}{45}(i)).
\item The shifts reduce: reading the recurrence against $(j+h)^{-s}$ and moving the cuts (Prop.~\hb{5.3}{28}) writes every coordinate with a shift beyond $\tau$ in its coset as a $\mathbb Q$-combination of coordinates with shifts below $\tau$, of lower indices, of the values $f(N+1),\dots,f(N+\tau)$, and, when $s\le d$, of the moments $(\sT_{-e;(1)}[x]\,f\,\IZ)((N,\infty))$, $e\le d-s$, Trailmoments at nonpositive index, the form of Thm.~\hb{29.18}{260} (Thm.~\hb{6.9}{45}(ii) is the Breta input of Thm.~\hb{6.8}{45}).
\item Each move of a cut releases a head on a coset $\frac1u\mathbb Z$, cleared by $u^{\,w}d_{uJ}^{\,w}$ (Lemma~\hb{34.2}{300}), and each use of the recurrence divides by a value of $p_0$ or $p_\tau$: through the reach $J$ the row is cleared by $d_{J}^{\,w}$ times the $w$-th power of the product of the extreme coefficients over the shifts used, the shape of Thm.~\hb{6.9}{45}(iv). At $f\equiv1$, $\tau=0$, this is Prop.~\hb{34.10}{302}.
\end{enumerate}
\end{claim}

Along a ray the reach is $J\sim k\max_\alpha|u_\alpha|$, so the heads cost $d_{ck}^{\,w}$ as in Thm.~\hb{34.5}{301}; the new cost is the product of the extreme coefficients, which is the frame of the live input's recurrence (Prop.~\hb{35.8}{309}): a power times a stacked Gamma. The power is the part that cannot cancel against anything, and \S8.7 measures it.

\subsection{Three marks: a live input of tier one, and a frequency $\frac13$}\status{derived; the value computed}

On the centre ray of $\binom{\cE}1$ on three marks (\S6.1) the pins $a_e=b_e$ leave the word $\Gamma(k)^{-2}\prod_{e}\Gamma(a_e)^2$ on $a_1+a_2+a_3=k$, $D=2$. The share $a_1$ is not closed but merged, by Cor.~\hb{4.7}{26} with the shares $(a_1,\,k-a_2-a_1)$: $\frac1{2\pi i}\int\Gamma(a_1)^2\Gamma(k-a_2-a_1)^2\,da_1=\Gst41{c,\,c,\,c,\,c}{2c}$, $c=k-a_2$. One share is left, with the word
\[
P_k=\frac1{\Gamma(k)^2}\ \frac1{2\pi i}\int_{(c)}\Gamma(a)^2\,\frac{\Gamma(k-a)^4}{\Gamma(2k-2a)}\,da,
\]
not paired: the Gamma below has the rate $2$ in the share. Closing to the left on the double poles $a=-j$ of $\Gamma(a)^2$ converges at every $k$ (the word is $O(4^{\re a}|a|^{2k-3})$ on the shifted contours), and the residue (Prop.~\hb{5.14}{33} at $p=2$) is
\begin{multline*}
P_k=\frac{2}{\Gamma(k)^2}\sum_{j\ge0}\frac{\Gamma(k+j)^4}{(j!)^2\,\Gamma(2k+2j)}\Big[\Th{1}{\IZ}{0}{j}\\
-2\,\Th{1}{\IZ}{0}{k+j-1}+\Th{1}{\IZ}{0}{2k+2j-1}\Big],
\end{multline*}
the regularized index-one values cancelling by balance, $2-4+2=0$. The live input is the Gamma ratio $f(j)=(j!)^2/(2j+1)!$, of tier one with $f(j+1)/f(j)=\frac{(j+1)^2}{(2j+2)(2j+3)}$, root rate $\frac14$; the rest is $2f(j)\,r_k(j)$ with the rational function $r_k(j)=(j+1)_{k-1}^4\big/\big((k-1)!^2\,(2j+2)_{2k-2}\big)$, whose poles sit on the coset $\frac12+\mathbb Z$ (rate $2$), and the bracket is a head of reach $k-1$ in $j$ less a head of reach $j+k$, the trail of \S8.3(b). At $k=1$,
\[
P_1=2\sum_{j\ge0}\frac{(j!)^2}{(2j+1)!}\,\Th{1}{\IZ}{j}{2j+1}=3\,\Tx{2}{(\Ilat{1+3\mathbb Z}-\Ilat{2+3\mathbb Z})}{0}=2.3439072\ldots,
\]
the sum and the integral of \S6.1 agreeing to $40$ digits, and the identification of the Trailmoment found by PSLQ with the basis $(P_1,\,\Tx{2}{(\Ilat{1+3\mathbb Z}-\Ilat{2+3\mathbb Z})}{0},\,1)$: a Trailmoment at index $2$ of a lattice input with the frequency $\frac13$, the input the remark after Thm.~\hb{19.57}{159} said a lattice sum can make and no bunch or leaf shows. Its irrationality is not known. Two remarks. The index-$0$ Trailmoment of the live input is already a lattice Trailmoment one index up, $\sum_jf(j)=2\,\Tx{1}{(\Ilat{1+3\mathbb Z}-\Ilat{2+3\mathbb Z})}{0}$ (computed), while $\sum_jf(j)/(2j+1)$ is not a rational combination of the index-$2$ value and $1$ (computed, $40$ digits): the single coordinates of Claim~\ref{cl:liverow} are more numbers than the closing's combination of them, and the family's word is a combination that is pure. And the frequency $\frac13$ is not a level: $3$ is the number of marks, and $K(\alpha)=\alpha_1+\alpha_2+\alpha_3$ vanishes at $\alpha=(1,\omega,\omega^2)$, $\omega^3=1$, the gauges removed from the letter complement $X_F$ (Def.~\hb{35.1}{307}); at two and four marks, where $K$ vanishes at a sign, the frequency is $\frac12$ (\S8.6). Not derived.

\subsection{Four marks: the pairs merge, the coset $\frac12$ gives $7\,\Tx3{\IZ}0$}\status{derived}

On four marks the word is $\Gamma(k)^{-2}\prod_{e=1}^4\Gamma(a_e)^2$ on $\sum_ea_e=k$, $D=3$. Cor.~\hb{4.7}{26} merges the pairs $\{1,2\}$ and $\{3,4\}$ exactly, $c=a_1+a_2$:
\[
P_k=\frac1{\Gamma(k)^2}\ \frac1{2\pi i}\int_{(c)}\frac{\Gamma(c)^4}{\Gamma(2c)}\ \frac{\Gamma(k-c)^4}{\Gamma(2k-2c)}\,dc,
\]
the shape integrals $\Gst21{c,\ c}{2c}$ of the two pairs (Thm.~\hb{3.9}{21}) times $\Gamma(c)^2\Gamma(k-c)^2$, the two-mark word on the pairs. At $k=1$ the reflections (Prop.~\hb{1.3}{11}) turn it into $\frac{2\pi^3\cos\pi c}{(1-2c)\sin^3\pi c}$: triple poles at the integers with periodic Laurent data, $\frac{\cos\pi\varepsilon}{\sin^3\pi\varepsilon}=\frac1{\pi^3\varepsilon^3}\big(1+0\cdot\varepsilon^2+O(\varepsilon^4)\big)$, and the one rational factor $(1-2c)^{-1}$, whose pole is on the coset $\frac12+\mathbb Z$ because the mass pin of a pair at equal shares reads $2c$ (Claim~\ref{cl:width} in the flesh). Closing to the right on $c=m\ge1$,
\begin{multline*}
P_1=-\sum_{m\ge1}\partial_c^2\frac{1}{1-2c}\Big|_{c=m}=8\sum_{m\ge1}\frac1{(2m-1)^3}\\
=8\,\Tx{3}{\Ilat{1+2\mathbb Z}}{0}=\big(2^3-1\big)\Tx{3}{\IZ}{0}=7\,\Tx3{\IZ}0,
\end{multline*}
the last step the coset relation Thm.~\hb{5.12}{32}(i) at $q=2$: the recalled value of \S6.3 is derived, and the $7$ is $2^3-1$ for the same reason as the $2^{2j+1}-1$ of Thm.~\hb{33.2}{297}. At $k\ge2$ the step factor $g_k(c)=(1-c)_{k-1}^4\big/\big((2-2c)_{2k-2}\Gamma(k)^2\big)$ has degree $2k-2$: the shifted contours do not vanish and the one-share closing diverges in both directions, the chamber of \S8.3(a) closed. The rows along the ray come from the ray's own recurrence, \S8.7.

\subsection{The frame is the denominator along a ray}\status{derived at the power $1/h$; computed in full on two rays}

Along a ray the values satisfy $\sum_{i=0}^hp_i(k)x_{k+i}=0$ (Thm.~\hb{35.19}{313}(ii)). When the rows are $x_k=\sum_iA_{k,i}v_i-B_k$ with $1,v_1,\dots$ independent over $\mathbb Q$, the recurrence separates by coordinates and $A_{\cdot,i}$, $B$ are solutions. Their frame (Prop.~\hb{35.8}{309}) is
\[
C_k=C_{k_0}\big((-1)^hc\big)^{k-k_0}\ \Gst{\cdot}{\cdot}{k-a_1,\ \dots}{k-b_1,\ \dots}\Big/\Gst{\cdot}{\cdot}{k_0-a_1,\ \dots}{k_0-b_1,\ \dots},\qquad c=\frac{\operatorname{lc}(p_0)}{\operatorname{lc}(p_h)},
\]
and when $\deg p_0=\deg p_h$ the number $c$ is, up to sign, the product of the roots of the symbol (Cor.~\hb{35.6}{308}), which are native values of the ray's gauges and include every conjugate of $f(\alpha^\star)$ (Prop.~\hb{35.7}{309}). If $D_k$ clears the $h$ rows at $k,\dots,k+h-1$, then $D_k^{\,h}C_k$ is a nonzero integer, so
\[
\liminf_k\frac{\log D_k}{k}\ \ge\ \frac1h\log\frac1{|c|}.\qquad\text{(derived)}
\]
Computed on the two centre rays, the recurrences fitted to the exact rows ($k\le13$ at three marks, $k\le9$ at four, PSLQ at $50$--$60$ digits) and checked on every further row, the rows then extended exactly to $k=200$:
\begin{gather*}
9(k+1)^2x_{k+2}-(10k^2+10k+3)\,x_{k+1}+k^2x_k=0\qquad\text{(three marks)},\\
64(k+1)^3x_{k+2}-2(2k+1)(5k^2+5k+2)\,x_{k+1}+k^3x_k=0\qquad\text{(four marks)},
\end{gather*}
with the symbols $(9\lambda-1)(\lambda-1)$, $c=\frac19$, and $(16\lambda-1)(4\lambda-1)$, $c=\frac1{64}$.
The roots are exactly the native values $K(\eta)^{-2}$ of the sign classes with $K(\eta)\neq0$ (Thm.~\hb{35.12}{310}(ii)): $K(\eta)\in\{3,\pm1\}$ at three marks, $\{4,\pm2\}$ at four, the sign $(1,1,-1,-1)$ with $K=0$ absent. Cor.~\hb{35.6}{308} is confirmed to the letter on both rays. The rows: with $v=\Tx{2}{(\Ilat{1+3\mathbb Z}-\Ilat{2+3\mathbb Z})}{0}$,
\[
x_1=3v,\quad x_2=v-\tfrac23,\quad x_3=\tfrac59v-\tfrac{23}{54},\quad x_4=\tfrac{31}{81}v-\tfrac{145}{486},\quad x_5=\tfrac{71}{243}v-\tfrac{1331}{5832},\ \dots
\]
and with $v'=7\,\Tx3{\IZ}0$,
\[
x_1=v',\quad x_2=\tfrac1{16}v'-\tfrac38,\quad x_3=\tfrac{7}{1024}v'-\tfrac{27}{512},\quad x_4=\tfrac1{1024}v'-\tfrac{37}{4608},\quad x_5=\tfrac{679}{4194304}v'-\tfrac{25555}{18874368},\ \dots
\]
One coordinate each: both rows are pure (\S7), at index $2$ and at index $3$. The clearing $D_k=\operatorname{lcm}(\operatorname{den}A_k,\operatorname{den}B_k)$:
\begin{center}\small
\begin{tabular}{lcccc}
\toprule
& $3$-part at $k=200$ & rest, $\log/k$ at $k=200$ & $\log D_k/k$ & decay $-\log f(\alpha^\star)$\\
\midrule
three marks & $3^{392}\approx9^{k}$ & $2.006\approx2=$ the index & $4.16\to\log9+2$ & $\log9=2.197$\\
& $2$-part & & & \\
four marks & $2^{1179}\approx64^{k}$ & $3.003\approx3=$ the index & $7.09\to\log64+3$ & $\log16=2.773$\\
\bottomrule
\end{tabular}
\end{center}
The rest is $d_k^{\,w}$, the primes up to $k$ to the power of the index, the heads of Claim~\ref{cl:liverow}(iii). The power is the \emph{full} $|c|^{-k}$, not the derived $|c|^{-k/2}$: the $k$-th power of the product of all the native values. On three marks the frame alone exactly cancels the decay, $9^kx_k\approx C/k$, and the heads finish it; on four marks $64^k\cdot16^{-k}=4^k$ before the heads. Both raw centre rays lose the gap, and the reason is not purity.

\subsection{What the frame says about rays}\status{derived from \S8.7, modulo the full power}

Write the symbol with coprime integer coefficients, $c_h\lambda^h+\dots+c_0$; then $c_hf(\alpha^\star)$ is an algebraic integer, and $c_h$ is what the recurrence divides by at every step. On the two rays computed $D_k\approx c_h^{\,k}d_k^{\,w}$. Read against the decay $f(\alpha^\star)^k$ of Thm.~\hb{35.11}{310}, the gap (Thm.~\hb{34.3}{300}) on a ray asks for
\[
\big|c_h\,f(\alpha^\star)\big|<e^{-w\,c_{\mathrm{lcm}}},\qquad c_{\mathrm{lcm}}=1.1055\ldots\ (\text{Lemma~\hb{34.2}{300}}),
\]
an algebraic integer of modulus below $e^{-w c_{\mathrm{lcm}}}$, $w$ the index of the coordinate. Three consequences.
\begin{enumerate}[label=(\roman*)]
\item \emph{No centre ray of a list with integer weights passes.} Its native value is $f(\one)=\prod_jK_j(\one)^{-2\sigma_j}$, a rational root of the symbol, so $c_h$ is a multiple of its denominator and $|c_hf(\alpha^\star)|\ge1$. The other roots, $\prod_jK_j(\eta)^{-2\sigma_j}$ at the signs, are at most one in modulus as well: the centre ray has no gauge where the stride is large.
\item \emph{The two-shape word passes because its native value is a unit.} Its gauges are the roots of $a^2+2a-1$ (Prop.~\hb{34.8}{302}): $a=\sqrt2-1$ inside the cube with $f=(\sqrt2-1)^4$, and $a=-\sqrt2-1$ outside it with $f=(\sqrt2+1)^4$; with both gauges nondegenerate the symbol vanishes at both (Cor.~\hb{35.6}{308}) and so contains $\lambda^2-34\lambda+1$, which is monic: $c_h=c_0=1$, the frame a pure stacked Gamma, the clearing the heads $d_n^3$ alone, and $(\sqrt2-1)^4=0.0294<e^{-3c_{\mathrm{lcm}}}=0.0363$ (Thm.~\hb{34.9}{302}). The cube insertion does not change the row; it changes the native value into a unit.
\item \emph{The search criterion.} A (family, direction) pair is a candidate exactly when the native value of the direction (Thm.~\hb{8.18}{58}, Thm.~\hb{35.11}{310}) is an algebraic number $\lambda$ whose integer multiple $c_h\lambda$ is small, and the ray has a second gauge, off the real orthant, where the stride is large, the conjugate. Both are read from the stride $f$ alone, before any closing: the native gauges are the critical points of $f$ (Thm.~\hb{35.19}{313}(i)). This is the internal statement of what the cube does, and it is where the arena is: directions whose critical values are units.
\end{enumerate}

\section{The content by the pins; the coaction by the tree}

Def.~\hb{19.56}{158} fixes the content before any pin is read: $\mathcal Z$, the counting input at the cut $0$, products of single Trailmoments. \S3(a) marked that as the collapse and \S8 produced, from the pins of two lists, a coset $\frac12$, a frequency $\frac13$, a live input and a trail. This entry replaces the definition and states what is then left to prove. The company is Chapters 1--30 with Def.~\hb{36.1}{315}: the admitted words $\mathfrak W$, the relation ideal $\mathfrak I$, the evaluation. Everything motivic that the handbook has lives in $\mathfrak W/\mathfrak I$, not in $\mathbb C$.

\subsection{The content, by the pins}\status{definition; replaces Def.~19.56's $\mathcal Z$}

Let $\mathcal C$ be the $\mathbb Q$-span in $\mathfrak W/\mathfrak I$ of the admitted trails of lattice inputs with rational cosets, rational frequencies and rational cuts,
\[
\big(\sT_{\boldsymbol\nu;I_d}[y]\ e^{2\pi i\ip{\mathbf a}{y}}\,\Ilat{\mathbb L}\big)\big(\{N<y_1<\cdots<y_d\}\big),\qquad \mathbb L=\textstyle\prod_i(b_i+\mathbb Z),\ \ \mathbf a,\mathbf b\in\mathbb Q^d,\ N\in\mathbb Q,\ \boldsymbol\nu\in\mathbb N^d,
\]
including the heads (finite trail regions). It is an algebra by the Struffle (Thm.~\hb{5.25}{38}), graded by the total index, which is the integration count of Def.~\hb{19.56}{158}, and filtered by the depth $d$; its relations are the trail relations already in $\mathfrak I$: cut (Prop.~\hb{5.3}{28}), transport (Thm.~\hb{5.9}{31}), cosets and frequencies (Thm.~\hb{5.12}{32}), parity (Thm.~\hb{5.17}{34}), Struffle, Shruffle (Thm.~\hb{5.26}{39}), splitting (Thm.~\hb{5.27}{39}). $\mathcal Z$ is the corner at frequency $0$, coset $0$, cut $0$, depth one and its products. The content of a childcule, of a bunch and of a room is defined as in Def.~\hb{19.56}{158} with $\mathcal C$ for $\mathcal Z$, and the words of Thm.~\hb{19.57}{159}(iv), (v) are then read in $\mathfrak W/\mathfrak I$, the closings being generators of $\mathfrak I$ (Def.~\hb{36.1}{315}). The three values of \S8 sit in $\mathcal C$ at frequency $0$ (index $3$), at coset $\frac12$ (index $3$), and at frequency $\frac13$ (index $2$). Nothing in Chapters 1--30 is changed by this; it is the arithmetic of a tree stated with the pins read.

\subsection{The coaction on periods is the hinge}\status{derived, under the chamber; by the recurrence where the chamber closes}

Thm.~\hb{18.13}{119} at a bunch $B$ of width one with room $R$,
\[
P_F(\mvec;s,t)=\frac1{2\pi i}\int\Gst42{a_d,\ b_d,\ a_r,\ b_r}{s,\ t}\,P_{F^B}(a_d,b_d)\,P_{F^R}(a_r,b_r)\,d\sigma_d,
\]
is one coproduct term written as an integral: the restriction $\cB|B$ on one side, the contraction $\cB/B$ on the other, the share kernel between. Close the hinge share on the lattice of the bunch-side Gammas (Thm.~\hb{4.5}{26}, in the chamber of Thm.~\hb{29.26}{264}). The $j$-th term is the Laurent data $\xi_j\in\mathcal C$ of the kernel at the point (Prop.~\hb{5.14}{33}), times the bunch face and the room face at the data the pin sets, both moving along a ray (Def.~\hb{35.1}{307}). The room's values along its ray are spanned by $h$ of them, $h$ its tier (Thm.~\hb{35.19}{313}(iii)); write $P_{F^R}(j)=\sum_{i<h}c_i(j)\,P_{F^R}(j_0+i)$. Then
\[
P_F(\mvec;s,t)=\sum_{i=0}^{h-1}P_{F^R}(j_0+i)\cdot\Big[\sum_{j}c_i(j)\,\xi_j\,P_{F^B}(j)\Big],
\]
a row over $h$ values of the contraction, with coefficients that are Trailmoments of the live input ``restriction along the hinge'' (\S8.3, \S8.4); and the same with $B$ and $R$ exchanged. Read down the Foreman tree this is Thm.~\hb{19.57}{159} as a statement about rows: the conjugates of a period are finitely many values of its minors, the cofactors are elements of $\mathcal C$ when the live inputs are of tier one (\S9.5), and at the leaves, single marks, only the Laurent data remain. Where the chamber is closed the lattice sum is replaced by the ray's recurrence (Thm.~\hb{35.3}{308}), which is the same statement read on the pencil (Thm.~\hb{35.10}{309}). On the four-mark list with $B=\{1,2\}$ the room is a Beta corner, $h=1$, the one conjugate is rational and the whole value sits in the cofactor: $7\,\Tx3{\IZ}0$ (\S8.6). On the four-mark list at the centre the tree itself takes no bunch: the list is an atom (\S11.2), and $B=\{1,2\}$ is a hinge used to evaluate it.

\subsection{The coaction on the content is read on the Laplace side}\status{shape derived; the formula not yet derived}

A trail of $\mathcal C$ is, by the Shruffle (Thm.~\hb{5.26}{39}), the chain integral of a string whose letters are the Laplace transforms of its increments and the constant $1$. For a lattice input on the coset $b+\mathbb Z$ at the frequency $a$ the Laplace transform is
\[
L(U)=\sum_{m\in b+\mathbb Z}e^{2\pi iam}e^{-mU}=\frac{e^{2\pi iab}\,e^{-bU}}{1-e^{2\pi ia}e^{-U}},
\]
whose singular points on the Laplace side are the frequencies, $e^{-U}=e^{-2\pi ia}$, and the constant letter $1$ has its singular point at $U=\infty$. So the content is the algebra of chain integrals along the Laplace line with letters singular at $\infty$ and at the roots of unity the frequencies name; the two-sided reflection of Thm.~\hb{5.18}{35} reads a trail through those poles. The coaction on $\mathcal C$ is the deconcatenation of the string: keep a subsequence of the letters, the left factor is the chain of the kept letters, the right factor the product of the chains of the dropped segments, each read between the singular points of the two letters that bound it. The first half is the cut relation (Prop.~\hb{5.3}{28}) on the trail region; the second half, a chain read between two frequencies, is a Trailmoment whose two ends are both data, and it is the one verb of the content the handbook does not yet name. Writing it is the derivation. (Alias, once: Goncharov's coaction on iterated integrals on $\mathbb P^1$ minus roots of unity.)

\subsection{The one theorem}\status{hypothesis, with its falsifier}

Prop.~\hb{5.15}{33} selects the index $w$ inside a Gamma word and says it ``supplies no interchange with the sums of a closing''. That interchange is the theorem that makes the company usable outright: for a hinge closing,
\[
C_w=\operatorname{Port}_{\kappa_B}\big[\text{the word}\big]\qquad\text{(Thm.~\hb{35.17}{312}(ii))},
\]
the column sum at the top index, the coefficient of the top coordinate, equal to a port of the word at the wall of $B$: a single constant term in place of the lattice sum. If it holds, the coefficients of every row are ports, their denominators are read from the Gammas of the walls, the reach per step is the slope of the hinge kernel's poles (Lemma~\hb{18.15}{121}), and \S9.2 and \S9.3 are one coaction computed on the two sides. Falsifier, finite: on the four-mark centre ray $64^kA_k=448,\ 1792,\ 12544,\ 114688,\ 1216768,\dots$ must be the constant terms of the powers of the stride at the hinge wall, with the regularization of Def.~\hb{35.16}{312}; the torus integral of $f^k$ itself diverges at the sign $(1,1,-1,-1)$ where $K$ vanishes, so the port, not the torus, is the object to compute. For the two-shape word the same statement is that $A_n=\sum_la_l^2$ is a constant term of a power of a Laurent polynomial, which is true (recalled), and is the reason those $A_n$ are integers.

\subsection{Where the content ends: tier, not width}\status{hypothesis; three data points}

The lattice sum of a hinge is a Trailmoment of a live input of finite tier (\S8.3). Three marks: tier one, the sum is a trail at the frequency $\frac13$. Four marks: tier one, a trail at the coset $\frac12$. Five marks, hinge $\{1,2\}\,|\,\{3,4,5\}$: the room is the three-mark list along the hinge ray, tier two, and the value is the recalled modular $L$-value, which is a period of the room's pencil (Thm.~\hb{35.10}{309}) and no trail. Hypothesis: the lattice sum against a live input of tier one is in $\mathcal C$; against tier two it is a period of the pencil of the room's ray, outside $\mathcal C$. The boundary of the content is the tier of the faces along the hinge, which the tree reads before any value is computed, and width one is the case in which every face is rational.

\subsection{What the company produces of the gap's two numbers}\status{derived, correcting a claim made in discussion}

Thm.~\hb{34.3}{300} needs a decay and a clearing. Both are produced by the structure of \S9.2; the claim that they are not was wrong.
\begin{enumerate}[label=(\roman*)]
\item \emph{The decay} is a root of the symbol of the ray's recurrence, which comes from the relation space $\mathcal L$ (Thm.~\hb{35.19}{313}(ii)), a finite algebraic computation on the stride; Cor.~\hb{35.6}{308} and Prop.~\hb{35.7}{309} put $f(\alpha^\star)$ and its conjugates among the roots. The real orthant selects the root the simplex period follows (Prop.~\hb{35.7}{309}: the largest value of $f$ on $\mathbb R^{\cE}_{>0}$): the Betti cycle choosing among the de Rham side's roots.
\item \emph{The clearing} is the denominators of the lattice sum of the weight-zero parts of the cofactors $\xi_j$, which Prop.~\hb{5.13}{32} names: heads of the lattice, rational, cleared by $d_{\mathrm{reach}}^{\,w}$ (Lemma~\hb{34.2}{300}), the reach being the slope of the hinge kernel's poles (Lemma~\hb{18.15}{121}); times the faces' values at integer data; times the frame (\S8.7). Thm.~\hb{34.7}{301} is this computation for the two-shape word, and it is sharp there.
\item \emph{What is not yet derived, and where it sits.} The closed chamber of the hinge is not the end: Thm.~\hb{19.39}{153} and Cor.~\hb{19.40}{153} write the period of $\binom{\cE}1$ at every $k$ as nested lattice sums converging at the root rates $(n-1)^{-2},\dots,\frac14$, the last two marks giving $\Gst41{z,\,z,\,z,\,z}{2z}$: a Trailmoment of the Gramma live input $\Gamma(z)^4/\Gamma(2z)$ (Thm.~\hb{6.10}{46}), with the recurrence $2(2z+1)f(z+1)=z^3f(z)$. Its row and clearing have the shape of Claim~\ref{cl:liverow}: heads at the index of the coordinate, times the extreme coefficients. Computed exactly to $k=200$ on both centre rays: $\operatorname{den}(A_k)$ is a pure power of the coset prime ($2$ at four marks, $3$ at three), and the part of $\operatorname{den}(B_k)$ prime to it divides $d_k^{\,w}$, $w=3$ resp.\ $2$, equal to it prime by prime at $k=200$ except at one prime. So the apparent $(k!)^3$ of the recurrence and the odd numbers of the extreme coefficients cancel completely, exactly as \S9.4 predicts for coefficients that are ports and rational parts that are heads. Derived: no. The derivation is the interchange of \S9.4, or the summation of the Gramma input's row by the Gramma's own relations (Thm.~\hb{2.3}{13}).
\end{enumerate}
The comparison of the two numbers is not a production; it is done per family, on the produced data.

\subsection{The row of the four-mark centre ray, derived}\status{derived; checked against the identified rows for $k\le9$}

Read Chapter~33 as Thm.~\hb{18.13}{119} at the centre (an identity found in a parallel reading of the handbook): the word of \S8.6 is, after the shift, the critical-line word of Def.~\hb{33.3}{298},
\[
\langle1\rangle_k=\frac{\operatorname{crit}_1[R_k]}{\Gamma(k)^2},\qquad R_k(u)=\frac{(1-u)_{k-1}^4}{(2-2u)_{2k-2}}=Q_k(u)+\sum_{i=1}^{k-1}\frac{r_{k,i}}{u-i-\frac12},
\]
$Q_k\in\mathbb Z[\tfrac12][u]$ of degree $2k-2$ (the integer poles of the denominator cancel against the fourth-power zeros of the numerator), simple poles at the half-integers. Two reads, neither of which closes the whole word:
\begin{enumerate}[label=(\roman*)]
\item The polynomial part by Prop.~\hb{33.4}{298}, with $Q_k(u)=(1-2u)\widetilde Q_k(u)+Q_k(\frac12)$: $\operatorname{crit}_1[Q_k]=Q_k(\frac12)\,\Tx{3}{\Ilat{\frac12+\mathbb Z}}{0}+\oE[\widetilde Q_k'(\frac12+iY)]$, the moments of $\mathbb P_{\mathrm{crit}}$ rational with $d_{m+1}\oE[X^m]\in\mathbb Z$ (Thm.~\hb{33.2}{297}(iii)).
\item A simple pole at $i+\frac12$ sits on a zero of $\partial_u\Gst40{u,\,u,\,1-u,\,1-u}{}$. Closing to the right converges for a pole term, meets only the triple poles at the integers, and gives a head on the coset and nothing else:
\[
\operatorname{crit}_1\Big[\frac1{u-i-\frac12}\Big]=-\sum_{m\ge1}\partial_u^2\Big[\frac1{(1-2u)(u-i-\frac12)}\Big]_{u=m}=-\frac8i\,\big(\sT_{3;(1)}[x]\,\Ilat{1+2\mathbb Z}\big)\big((0,2i]\big).
\]
\end{enumerate}
So the coordinate comes from the polynomial part at the Cremona fixed point alone, $A_k=7\,Q_k(\frac12)/(k-1)!^2$, the port at infinity of $R_k(w)/(w-\frac12)$: the interchange of \S9.4 holds for critical-line words. The residues are explicit, $r_{k,i}/i=2^{-8(k-1)}\binom{2i}{i}^3\binom{2j}{j}^3(i!\,j!)^2$, $j=k-1-i$, with $R_k(\frac12)$ the $i=0$ term, and $Q_k(\frac12)=R_k(\frac12)+\sum_ir_{k,i}/i$ gives
\[
A_k=\frac{7}{2^{8(k-1)}}\sum_{i+j=k-1}\Big(\frac{(2i)!\,(2j)!}{i!\,j!\,(i+j)!}\Big)^2\binom{2i}{i}\binom{2j}{j},
\]
every term an integer (the inner ratio is one; alias, once: a super Catalan number). Hence $\operatorname{den}(A_k)$ is a power of $2$ for every $k$: the first half of the cancellation of \S9.6(iii), derived. The rational part,
\[
(k-1)!^2B_k=-\oE[\widetilde Q_k'(\tfrac12+iY)]+8\sum_{i=1}^{k-1}\frac{r_{k,i}}{i}\,\big(\sT_{3;(1)}[x]\,\Ilat{1+2\mathbb Z}\big)\big((0,2i]\big),
\]
is cleared, as derived, by $2^{\ast}(k-1)!^2\,d_{2k-3}\operatorname{lcm}(\text{odd}\le2k-3)^3$; observed (\S9.6(iii)), by $2^{\ast}d_k^{\,3}$. The difference is two congruences with no analysis left in them: the odd part of $(k-1)!^2$ divides the moment part, and for every prime $p=2n-1$ with $k<p<2k$ the coefficient $\sum_{i\ge n}\binom{2i}i^3\binom{2j}j^3\binom{k-1}i^{-2}$ of $(2n-1)^{-3}$, with the moment part, vanishes mod $p^3$. True to $k=200$, not derived. Both $A_k$ and $B_k$ were checked exactly against the PSLQ rows for $k\le9$ before this entry was written. (Both congruences, and the clearing they stand for, are proved in \S10.)

\section{The four-mark centre ray, proved: the row, the port, the clearing}

\S9.7 left the clearing of the four-mark centre ray at two congruences, and \S9.4 left the port as a finite falsifier. Both are settled here for every $k$, from the word of \S8.6 and nothing else. Write $n=k-1$, $p_l=l+\frac12$, $j=n-l$,
\[
c_l=\binom{2l}{l}\binom{2j}{j},\qquad S_l=\frac{c_l}{\binom nl}=\frac{(2l)!\,(2j)!}{l!\,j!\,n!},\qquad w_l=S_l^{\,2}\,c_l,\qquad F(\sigma)=\binom{\sigma-1}{n}^{\!2},
\]
$F[\sigma,p]=\big(F(\sigma)-F(p)\big)/(\sigma-p)$, and, for a polynomial $g$, the critical-line functional
\[
\mathcal J[g]=\frac1{2\pi i}\int_{(\frac12)}g(\sigma)\sum_{m\in\mathbb Z}\frac{d\sigma}{(\sigma+m)^3}=\frac12\,\oE\big[g'(X)\big],\qquad X=\tfrac12+iY,\quad Y\sim\mathbb P_{\mathrm{crit}},
\]
the second form by parts, since $\sum_m(\sigma+m)^{-3}=-\frac12\,\partial_\sigma\Gst40{\sigma,\ \sigma,\ 1-\sigma,\ 1-\sigma}{}$ (Lemma~\hb{5.16}{34}) and that stacked Gamma on the line is $2\pi p_{\mathrm{crit}}$ (Def.~\hb{33.1}{297}); Thm.~\hb{33.2}{297}(iii) gives $\oE[e^{tX}]=t/(1-e^{-t})$. $S_l$ is an integer (alias, once: a super Catalan number) by Legendre's formula and the floor inequality $\lfloor2x\rfloor+\lfloor2y\rfloor\ge\lfloor x\rfloor+\lfloor y\rfloor+\lfloor x+y\rfloor$. The $c_l$ are the constant terms of bunch and room, $c_l=\CT[(XX')^l]\,\CT[(YY')^{j}]$ in the notation of \S10.2. Every step below is also checked in exact arithmetic for $n\le45$; the proofs do not rest on the checks.

\subsection{The row at every step}\status{derived}

\begin{claim}\label{cl:row4}
For every $n\ge0$,
\[
256^n\langle1\rangle_{n+1}=\sum_{l=0}^{n}w_l\,\Tx{3}{\Ilat{\frac12+\mathbb Z}}{l}\;-\;16^n\sum_{l=0}^{n}c_l\;\mathcal J\big[F[\,\cdot\,,p_l]\big].
\]
Hence $\langle1\rangle_{n+1}=A_{n+1}\Tx3\IZ0-B_{n+1}$ with $256^nA_{n+1}=7\sum_lw_l$ and
\[
256^nB_{n+1}=\sum_{l=0}^{n}w_l\,\Th{3}{\Ilat{\frac12+\mathbb Z}}{0}{l}\;+\;16^n\sum_{l=0}^{n}c_l\;\mathcal J\big[F[\,\cdot\,,p_l]\big].
\]
\end{claim}

\begin{enumerate}[label=(\alph*)]
\item The word of \S8.6 is $\langle1\rangle_k=\Gamma(k)^{-2}\frac1{2\pi i}\int\Gamma(\sigma)^4\Gamma(k-\sigma)^4\big/\big(\Gamma(2\sigma)\Gamma(2k-2\sigma)\big)\,d\sigma$ on $0<\re\sigma<k$. Duplication and reflection (Prop.~\hb{1.3}{11}) turn it into
\[
\langle1\rangle_k=\frac{1}{4^n\,n!^2}\,\frac1{2\pi i}\int\rho(\sigma)\sum_{m}\frac{d\sigma}{(\sigma+m)^3},\qquad\rho(\sigma)=\frac{(1-\sigma)_n^{\,3}}{(\frac12-\sigma)_{n+1}}.
\]
The triple poles of the kernel at $1,\dots,n$ meet triple zeros of $\rho$, and the simple poles of $\rho$ at the $p_l$ meet zeros of the kernel, so the line may be $\re\sigma=\frac12$.
\item With $\Omega(\sigma)=\prod_{i=1}^n(\sigma-i)=n!\binom{\sigma-1}{n}$ one has $(1-\sigma)_n=(-1)^n\Omega$ and
\[
\frac{\rho}{\Omega^2}=\frac{(-1)^n\,\Omega(\sigma)}{(\frac12-\sigma)_{n+1}}=-\frac1{4^n}\sum_{l=0}^{n}\frac{c_l}{\sigma-p_l}:
\]
the residue at $p_l$ is $(-1)^{n+1}\Omega(p_l)/\prod_{i\ne l}(i-l)$, with $\Omega(p_l)=2^{-n}(-1)^{n-l}(2l-1)!!\,(2j-1)!!$. Multiplying by $\Omega^2=n!^2F$ splits $\rho$ into a polynomial and its poles,
\[
\rho=P+\sum_{l=0}^n\frac{r_l}{\sigma-p_l},\qquad P=-\frac{n!^2}{4^n}\sum_{l}c_l\,F[\sigma,p_l],\qquad r_l=-\frac{c_l\,\Omega(p_l)^2}{4^n}=-\frac{n!^2\,w_l}{64^n}.
\]
\item A pole term closes to the right on the triple poles of the kernel at the positive integers (Thm.~\hb{4.5}{26}); the lines $\re\sigma=N+\frac12$ carry $O(N^{-1})$. So
\[
\frac1{2\pi i}\int_{(\frac12)}\frac1{\sigma-p_l}\sum_m\frac{d\sigma}{(\sigma+m)^3}=-\sum_{i\ge1}\frac1{(i-p_l)^3}=-\Tx{3}{\Ilat{\frac12+\mathbb Z}}{l},
\]
the terms $i\le l$ removing the head $\Th{3}{\Ilat{\frac12+\mathbb Z}}{0}{l}$. The polynomial part gives $\mathcal J[P]$.
\item Assembling (a)--(c) gives the claim, and $\Tx3{\Ilat{\frac12+\mathbb Z}}0=7\,\Tx3\IZ0$ (Thm.~\hb{5.12}{32}(i) at $q=2$). At $n=1$: $w_0=w_1=8$, $c_0=c_1=2$, $\mathcal J[F[\cdot,\frac12]]=\mathcal J[F[\cdot,\frac32]]=\frac12$, and $256\langle1\rangle_2=112\,\Tx3\IZ0-96$, the row $(\frac7{16};\frac38)$.
\end{enumerate}

\subsection{The top coordinate is a port}\status{derived}

\begin{claim}\label{cl:port4}
For every $n\ge0$,
\[
\sum_{l=0}^{n}S_l^{\,2}\,c_l=4^n\sum_{l=0}^{n}\binom nl^{\!2}c_l=4^n\,\CT\big[L^n\big],\qquad L=\Big(\sum_e\alpha_e\Big)\Big(\sum_e\alpha_e^{-1}\Big),
\]
$L$ the inverse of the stride of the centre ray (Def.~\hb{35.9}{309}). Hence $A_k=7\cdot64^{-n}\CT[L^n]$ and $64^nA_k\in7\mathbb Z$ for every $k$: the falsifier of \S9.4 passes for the top coordinate, the coefficient being a finite sum across the hinge $\{1,2\}\,|\,\{3,4\}$ of bunch constant terms times room constant terms, with no chamber.
\end{claim}

In $L^n=(X+Y)^n(X'+Y')^n$, with $X=\alpha_1+\alpha_2$, $X'=\alpha_1^{-1}+\alpha_2^{-1}$ and $Y,Y'$ the same on $\{3,4\}$, only equal powers of $X$ and $X'$ have a constant term in $\alpha_1,\alpha_2$, and $\CT[(XX')^l]=\sum_i\binom li^2=\binom{2l}l$: that is the middle sum. Put $T(n)=4^{-n}\sum_lS_l^2c_l$ and $U(n)=\sum_l\binom nl^2c_l$. Both satisfy
\[
(n+2)^3X_{n+2}-2(2n+3)(5n^2+15n+12)\,X_{n+1}+64(n+1)^3X_n=0\qquad(n\ge0),
\]
by telescoping in $l$. For the summand $t(n,l)$ of either sum,
\begin{multline*}
64(n+1)^3t(n,l)-2(2n+3)(5n^2+15n+12)\,t(n+1,l)+(n+2)^3t(n+2,l)\\
=G(n,l+1)-G(n,l),\qquad G=R\,t,
\end{multline*}
with the certificates
\[
R_T=\frac{4\,l\,(2j+1)^3\,\kappa_T}{(n+1)^2(j+1)(j+2)},\qquad R_U=-\frac{4\,(n+1)^2\,l^3\,(2j+1)\,\kappa_U}{(j+1)^3(j+2)^3},
\]
\begin{align*}
\kappa_T&=12j^3+10lj^2+2l^2j+59j^2+33lj+4l^2+94j+26l+49,\\
\kappa_U&=12j^3+10lj^2+2l^2j+62j^2+35lj+3l^2+104j+30l+56.
\end{align*}
The ratios of the summands are rational: $t_T(n,l+1)/t_T(n,l)=(2l+1)^3j/((l+1)(2j-1)^3)$, $t_T(n+1,l)/t_T(n,l)=2(2j+1)^3/((j+1)(n+1)^2)$, $t_U(n,l+1)/t_U(n,l)=j^3(2l+1)/((l+1)^3(2j-1))$, $t_U(n+1,l)/t_U(n,l)=2(n+1)^2(2j+1)/(j+1)^3$. Dividing by $t(n,l)$ turns each relation into an identity of polynomials in $(n,l)$ of total degree $9$, for $0\le l\le n-1$; summing it, with $G(n,0)=0$, leaves the boundary $G(n,n)$ plus the left side at $l=n,n+1,n+2$, which divided by $t(n,n)$ is an identity of polynomials in $n$. The four identities hold as polynomial identities (verified exactly; the certificates were read off the exact sums, and the proof uses only the identities). With $T(0)=U(0)=1$, $T(1)=U(1)=4$ and $(n+2)^3\ne0$, $T=U$. The recurrence that \S8.7 fitted is thereby derived for the top coordinate: $A_k=7\cdot64^{-n}U(n)$.

\subsection{The clearing}\status{derived}

\begin{claim}\label{cl:clear4}
For every $k\ge1$ the odd part of $\operatorname{den}(B_k)$ divides $d_{k-1}^{\,3}$, and $256^{k-1}d_{k-1}^{\,3}B_k\in\mathbb Z$. With Claim~\ref{cl:port4}, the row of $\langle1\rangle_k$ is cleared by $2^{\ast}d_{k-1}^{\,3}$: the frame on the coordinate and the heads on the rational part. The two congruences of \S9.7 are consequences: neither the odd part of $(k-1)!^2$ nor the primes between $k$ and $2k$ survive.
\end{claim}

Fix an odd prime $q$, let $e\ge1$ with $q^{e-1}\le n<q^e$, and write $v$ for the $q$-adic valuation. By Claim~\ref{cl:row4} it suffices that each of the two sums in $256^nB_{n+1}$ has $v\ge-3(e-1)$.

\emph{The heads.} $\Th{3}{\Ilat{\frac12+\mathbb Z}}{0}{l}=8\sum_{i=1}^{l}(2i-1)^{-3}$, and $2l-1<2n<q^{e+1}$, so the only odd number up to $2l-1$ with $v\ge e$ is $q^e$ itself, when $q^e\le2l-1$. In that case Legendre's formula at the level $q^e$, where $\lfloor l/q^e\rfloor=\lfloor j/q^e\rfloor=\lfloor n/q^e\rfloor=0$ and $\lfloor2l/q^e\rfloor=1$, gives $v(S_l)\ge1$ and $v\binom{2l}l\ge1$, every other level contributing $\ge0$ by the floor inequality; so $v(w_l)\ge3$, which pays for $q^{-3e}$. Otherwise the head already has $v\ge-3(e-1)$.

\emph{The moments.} Write a polynomial in the binomial basis, $g=\sum_i(\Delta^ig)(0)\binom\sigma i$. From $\oE[(1+z)^X]=(1+z)\log(1+z)/z$,
\[
\mathcal J\Big[\binom\sigma i\Big]=\lambda_i=(-1)^i\,\frac{H_{i-1}-1}{i(i+1)}\quad(i\ge1),\qquad\lambda_0=0,
\]
$H_m$ the harmonic number. Vandermonde twice gives the binomial coefficients of the divided difference,
\[
\big(\Delta^iF[\,\cdot\,,p]\big)(0)=\sum_{m}f_m\sum_{t=i+1}^{m}\frac1t\binom{p}{m-t}\binom{-p-1}{t-1-i},\qquad F=\sum_mf_m\binom\sigma m,
\]
with $f_m\in\mathbb Z$ because $F$ is integer-valued, and $q\mid f_m$ for $m\ge q^e$: $f_m=\sum_{i,i'\le n}(-1)^{i+i'}\binom mi\binom i{m-i'}$, and when $m\ge q^e>n$ the addition $i+(m-i)$ of two numbers below $q^e$ carries, so $q\mid\binom mi$ (Kummer). At $p=p_l\in\mathbb Z_q$ the binomials of $p$ and of $-p-1$ are $q$-integral. Since $t\le2n<2q^e$, the only $t$ with $v(t)\ge e$ is $t=q^e$, where $m\ge q^e$ and $q\mid f_m$. So $v\big((\Delta^iF[\cdot,p_l])(0)\big)\ge1-e$ for every $i$, and $\ge2-e$ for $i\ge q^e$, where every $t$ exceeds $q^e$. For $1\le i\le2n-1$: $v(\lambda_i)\ge2-2e$ for $i<q^e$, and $v(\lambda_i)\ge1-2e$ for $i\ge q^e$, using $v(H_{q^e-1})\ge2-e$, which holds because $H_{q^e-1}=\sum_{s<e}q^{-s}\sum_{r<q^{e-s},\,q\nmid r}1/r$ and $\sum_{r<q^a,\,q\nmid r}1/r$ has $v\ge a$ (pair $r$ with $q^a-r$). In each case the product has $v\ge3-3e$, and $16^n$ and the $c_l$ are $q$-integral.

\emph{The prime $2$.} The same count with $v_2\binom pr\ge-2r$ for $p\in\frac12+\mathbb Z$ gives $v_2\ge1-3\lfloor\log_2n\rfloor$ for the moment sum, the factor $16^n$ absorbing the $2^{-4n}$ of the binomials, and the heads are $2$-integral; hence $256^nd_n^{\,3}B_{n+1}\in\mathbb Z$. The sharper $64^{k-1}d_{k-1}^{\,3}B_k\in\mathbb Z$ is computed ($k\le40$), not derived.

\subsection{What this settles}\status{derived, with what is not}

\begin{enumerate}[label=(\roman*)]
\item For every $k$: the row of the four-mark centre ray, its coordinate as a port with $\operatorname{den}(A_k)\mid64^{k-1}$, and its clearing, $\operatorname{den}(B_k)\mid256^{k-1}d_{k-1}^{\,3}$ with odd part dividing $d_{k-1}^{\,3}$.
\item The mechanism is internal. The residues of $\rho/\Omega^2$ on the coset $\frac12+\mathbb Z$ are the products $c_l$ of the bunch and room constant terms, the same numbers the torus hinge sums in Claim~\ref{cl:port4}: the closing's poles and the port read the same faces. The polynomial part of the word is those poles carried by $\Omega^2=n!^2F$, so the $n!^2$ of $\Gamma(k)^{-2}$ cancels before any expectation is taken, and what the critical-line law contributes is the numbers $\lambda_i$, whose denominators are heads. The primes between $k$ and $2k$ are paid by the constant terms $\binom{2l}l$ inside $S_l$.
\item Not settled: the $2$-adic sharpness for $B_k$; the recurrence of the rational parts (computed, \S8.7); the same derivation on three marks (\S8.5), where the coset is replaced by a frequency $\frac13$.
\end{enumerate}
Checks: \texttt{clearing\_proof\_checks.py} (every step, $n\le45$, every odd prime $q\le2n$) and \texttt{wz\_symbolic.py} (the four polynomial identities of \S10.2).

\section{The gap is a theorem; the four-mark tree is one atom}

Two questions were asked of \S9 and \S10.
\begin{itemize}
\item \emph{Is the gap read off the pencil settled?} It is a theorem, and \S10 proved its inputs (\S11.1). Calling it a hypothesis in discussion was a drift.
\item \emph{Is the tree the coaction?} The Foreman tree of the four-mark centre ray is one atom (\S11.2). Def.~\hb{13.6}{78} returns no bunch, and the tree records no hinge. The coaction it predicts is the coaction of the content alone, and that agrees with the period.
\end{itemize}
The hinges at $B=\{1\}$ and $B=\{1,2\}$, used below and in \S10, evaluate that atom; neither is a node of the tree. An earlier version of this entry drew conclusions about the tree from the hinge at $B=\{1\}$. Those conclusions are withdrawn. The computations are kept, as facts about an evaluation.

\subsection{The gap, read off the pencil}\status{derived; inputs for four marks from \S10}

\begin{claim}\label{cl:gappencil}
Let a centre ray (Def.~\hb{35.9}{309}) have values $x_k=\langle1\rangle_k$ with a row $(A_k;B_k)$ in one coordinate $v$ (Def.~\hb{34.1}{300}). Let $\lambda=\min_{\mathbb R^{\cE}_{>0}}L$ be the orthant fibre of its pencil (Thm.~\hb{35.10}{309}). Let $D_k=c^{k-1}d_{k-1}^{\,w}$ clear the row for every $k$. If
\[
c\,e^{w\,c_{\mathrm{lcm}}}<\lambda,
\]
then $v$ is irrational.
\end{claim}

\emph{Proof.} The read is strictly positive, so $x_k>0$. Since $L\ge\lambda$ on the orthant, $L^{-k}\le\lambda^{-(k-1)}L^{-1}$, and so $x_k\le\lambda^{-(k-1)}x_1$. Then
\[
0<D_kx_k\le x_1\Big(\frac{c\,d_{k-1}^{\,w/(k-1)}}{\lambda}\Big)^{k-1}\longrightarrow0
\]
by Lemma~\hb{34.2}{300}. Thm.~\hb{34.3}{300}(i) concludes. $\square$

\emph{Four marks.} The inputs are proved: $\lambda=16$ ($L\ge16$, with equality at the centre); $c=256$ and $w=3$ (Claims~\ref{cl:port4}, \ref{cl:clear4}). The inequality is
\[
\frac{c\,e^{3c_{\mathrm{lcm}}}}{\lambda}=16\,e^{3c_{\mathrm{lcm}}}>1,
\]
and it stays above $1$ with $c_{\mathrm{lcm}}$ replaced by $1$. So the four-mark ray does not pass the criterion, and the verdict is a theorem. A stronger statement would be that no clearing of this row can pass. It needs only $c\ge64\cdot2^{-o(1)}$, i.e.\ that $\operatorname{den}(A_k)$ is $64^{k-1}$ up to $2^{o(k)}$, since $64>16$. Equivalently $v_2\big(\CT[L^{k-1}]\big)=o(k)$, which is computed ($\le15$ for $k\le40$), not derived.

\subsection{The Foreman tree of the four-mark centre ray}\status{derived}

At the table's gauge, every permutation $g$ of the four marks fixes both Foreman polynomials,
\[
K=e_1(\alpha),\qquad \Psi=e_3(\alpha).
\]
So $g$ carries each view of a mark $e$ (Def.~\hb{10.1}{64}) to the view of $g(e)$ with the partners moved by $g$. The readings $\operatorname{read}^{(3)}(e)$ are therefore one multiset for all four marks, and
\[
\operatorname{Cl}(\Upsilon_{\boldsymbol\sigma,\boldsymbol\delta;\one})=\{\cE\}\qquad\text{at the centre and at every lean,}
\]
the lean being read through $K$ and $\Psi$, which $g$ fixes separately. The realization is atomic, so $\operatorname{Tree}(F;\boldsymbol\sigma,\boldsymbol\delta)$ (Def.~\hb{14.4}{81}) is the one atom. Def.~\hb{13.6}{78} has no class to return except $\cE$ itself, which is not the bunch of a peel. The tree takes no bunch and records no hinge.

The reference evaluates this atom by Cor.~\hb{19.40}{153}. Here $\binom{\cE}1$ is the table that is one circuit (Cor.~\hb{17.20}{113}) with its two Foreman polynomials exchanged, which at the centre $s=t$ changes nothing. With unit weights the evaluation removes one mark at a time, at the root rates $\frac19$ and $\frac14$, and ends at four Gammas over one.

\emph{What the tree says of the coaction.} \S9.2 reads one coaction term per hinge down the tree, and here there is none. So the tree predicts that the only coaction is the coaction of the content. The value is the row of Claims~\ref{cl:row4} and~\ref{cl:port4},
\[
\langle1\rangle_k=A_k\cdot7\,\Tx3{\IZ}0-B_k .
\]
Its one coordinate is a trail of depth one at the cut $0$, whose deconcatenation (\S9.3) has no proper cut. The tree's prediction is therefore that the reduced coaction of $\langle1\rangle_k$ vanishes. On four marks the tree and the coaction agree. (Alias, once: this is the primitivity of $\zeta^{\mathfrak m}(3)$, recalled.)

\subsection{Evaluating the atom: two closings}\status{derived; both sides checked to $200$ digits}

The first lattice of Cor.~\hb{19.40}{153} removes the mark $1$. Take the centre ray at $k=1$ and write that removal as the hinge of Thm.~\hb{18.13}{119} at $B=\{1\}$, with room $R=\{2,3,4\}$. Put $\alpha_1=t$ and give the room projective coordinates $v$. Then
\[
L=(t+V)\big(t^{-1}+V'\big),\qquad V=v_2+v_3+v_4,\quad V'=v_2^{-1}+v_3^{-1}+v_4^{-1},\quad VV'=L_3,
\]
where $L_3$ is the room's own stride reciprocal. Integrate the scale by partial fractions:
\[
\int_0^\infty\frac{dt}{t}\,\frac1L=\int_0^\infty\frac{dt}{(t+V)(1+tV')}=\frac{\log L_3}{L_3-1}.
\]
This gives the \emph{left closing}:
\[
7\,\Tx3{\IZ}0=\langle1\rangle_1=\int\frac{\log L_3}{L_3-1}\,d\mu_3 .
\]
It is Lemma~\hb{5.16}{34}(ii), applied under the room integral. To see this, write the share integral of Thm.~\hb{18.13}{119} at $B=\{1\}$:
\[
\langle1\rangle_1=\frac1{2\pi i}\int_{(1/2)}\Gst40{a,\ a,\ 1-a,\ 1-a}{}\,P_3(a)\,da,\qquad P_3(a)=\int L_3^{-a}\,d\mu_3 .
\]
Lemma~\hb{5.16}{34}(i), the right closing, gives the \emph{right closing} of the same integral:
\[
\langle1\rangle_1=-\sum_{c\ge1}P_3'(c).
\]
The left closing gives the same line through the expansion $\log L_3/(L_3-1)=\sum_{c\ge1}L_3^{-c}\log L_3$, valid because $L_3\ge9$.

The left closing is an integral, over a semi-algebraic domain, of an algebraic function times one logarithm. It is reached from $\langle1\rangle_1$ by three steps: a change of variables (the scale of the bunch), Fubini, and an elementary integral. Each step is a relation of Def.~\hb{36.1}{315}. So the hinge with its left closing is an identity in $\mathfrak W/\mathfrak I$ with no further hypothesis. At $B=\{1,2\}$, the hinge of Claim~\ref{cl:row4}, the same steps give the symmetric kernel
\[
\langle1\rangle_1=\iint\frac{\log L_B-\log L_R}{L_B-L_R}\,d\mu_B\,d\mu_R,\qquad L_B=\frac{(1+x)^2}x,\ \ L_R=\frac{(1+y)^2}y,
\]
which is the peel formula with $L_B$, $L_R$ in place of $\langle\cdot\rangle$. (Recalled, as an import: relations of this kind hold between motivic periods, by Huber--M\"uller-Stach.)

\subsection{The room's coordinates}\status{computed: integer relations at $240$--$760$ digits; $\alpha_c=s_c/3$ derived from them}

Put
\[
L:=\tfrac13P_3(1)=\Tx{2}{(\Ilat{1+3\mathbb Z}-\Ilat{2+3\mathbb Z})}{0},\qquad \Lambda:=P_3'(1)=-\int L_3^{-1}\log L_3\,d\mu_3 .
\]
Every value and every jet of the room has a row in $\{1,L,\Lambda\}$:
\[
P_3(c)=r_c+s_cL,\qquad P_3'(c)=\beta_c+\gamma_cL+\alpha_c\Lambda .
\]
\[
\begin{array}{c|cc|ccc}
c&r_c&s_c&\beta_c&\gamma_c&\alpha_c\\\hline
1&0&3&0&0&1\\
2&-\frac23&1&\frac43&\frac43&\frac13\\[2pt]
3&-\frac{23}{54}&\frac59&\frac{13}{18}&\frac{26}{27}&\frac5{27}\\[2pt]
4&-\frac{145}{486}&\frac{31}{81}&\frac{2033}{4374}&\frac{526}{729}&\frac{31}{243}\\[2pt]
5&-\frac{1331}{5832}&\frac{71}{243}&\frac{35555}{104976}&\frac{1253}{2187}&\frac{71}{729}
\end{array}
\]
The ratio $\alpha_c=s_c/3$ holds for every $c\le30$. It is also derived, from three inputs:
\begin{itemize}
\item the room's ray has order two (Thm.~\hb{35.19}{313}(iii));
\item its recurrence, differentiated in $c$, has the same homogeneous solutions $u_c,v_c$, so that $P_3'(c)=u_c\Lambda+v_cP_3'(2)+(\text{terms in }1,L)$;
\item $P_3'(2)=\frac13\Lambda+\frac43+\frac43L$.
\end{itemize}
From these, $\alpha_c=u_c+\frac13v_c$, and $s_c=3u_c+v_c$.

The coordinates fall like $c^{-1}$, while the values fall like $9^{-c}$:
\[
\begin{array}{c|cccc|c}
c&c\,s_c&c\,\alpha_c&c\,\gamma_c&c\,\beta_c&|P_3'(c)|\\\hline
10&1.3412&0.4471&2.8018&1.4210&1.1\cdot10^{-9}\\
20&1.2888&0.4296&2.7649&1.3089&1.5\cdot10^{-19}\\
30&1.2723&0.4241&2.7520&1.2745&2.9\cdot10^{-29}
\end{array}
\]
At $c=30$ each of the three terms of $P_3'(c)$ has size about $10^{-2}$, and they cancel to $10^{-29}$.

Cor.~\hb{19.40}{153} sums from the inside out. The inner lattices produce these room values, and the outer lattice cancels their $L$- and $\Lambda$-parts. The cancellation is analytic: the pieces fall like $9^{-c}$, their coordinates like $c^{-1}$, and the coordinate sums diverge.

Two readings of where the room's coordinates come from (derived):
\begin{itemize}
\item \emph{The room's pencil.} It has the critical value $1$,
\[
L_3-1=\frac{(v_2+v_3)(v_2+v_4)(v_3+v_4)}{v_2v_3v_4},
\]
and a one-mark removal, with $L_B=1$, reads it there.
\item \emph{The arrangements.} The room's $\{e_1=e_2=0\}=\{(1:\omega:\omega^2),(1:\omega^2:\omega)\}$ meets mod $3$. The list's $e_1\cap e_3$, the triangle $\{\alpha_i+\alpha_j=\alpha_k+\alpha_l=0\}$, meets only mod $2$.
\end{itemize}
None of this is a statement about the tree, which removes nothing here.

Alias, once: $L=L(\chi_{-3},2)$, and (computed, $254$ digits)
\[
\Lambda=\sqrt3\Big(8\,\im\operatorname{Li}_3\big(e^{i\pi/6}/\sqrt3\big)-\frac{35}{162}\pi^3-\frac{\pi}6\log^23\Big).
\]
Checks: \texttt{room\_jets.py}, \texttt{room\_jets\_ident.py}, \texttt{room\_jets\_coeffs.py}.

\subsection{The first test on a tree that peels}\status{the classes, the bunch and the peel derived; the value and the comparison not yet computed}

Take five marks and the list of all pairs except $\{1,2\}$: rank $2$, with $9$ bases. Its two Foreman polynomials are
\[
K=e_2(\alpha)-\alpha_1\alpha_2,\qquad \Psi=e_3(\alpha)-\alpha_3\alpha_4\alpha_5 .
\]
\emph{Classes.} The swap of $1,2$ and the permutations of $3,4,5$ fix the list, so $\{1,2\}$ and $\{3,4,5\}$ each lie in one class. The occupations at the centre are
\[
R_1(\one)=\tfrac39=\tfrac13,\qquad R_3(\one)=\tfrac49 .
\]
Since $\frac49\notin\{\frac13,\frac23\}$, Prop.~\hb{10.9}{66}(ii) separates the two:
\[
\operatorname{Cl}=\big\{\{1,2\},\ \{3,4,5\}\big\},\qquad B(M;\boldsymbol\sigma,0)=\{1,2\}
\]
by Def.~\hb{13.6}{78}, the class with the fewest marks.

\emph{Width one.} A basis meets $B$ in $0$ or $1$ marks, and a cobasis in $1$ or $2$. So Thm.~\hb{18.13}{119} applies at the Foreman bunch.

\emph{The peel} (Def.~\hb{14.1}{81}). The patterns of $B$ are $\emptyset$, $\{1\}$, $\{2\}$, each with probability $\frac13$. The children are:
\begin{itemize}
\item at $\emptyset$, the pairs on $R=\{3,4,5\}$, with exponents $(\frac s3,\frac t3)$;
\item at $\{1\}$ and $\{2\}$, the singles on $R$; these have the same bases, so they merge into one table with exponents $(\frac{2s}3,\frac{2t}3)$.
\end{itemize}
The singles have $(K,\Psi)=(e_1,e_2)$ and the pairs have $(e_2,e_1)$. So at the centre $s=t$ the two children multiply to
\[
(e_1e_2)^{-s}\quad\text{on }R .
\]
That is the three-mark list at its own centre: an atom, whose content is $L=\Tx{2}{(\Ilat{1+3\mathbb Z}-\Ilat{2+3\mathbb Z})}{0}$ and its jets (\S11.4). The bunch family on $\{1,2\}$ is a Beta corner.

\emph{The tree's prediction.} There is one hinge term, which pairs the Beta data of the bunch with the three-mark room. So the row of the five-mark value has its coordinates in the room's content, entering through that one term.

\emph{The test.} Compute the value and its row by the hinge at the Foreman bunch, then compare the row's coaction with that term.

\section{A width-one node separates; a dual pair factors}

Notation for this entry: $M$ is a list on $\cE$ of rank $r$, $B$ a bunch with room $R$, and for a list $N$ write $\langle N\rangle=K_N\Psi_N$. Width one means
\[
r(B)+r(R)-r=1;\qquad\text{put}\quad d=r-r(R),\quad\text{so } r(B)=d+1 .
\]
Since $K_{N^\ast}=\Psi_N$ and $\Psi_{N^\ast}=K_N$,
\[
\langle N^\ast\rangle=\langle N\rangle .
\]

\subsection{Two levels}\status{derived}

\begin{claim}\label{cl:twolevels}
At a bunch of width one,
\begin{align*}
K(\tau\alpha_B,\alpha_R)&=\tau^{d}\big(K_{M/R}(\alpha_B)\,K_{M|R}(\alpha_R)+\tau\,K_{M|B}(\alpha_B)\,K_{M/B}(\alpha_R)\big),\\
\Psi(\tau\alpha_B,\alpha_R)&=\tau^{|B|-d-1}\big(\Psi_{M|B}(\alpha_B)\,\Psi_{M/B}(\alpha_R)+\tau\,\Psi_{M/R}(\alpha_B)\,\Psi_{M|R}(\alpha_R)\big).
\end{align*}
\end{claim}

\emph{Proof.} For a basis $T$, $T\cap R$ is independent in $M|R$, so $|T\cap B|\ge d$. Equality holds exactly when $T\cap R$ is a basis of $M|R$ and $T\cap B$ a basis of $M/R$. Likewise $|T\cap B|\le r(B)=d+1$, with equality exactly when $T\cap B$ is a basis of $M|B$ and $T\cap R$ a basis of $M/B$. Width one leaves no other level. $\Psi$ sums $\alpha^{\cE\setminus T}$ over the same bases, and $|B\setminus T|=|B|-|T\cap B|$. $\square$

\subsection{The node separates}\status{derived}

\begin{claim}\label{cl:separation}
Let $s,t>0$, let the masses lie on the balance wall, and put
\[
\kappa=m_B-sd-t(|B|-d-1),\qquad 0<\kappa<s+t .
\]
Then
\[
P_F(\mvec;s,t)=\int_{\mathbb P(\mathbb R^B_{>0})}\int_{\mathbb P(\mathbb R^R_{>0})}F_B\,F_R\;J\big(x_B\,y_R\big)\;\Omega_B\,\Omega_R ,
\]
with the two faces' data
\begin{gather*}
F_B=\alpha_B^{m_B}K_{M/R}^{\kappa-s}K_{M|B}^{-\kappa}\Psi_{M|B}^{-t},\qquad F_R=\alpha_R^{m_R}K_{M|R}^{\kappa-s}K_{M/B}^{-\kappa}\Psi_{M/B}^{-t},\\
x_B=\frac{\langle M/R\rangle}{\langle M|B\rangle}\ \text{on }B,\qquad y_R=\frac{\langle M|R\rangle}{\langle M/B\rangle}\ \text{on }R,
\end{gather*}
both of degree zero, and the kernel the scale integral of two linear letters, one Breta of one bunch (Prop.~\hb{15.1}{87}):
\[
J(w)=\int_0^\infty u^{\kappa-1}(1+u)^{-s}(1+wu)^{-t}\,du=w^{-t}\,\Gst22{\kappa,\ s+t-\kappa}{s,\ t}\,\sB^{\{0,1\}}_{(s,t),\,s+t-\kappa;\,I}\big(1,w^{-1}\big).
\]
\end{claim}

\emph{Proof.} Write $\alpha_B=\tau\beta$. The masses and Claim~\ref{cl:twolevels} give the factor $\tau^{\kappa}$, so
\[
P_F=\int\!\!\int\int_0^\infty \tau^{\kappa}\,(K_{\mathrm{lo}}+\tau K_{\mathrm{hi}})^{-s}(\Psi_{\mathrm{lo}}+\tau\Psi_{\mathrm{hi}})^{-t}\,\frac{d\tau}\tau\,(\cdots),
\]
with $K_{\mathrm{lo}}=K_{M/R}K_{M|R}$, $K_{\mathrm{hi}}=K_{M|B}K_{M/B}$, $\Psi_{\mathrm{lo}}=\Psi_{M|B}\Psi_{M/B}$ and $\Psi_{\mathrm{hi}}=\Psi_{M/R}\Psi_{M|R}$. Substitute $\tau=(K_{\mathrm{lo}}/K_{\mathrm{hi}})\,u$. The prefactor $(K_{\mathrm{lo}}/K_{\mathrm{hi}})^\kappa K_{\mathrm{lo}}^{-s}\Psi_{\mathrm{lo}}^{-t}$ splits as $F_BF_R$, and the variable of the kernel is
\[
w=\frac{K_{\mathrm{lo}}\Psi_{\mathrm{hi}}}{K_{\mathrm{hi}}\Psi_{\mathrm{lo}}}=\frac{K_{M/R}\Psi_{M/R}}{K_{M|B}\Psi_{M|B}}\cdot\frac{K_{M|R}\Psi_{M|R}}{K_{M/B}\Psi_{M/B}}=x_B\,y_R .
\]
Both $F_B$ and $F_R$ have degree zero; for $F_R$ this is the balance wall. The kernel is the scale integral with the coefficients $1$ and $w^{-1}$ and the letter exponents $s$ and $t$, which is Prop.~\hb{15.1}{87}. $\square$

The share side of the kernel (Cor.~\hb{15.4}{88}, one share) is
\[
J(w)=\frac1{2\pi i}\int\Gst42{a,\ t-a,\ \kappa-a,\ s-\kappa+a}{s,\ t}\,w^{-a}\,da .
\]
Put into Claim~\ref{cl:separation}, it gives the share integral of Thm.~\hb{18.13}{119}: the bunch family and the room family are read at $F_B\,x_B^{-a}$ and $F_R\,y_R^{-a}$. So the claim is the hinge read before its share is opened. At $s=t=\kappa=1$ the share integrand is $\Gst40{a,\ a,\ 1-a,\ 1-a}{}\,w^{-a}$, and $J(w)=\log w/(w-1)$ is its left closing (Lemma~\hb{5.16}{34}(ii)); at integer data in general the poles of the share integrand coincide, and the closing reads Laurent data of the Gammas (Prop.~\hb{5.14}{33}). The claim is then reached from $P_F$ by a birational change of variables and Fubini, both relations of Def.~\hb{36.1}{315}. (Recalled, as an import: such relations hold between motivic periods, Huber--M\"uller-Stach.)

\subsection{A dual pair factors}\status{derived}

\begin{claim}\label{cl:dualpair}
If $M/B=(M|R)^\ast$ as lists on $R$, then $y_R\equiv1$ and
\[
P_F(\mvec;s,t)=\Big[\int F_B\,J(x_B)\,\Omega_B\Big]\cdot\Big[\int F_R\,\Omega_R\Big].
\]
At integer data both factors are periods. The identity is a relation of Def.~\hb{36.1}{315}, so the motivic period of the node is the product of the two motivic periods, and its coaction is the product of their coactions. The same holds with the roles exchanged when $M/R=(M|B)^\ast$.
\end{claim}

\emph{Proof.} $\langle M/B\rangle=\langle(M|R)^\ast\rangle=\langle M|R\rangle$, so $y_R\equiv1$ in Claim~\ref{cl:separation}, and the integral over $R$ splits off. The coaction is a homomorphism of algebras. $\square$

(Later: \S14.5 replaces the one cell by several, Claim~\ref{cl:pinned14}, and the width of the bunch drops out.)

\subsection{The five-mark list of \S11.5}\status{derived}

Take all pairs except $\{1,2\}$, with the Foreman bunch $B=\{1,2\}$. The four minors are:
\[
M|R=\binom R2,\qquad M/B=\binom R1=\Big(\binom R2\Big)^{\!\ast},\qquad M|B=\binom B1,\qquad M/R=\binom B0 .
\]
So Claim~\ref{cl:dualpair} applies.

On the centre ray at $k$ every mass is $k$, and $\kappa=k$. Then
\begin{gather*}
F_R=\alpha_R^{k}\,(e_1e_2)^{-k}\quad\text{(the three-mark centre ray at $k$)},\\
F_B=x_B^{\,k},\qquad x_B=\frac{\alpha_1\alpha_2}{(\alpha_1+\alpha_2)^2},\qquad J(w)=w^{-k}\,\sB^{\{0,1\}}_{(k,k),\,k;\,I}\big(1,w^{-1}\big),
\end{gather*}
the stacked Gamma $\Gst22{k,\ k}{k,\ k}$ being $1$. At $k=1$, $J(w)=\log w/(w-1)$ (Lemma~\hb{5.16}{34}(ii)) and
\[
\int x_B\,\frac{\log x_B}{x_B-1}\,\Omega_B=\sum_{c\ge1}-\partial_c\Gst21{c,\ c}{2c}=\sum_{j\ge0}2\,\Gst21{j+1,\ j+1}{2j+2}\Th{1}{\IZ}{j}{2j+1},
\]
the derivative of the Beta corner reading the head of the counting input between $j$ and $2j+1$ (Prop.~\hb{5.13}{32}),
which is the series of the three-mark value $P_3(1)=3L$ (\S11.4). Hence
\[
\langle1\rangle_1=(3L)^2=9\,\Tx{2}{(\Ilat{1+3\mathbb Z}-\Ilat{2+3\mathbb Z})}{0}^{\,2}.
\]
The one hinge term of the tree is the whole value, a product of the room's content with itself, and the coaction is the coaction of a square. (Alias, once: $9\,L(\chi_{-3},2)^2$.) The two levels for this list are checked in \texttt{two\_levels\_check.py}.

\subsection{The coupled node}\status{derived; Brown's coaction formula recalled; the boundary list and the open part corrected in \S14.4}

\begin{claim}\label{cl:nodeclosure}
At a node of width one, every left factor of the coaction of $P_F$ is a finite sum of terms of two kinds.
\begin{enumerate}[label=(\roman*)]
\item \emph{Faces.} A period of a face of the bunch family times a period of a face of the room family. The two corners $u\to0$ and $u\to\infty$ are the hinge term, the two levels of Claim~\ref{cl:twolevels}:
\[
F_B\cdot F_R\qquad\text{and}\qquad\big(F_B\,x_B^{-t}\big)\cdot\big(F_R\,y_R^{-t}\big).
\]
\item \emph{Pairings of the two laws.} With $'$ marking an insertion (a monomial on $B$, a monomial on $R$, integer shifts of $s,t,\kappa$),
\[
\iint F'_B\,F'_R\;J'(x_By_R)\,\Omega_B\,\Omega_R=\frac1{2\pi i}\int\hat J'(a)\,\varphi'_B(a)\,\varphi'_R(a)\,da,
\]
where $\varphi'_B(a)=\int F'_B\,x_B^{-a}\,\Omega_B$ and $\varphi'_R(a)=\int F'_R\,y_R^{-a}\,\Omega_R$ are the bunch and room families with the insertion, and $\hat J'$ is the Gamma kernel at the shifted data.
\end{enumerate}
The classes on $D_{B,R}=\{x_By_R=1\}$ are of kind (ii).
\end{claim}

\emph{Proof.}
\begin{itemize}
\item The coaction is $\Delta[H,\sigma,\omega]^{\mathfrak m}=\sum_i[H,\sigma,e_i]^{\mathfrak m}\otimes[H,e_i^\vee,\omega]^{\mathfrak{dr}}$ over a basis $(e_i)$ of the de Rham realization (recalled, Brown). So its left factors are the periods, on the same chain $\sigma$, of the basis classes.
\item Work in the separated coordinates $(\beta,\alpha_R,u)$ of Claim~\ref{cl:separation}. A class of the relative de Rham cohomology is a form on the open part together with forms on the boundary strata. Its period on $\sigma$ is the period of the open part plus the periods of the boundary forms on the pieces of $\partial\sigma$.
\item \emph{The open part.} The form has poles only on the letters, so it is a finite sum of monomials times integer powers of the letters. By Claim~\ref{cl:twolevels} each term reads
\[
F'_B\,F'_R\;u^{\kappa'-1}(1+u)^{-s'}(1+x_By_Ru)^{-t'}.
\]
Integrating $u$ gives kind (ii): Claim~\ref{cl:separation} at the shifted data, continued in the data where needed.
\item \emph{The boundary.} The strata are the faces of $\mathbb P(\mathbb R^B_{>0})$ times $\mathbb P(\mathbb R^R_{>0})$, the faces of $\mathbb P(\mathbb R^R_{>0})$ times $\mathbb P(\mathbb R^B_{>0})$, and $u=0$, $u=\infty$. Each is a product of a face on the bunch side and a face on the room side, so its periods are kind (i). As $u\to0$ the integrand tends to $u^{\kappa-1}F_BF_R$; as $u\to\infty$ it tends to $u^{\kappa-s-t-1}\,(F_Bx_B^{-t})(F_Ry_R^{-t})$.
\item $D_{B,R}$ is not a stratum of the domain. It is where the two punctures of the fibre meet, so its classes are classes of the open part, of kind (ii). $\square$
\end{itemize}

\emph{What the coupled node does.} At a dual pair a pairing is a product (Claim~\ref{cl:dualpair}), and the coaction factors through the two sides. At a coupled node a pairing need not be a product of the two sides' periods. For example, the hinge that proves Claim~\ref{cl:row4} pairs two Beta-corner laws, whose moments at integer data are rational, and the pairing is $7\,\Tx3{\IZ}0$. That is content created at the hinge: the cofactor of \S9.2, carried by the pairing of the two laws. So at every node of width one the coaction is carried by three things, all of which the tree holds: the hinge term, the faces of the two sides (handled by their own nodes, down the tree), and the pairings of the two laws.

\subsection{Any width}\status{the boundary step as corrected in \S14.4; the generation made precise in \S14}

At a bunch of width two or more the two levels are replaced by the patterns. With $\alpha_R$ fixed, the bunch integral is the port (Cor.~\hb{15.17}{92}): a Breta of the children's coefficients $\Lambda_{j,(p)}(z_R)$. The proof of Claim~\ref{cl:nodeclosure} then runs unchanged:
\begin{itemize}
\item a form on the open part has as its period on $\sigma$ an integral over the room of a Breta of the children, with an insertion;
\item a boundary class is a face of the bunch side times a face of the room side.
\end{itemize}
By induction down the tree, every left factor of the coaction of a list's period is built from the atoms by the tree's own operations: faces, hinge terms and pairings. (\S14 proves this in general and says what it carries: read without weights it is automatic and empty (\S14.2); its content is the weight-graded faces (\S14.2), the hinge form (\S14.4) and the dual pairs (\S14.5).)

\section{The coaction on the content}

Claim~\ref{cl:nodeclosure} says where the factors of the coaction come from. This entry writes the coaction on the content itself:
\begin{itemize}
\item a trail is a chain on the Laplace side (\S13.1);
\item a chain read between two singular points is again content (\S13.2);
\item the coaction cuts the chain where its letters are singular (\S13.3);
\item at tier one the singular points are the frequencies of the lattices;
\item at tier two the live input's letter is singular at the critical values of the pencils, and a kept letter opens into the solutions of its ray (\S13.4, \S13.5);
\item the critical values read the primes at which the letters degenerate; the content can see fewer of them (\S13.6--\S13.8).
\end{itemize}
Throughout, a chain is
\[
\mathrm I(a_0;a_1,\dots,a_w;a_{w+1})=\int_{a_0<t_1<\dots<t_w<a_{w+1}}\frac{dt_1}{t_1-a_1}\cdots\frac{dt_w}{t_w-a_w},
\]
on the straight path from $a_0$ to $a_{w+1}$, with $\mathrm I(a;\varnothing;b)=1$ and $\mathrm I(a;w;a)=0$ for $w\neq\varnothing$.

\subsection{Trails are chains}\status{derived; checked at depth two}

\begin{claim}\label{cl:trailchain}
Let $\zeta_i=e^{2\pi ia_i}$ be frequencies and $\nu_i\ge1$ indices, with $(\nu_d,\zeta_d)\neq(1,1)$. Then
\begin{multline*}
\big(\sT_{\boldsymbol\nu;I_d}[y]\ e^{2\pi i\ip{\mathbf a}{y}}\,\Ilat{\mathbb Z^d}\big)\big(\{0<y_1<\dots<y_d\}\big)\\
=(-1)^d\,\mathrm I\big(0;\ c_1^{-1},0^{\nu_1-1},\ c_2^{-1},0^{\nu_2-1},\ \dots,\ c_d^{-1},0^{\nu_d-1};\ 1\big),
\end{multline*}
where the letters are listed in order and $c_i=\zeta_i\zeta_{i+1}\cdots\zeta_d$.
\end{claim}

\emph{Proof.} Put $z=e^{-U}$, the variable of the Laplace side (\S9.3), and let $S(z)$ be the trail with $z^{n_d}$ inserted at its last step, so that $S(0)=0$ and the trail is $S(1)$.
\begin{itemize}
\item If $\nu_d\ge2$, then $z\,\partial_zS$ is the same trail with $\nu_d$ lowered by one.
\item If $\nu_d=1$, summing the geometric tail over $n_d>n_{d-1}$ gives
\[
\partial_zS=\frac{\zeta_d}{1-\zeta_dz}\,S^\flat(z),
\]
where $S^\flat$ is the trail of depth $d-1$ whose last frequency is $\zeta_{d-1}\zeta_d$.
\end{itemize}
Since $\zeta\,dz/(1-\zeta z)=-dz/(z-\zeta^{-1})$, integrating from $0$ to $1$ one letter at a time gives the chain. $\square$

Cosets reduce to frequencies by rescaling the lattice (Thm.~\hb{5.12}{32}). A cut $N\in\mathbb Z$ reduces to the cut $0$ plus a head (Prop.~\hb{5.3}{28}), and a head is a finite sum in $\mathbb Q(\mu_M)$, of weight zero. So the content at level $M$ is spanned by chains whose letters lie in $\{0\}\cup\mu_M$, together with heads. Checked at depth two, $\boldsymbol\nu\in\{(1,1),(2,1),(1,2)\}$, against an independent term-by-term expansion of the chain: \texttt{trail\_chain\_check.py}.

\subsection{A trail read between two singular points}\status{derived}

The cut slot of a Trailmoment, the region $\{N<y_1<\cdots\}$, is on the Laplace side the path from $0$ to $1$. The new verb reads the same letters on the path between two singular points $s,s'\in\{0\}\cup\mu_M$:
\[
\mathrm I(s;a_1,\dots,a_w;s'),\qquad a_i\in\{0\}\cup\mu_M .
\]
It is regularized with the tangent at each end pointing along the path, and by the shuffle. On the Trailmoment side this sets the counting input at index one, cut $0$, to zero.

\begin{claim}\label{cl:betweenclosed}
Every chain read between two singular points is a polynomial, with coefficients in $\mathbb Q(\mu_M)$, in content of level $M$.
\end{claim}

\emph{Proof.}
\begin{itemize}
\item \emph{Rotation.} For $s'\in\mu_M$, $t\mapsto s't$ gives $\mathrm I(0;a;s')=\mathrm I(0;a/s';1)$, and $a/s'\in\{0\}\cup\mu_M$.
\item \emph{Reversal.} $\mathrm I(s';a_1,\dots,a_w;s)=(-1)^w\,\mathrm I(s;a_w,\dots,a_1;s')$.
\item \emph{Path composition.} $\mathrm I(s;a;s')=\sum_{a=uv}\mathrm I(s;u;0)\,\mathrm I(0;v;s')$.
\end{itemize}
This reduces every chain to products of chains from $0$ to $1$. A chain from $0$ to $1$ whose first letter is not $0$ and whose last letter is not $1$ parses as blocks $(c^{-1},0,\dots,0)$, so it is a trail by Claim~\ref{cl:trailchain}, with frequencies $\zeta_i=c_i/c_{i+1}\in\mu_M$. The shuffle regularization writes the remaining chains through these. $\square$

\subsection{The cutting rule}\status{the rule and its coassociativity derived; checked on 21844 chains; its compatibility with the shuffle and with path composition recalled}

The rule is written with the de Rham factor on the right, as in \S12.5:
\[
\Delta\,\mathrm I^{\mathfrak m}(a_0;a_1,\dots,a_w;a_{w+1})=\sum_{S}\ \mathrm I^{\mathfrak m}(a_0;a_S;a_{w+1})\ \otimes\ \prod_{p=0}^{|S|}\mathrm I^{\mathfrak{dr}}\big(a_{s_p};a_{s_p+1},\dots,a_{s_{p+1}-1};a_{s_{p+1}}\big).
\]
Here the sum runs over the kept subsequences $S$ of $1,\dots,w$, and $s_0=0<s_1<\dots<s_{|S|}<s_{|S|+1}=w+1$ lists $S$ together with the two ends.

In the slots of the Trailmoment:
\begin{itemize}
\item the letters kept form the left factor: a trail with the same cut, read on the same path;
\item each run of dropped letters forms a right factor: the verb of \S13.2, read between the singular points of the two kept letters that bound it, on the de Rham side.
\end{itemize}

\begin{claim}\label{cl:cutting}
\begin{enumerate}[label=(\roman*)]
\item The rule is coassociative.
\item It commutes with rotation and reversal.
\item It is compatible with the shuffle product and with path composition. So it is well defined on chains modulo the relations of \S13.2, and it is the motivic coaction there (recalled; alias below).
\end{enumerate}
\end{claim}

\emph{Proof of (i).} Both $(\Delta\otimes1)\Delta$ and $(1\otimes\Delta)\Delta$ are sums over nested kept sets $T\subseteq S$.
\begin{itemize}
\item In the first, $S$ is kept and then $T$ is kept inside $S$.
\item In the second, $T$ is kept, and $S\setminus T$ is what is kept inside each gap of $T$.
\end{itemize}
In both, the three factors are the same: the chain of $T$; for each gap of $T$, the chain of the positions of $S$ inside it, read between its ends; and for each gap of $S$, its own chain. (ii) is read off the formula. $\square$

Checked by machine: (i) holds on all $21844$ chains of length at most six over four letters, with four pairs of ends (\texttt{cutting\_rule\_check.py}).

\emph{Examples.}
\begin{itemize}
\item Depth one at frequency $0$, index $3$: $\Tx3{\IZ}0=-\mathrm I(0;1,0,0;1)$. Every proper kept set gives one of two things: a kept chain $\mathrm I(0;1;1)$ or $\mathrm I(0;0,\dots;1)$, zero by the regularization; or a gap with equal ends or of the form $\mathrm I(0;0;1)$, also zero. So the trail has no proper cut. The four-mark content $7\,\Tx3{\IZ}0$ is therefore primitive, as the atom of \S11.2 requires.
\item At the frequency $\frac13$ the one possible cross term pairs the kept letter $\mathrm I^{\mathfrak m}(0;1;\omega)=\log(1-\omega)$ with the gap $\mathrm I^{\mathfrak{dr}}(1;0;\omega)=\log^{\mathfrak{dr}}\omega$. That gap is the constant letter read between two singular points that are roots of unity, and it vanishes on the de Rham side. At tier one, cuts between frequencies leave no logarithm. That is what changes at tier two (\S13.6).
\end{itemize}

\subsection{The letter of a live input}\status{derived; computed for the three-mark room}

At a coupled node the cofactor of \S9.2 is a trail of the live input
\[
\psi(j)=\xi_j\cdot P_{F^B}(j)\cdot c_i(j),
\]
the stacked Gamma's Laurent data times the bunch's value times the room's coordinate. Each factor satisfies a recurrence with polynomial coefficients:
\begin{itemize}
\item the Laurent data, a Gamma ratio, satisfy one of order one;
\item the two sides' rays satisfy theirs by Thm.~\hb{35.19}{313}(ii).
\end{itemize}
So the letter $\sum_j\psi(j)z^j$ satisfies a linear differential equation in $z$. Its singular points are the reciprocals of the symbols' roots, multiplied together: the products of the two sides' critical values in the hinge's variables. (Alias, once: Hadamard's multiplication theorem.)

\emph{The three-mark room, computed.} From the coordinates $P_3(c)=r_c+s_cL$ for $c\le26$, both coordinates satisfy one recurrence, unique with coefficients of degree two:
\[
c^2\,P_3(c)-(10c^2+10c+3)\,P_3(c+1)+9(c+1)^2\,P_3(c+2)=0 .
\]
Its symbol is $9\lambda^2-10\lambda+1=(9\lambda-1)(\lambda-1)$.
\begin{itemize}
\item The value follows $\lambda=\frac19$, the orthant fibre $9$.
\item The coordinates follow $\lambda=1$, the critical value $1$ of \S11.5.
\end{itemize}
So the letters $\sum_cs_cq^{c-1}$ and $\sum_cr_cq^{c-1}$ are singular exactly at the room's critical values $q=1$ and $q=9$ (\texttt{tier\_two\_letter.py}).

\subsection{The cutting rule at tier two}\status{derived; Brown's formula recalled}

\begin{claim}\label{cl:tiertwo}
Let a chain carry one live letter $G$, whose equation has order $h$ and solutions $g_1,\dots,g_h$, among rational letters. Its coaction is the rule of \S13.3 with two changes.
\begin{enumerate}[label=(\roman*)]
\item A kept live letter opens into its solutions: the term with $G$ kept becomes the sum over $k$ of the chain with $g_k$ in place of $G$, tensored with the de Rham coordinate of $G$ on $g_k$.
\item Dropped runs are read between singular points, which now include the singular points of $G$: the products of critical values.
\end{enumerate}
\end{claim}

\emph{Proof.}
\begin{itemize}
\item $G$ solves, with rational data at $z=0$, the Gauss--Manin equation of the node's family over the $z$-line: the two sides' pencils (Thm.~\hb{35.10}{309}) joined by the hinge (Claim~\ref{cl:separation}).
\item So the chain is an entry of the transport of that system along the path, from the fibre over $0$ to the fibre over the end. That is a period of the relative cohomology of the family over the path, relative to the fibres over the two ends.
\item Filter it by the rational letters. Each graded piece is the rational letters tensored with the de Rham realization of the local system of $G$, whose basis is $g_1,\dots,g_h$. The boundary pieces at the singular points carry that system's fibres there.
\item Brown's formula $\Delta[H,\sigma,\omega]^{\mathfrak m}=\sum_i[H,\sigma,e_i]^{\mathfrak m}\otimes[H,e_i^\vee,\omega]^{\mathfrak{dr}}$ in this basis gives (i) from the classes $g_k$ and (ii) from the boundary pieces.
\item At tier one the local system has rank one and is constant, so (i) is the identity and the singular points are the frequencies: the rule of \S13.3. $\square$
\end{itemize}

\emph{The row of \S9.2 is this rule.} The room's coordinates $c_i(j)$, $i<h$, are the fundamental solutions of its ray, and the cofactor of the $i$-th room value is the trail of the live letter built on $c_i$. So
\[
P_F=\sum_{i<h}P_{F^R}(j_0+i)\cdot\Big[\sum_j c_i(j)\,\xi_j\,P_{F^B}(j)\Big]
\]
has the shape of (i): the room's $h$ values against the kept live letter opened into its $h$ solutions.

\subsection{The level read off the pencil}\status{matches at three marks; matches at five (recalled); fails at four (\S13.7)}

A dropped run of constant letters read between two singular points $s,s'$ is $\log(s'/s)$.
\begin{itemize}
\item Between two frequencies the ratio is a root of unity, and the logarithm vanishes on the de Rham side (\S13.3).
\item Between two critical values it does not vanish.
\end{itemize}
For $\binom{\cE}1$ on $n$ marks the critical values are $(n-2j)^2$. Their primes, set against the level of the content:
\[
\begin{array}{c|c|c|l}
n&\text{critical values}&\text{primes}&\text{level of the content}\\\hline
3&9,\ 1&3&3:\ \text{frequency }\tfrac13\text{ of }L\ (\S11.4)\\
4&16,\ 4,\ 0&2&1:\ 7\,\Tx3{\IZ}0\ (\S13.7)\\
5&25,\ 9,\ 1&3,\ 5&15:\ \text{the recalled modular value }(\S9.5)
\end{array}
\]
For three marks the two readings meet directly. The constant letter read between the room's critical values $1$ and $9$ is $\log 9=2\log 3$, the logarithm that $\Lambda$ carries (\S11.4).

The statement entered here earlier, ``the primes of the content of a list are the primes of the ratios of the critical values of its pencils'', is false. Its four-mark row was entered as the coset $\frac12$, which collapses (\S13.7). \S13.8 proves what the critical values do read.

\subsection{The equality fails at four marks}\status{derived}

By Claims~\ref{cl:row4} and~\ref{cl:port4}, every value of the four-mark centre ray is
\[
\langle1\rangle_k=A_k\cdot 7\,\Tx3{\IZ}0-B_k,\qquad A_k,B_k\in\mathbb Q .
\]
The coset $\frac12$ of Claim~\ref{cl:row4} collapses by Thm.~\hb{5.12}{32}:
\[
\Tx3{\Ilat{\frac12+\mathbb Z}}0=\sum_{l\ge0}\frac8{(2l+1)^3}=7\,\Tx3{\IZ}0 .
\]
So the content of the whole ray is $\mathbb Q+\mathbb Q\,\Tx3{\IZ}0$, of level one, at every $k$. The nonzero critical values $16$ and $4$ carry the prime $2$. Hence the primes of the content are not, in general, the primes of the critical values. (Alias, once: $\zeta^{\mathfrak m}(3)$ is unramified.)

\subsection{What the critical values read}\status{derived}

\begin{claim}\label{cl:criticalprime}
Let $\alpha^\star$ be a critical point of $L$ on the torus with critical value $\lambda=L(\alpha^\star)\neq0,\infty$. Let $\mathfrak p$ be a prime, of a number field holding the coordinates of $\alpha^\star$, at which those coordinates are units. If $\mathfrak p\mid\lambda$, then at the reduction $\bar\alpha^\star$ the letters' divisor $\{K\Psi=0\}$ fails to have smooth components meeting transversally: either a letter is singular at $\bar\alpha^\star$, or the two letters meet there non-transversally.
\end{claim}

\emph{Proof.} Write $D_i=\alpha_i\partial_{\alpha_i}$. The critical equations of $L=K\Psi/\prod_e\alpha_e$ are
\[
\frac{D_iK}K+\frac{D_i\Psi}\Psi=1\qquad(i\in\cE),
\]
equivalently $\Psi\,D_iK+K\,D_i\Psi=K\Psi$. Put $K(\alpha^\star)=\pi^ak$ and $\Psi(\alpha^\star)=\pi^b\psi$ with $k,\psi$ units, so $a+b=v_{\mathfrak p}(\lambda)>0$, since $\prod\alpha^\star_e$ is a unit.
\begin{itemize}
\item If $a<b$, divide by $\pi^a$: $k\,D_i\Psi\equiv0$ modulo $\mathfrak p$ for every $i$. Since $\Psi(\bar\alpha^\star)=0$ and the $\alpha^\star_i$ are units, $\{\Psi=0\}$ is singular at $\bar\alpha^\star$.
\item If $a>b$, likewise $\{K=0\}$ is singular at $\bar\alpha^\star$.
\item If $a=b$, divide by $\pi^a$: $\psi\,D_iK+k\,D_i\Psi\equiv0$, so the two gradients are proportional at $\bar\alpha^\star$, and the letters meet non-transversally (or one is singular). $\square$
\end{itemize}

\emph{Three marks.} The critical point $(1,1,1)$ has value $9$. Modulo $3$ the line $e_1=0$ becomes tangent to the conic $e_2=0$ there: the two points $(1:\omega:\omega^2)$ and $(1:\omega^2:\omega)$, where the letters meet transversally in characteristic zero, both reduce to $(1:1:1)$.

\emph{Four marks.} The critical point $(1,1,1,1)$ has value $16$. Modulo $2$ the three vertices $(1,-1,-1,1)$, $(1,-1,1,-1)$, $(1,1,-1,-1)$ of the triangle $e_1\cap e_3$ (\S11.5) collapse onto it.

So the critical values read the primes at which the letters degenerate, and the place: the reduced critical point. The content is what the orthant cycle sees of that geometry, and it can see less. At four marks the degeneration at $2$ is there, and the content $7\,\Tx3{\IZ}0$ does not see it.

\emph{Where the prime goes when the content does not see it: the frame.} When the roots of the symbol of the ray's recurrence are the reciprocals of the nonzero critical values (Thm.~\hb{35.19}{313}(ii), Cor.~\hb{35.6}{308}), the ratio of the symbol's extreme coefficients is their product. That is the frame of \S8.7:
\[
\text{three marks: }9\lambda^2-10\lambda+1\ \Rightarrow\ 9=9\cdot1;\qquad\text{four marks: }64\lambda^2-20\lambda+1\ \Rightarrow\ 64=16\cdot4 .
\]
At four marks the prime $2$ of the critical values sits in the denominators, $64^{k-1}A_k\in7\mathbb Z$ (Claim~\ref{cl:port4}), not in the content. So the critical primes are the frame's primes. Whether the content carries them is the separate question of this subsection.

\emph{Status of the two remaining inclusions.}
\begin{enumerate}[label=(\alph*)]
\item The primes of the content lie among the primes of bad reduction of the list's motive (recalled: a motive with good reduction at $p$ has periods unramified at $p$).
\item Whether every prime of bad reduction of the letters is a critical prime of $L$ is not derived. A generic master function sees every change of the Euler characteristic through its critical points. $L$ has the fixed exponents $(1,1,-1)$, and nothing derived here excludes a degeneration that its critical points miss.
\end{enumerate}

Alias, once:
\begin{itemize}
\item Claim~\ref{cl:trailchain} is the iterated-integral form of multiple polylogarithms at roots of unity.
\item \S13.3 is Goncharov's coaction, motivic by Goncharov and Brown.
\item Claim~\ref{cl:tiertwo} is Brown's formula for periods with coefficients in a local system.
\end{itemize}

\section{The tree and the coaction, in general}

\S12.6 ended with the general statement: down the tree, every left factor of the coaction of a list's period is built from the atoms by faces, hinge terms and pairings. This entry proves it for every family of lists at integer data where the period converges, read at any of its realizations, and says exactly what it carries.

The proof, written for a reader outside the family, is the note \texttt{tree\_coaction\_note.tex}; its Theorems A, B, C, D are Claims~\ref{cl:weights}, \ref{cl:hinge14}, \ref{cl:pinned14}, \ref{cl:reach14} below, its Corollary~3.1 is Claim~\ref{cl:serif14}, and its Lemma~4.4 is Claim~\ref{cl:onehinge14}. Two corrections to Claim~\ref{cl:nodeclosure} come out of it (\S14.4).

\subsection{The motive of a family}\status{derived; the realization category, Brown's coaction formula and his affine model recalled}

At integer data with every wall $\kappa_S\ge1$, the period of Thm.~\hb{8.10}{56} is
\[
P_F=\int_{\sigma_{\cE}}\alpha^{\mvec}\prod_jK_j^{-s_j}\Psi_j^{-t_j}\,\frac{\Omega_{\cE}}{\alpha^{\one}} .
\]
The model:
\begin{itemize}
\item blow the projective space of the simplex up along every coordinate subspace, so that each face $\alpha_S\to0$, $\varnothing\neq S\subsetneq\cE$, becomes one boundary divisor $D_S$;
\item remove the zero loci of the letters;
\item remove also the letters of the lists $\binom J1$, $J\subseteq\cE$, $|J|\ge2$, whose $K$ is $\alpha_J=\sum_{j\in J}\alpha_j$.
\end{itemize}
Then
\[
P_F^{\mathfrak m}=\big[H^{n-1}(X_{\cE}\setminus Y_F,\ D\setminus Y_F),\ \widetilde\sigma_{\cE},\ \omega_F\big]^{\mathfrak m}.
\]
Three facts are read off one chart computation:
\begin{itemize}
\item the closed cell misses every letter;
\item $\omega_F$ vanishes to order exactly $\kappa_S-1$ along $D_S$, so the wall is the order;
\item with the $\binom J1$ adjoined the complement is affine.
\end{itemize}

\subsection{Faces are minors; one form carries every period}\status{derived; Brown's weight-graded face maps recalled}

Along the face $\alpha_S\to0$ the two letters of a table $\cB$, Kirchhoff $K$ (bases) and Symanzik $\Psi$ (cobases), factor through minors; on the wall these are the faces of Thm.~\hb{16.23}{100}:
\[
K\ \to\ K_{\cB/(\cE\setminus S)}(\alpha_S)\,K_{\cB\setminus S}(\alpha_{\cE\setminus S}),\qquad \Psi\ \to\ \Psi_{\cB|S}(\alpha_S)\,\Psi_{\cB/S}(\alpha_{\cE\setminus S}).
\]
Along a flag $S_1\subsetneq\cdots\subsetneq S_\ell$ the block $Z_i=S_i\setminus S_{i-1}$ carries
\[
K_{\cB\setminus S_{i-1}/(\cE\setminus S_i)}\qquad\text{and}\qquad\Psi_{\cB/S_{i-1}\setminus(\cE\setminus S_i)},
\]
together with the $\alpha_J$, $J\subseteq Z_i$.

\begin{claim}\label{cl:oneform}
On the affine complement, every class of the motive is the class of one global form. Every period of the motive on the cell, every left factor of the coaction of $P_F^{\mathfrak m}$ among them, is the period of a single form on the whole simplex.
\end{claim}

\emph{Proof.} A class is a tuple of top forms on the strata. Take the deepest nonzero entry. On the stratum one level up, which is affine, a form of one degree less restricts to any prescribed top forms on its boundary components; subtracting its coboundary removes the deepest entry. Induct. $\square$

So ``built from faces, hinge terms and pairings'' holds with no face at all: read without weights it carries nothing. The faces carry information through the weights.

\begin{claim}\label{cl:weights}
For $k<n-1$, a left factor of weight $\le k$ is a sum of products of periods of forms on the blocks of the flags whose stratum has dimension $k$. Only blocks with at most $k+1$ marks enter.
\end{claim}

\emph{Proof.} The weight-$\le k$ part comes from the $k$-dimensional strata (Brown's weight-graded face maps, which need the strata affine). Claim~\ref{cl:oneform} on each such stratum relative to its boundary gives one form there. Each stratum is a product of the blocks' complements, so the form is a sum of products. $\square$

\begin{claim}\label{cl:serif14}
Every period of a form $\eta$ on a block $Z$ is a rational combination of constant terms of continued periods of families of lists on $Z$. On the torus $\eta=\sum_tc_t\,\omega_{F_t}$, $c_t\in\mathbb Q$: the tables of $F_t$ are the lists of the block's letters and the $\binom J1$, its exponents are the powers of the letters, and its masses $\mathbf a_t\in\mathbb Z^Z$ lie on the balance wall. For a balanced direction $\mathbf v$ of the masses with $v_S\neq0$ for every nonempty proper $S$,
\[
\int_{\widetilde\sigma_Z}\eta=\sum_tc_t\,\CT_{\varepsilon=0}P_{F_t}(\mathbf a_t+\varepsilon\mathbf v;\mathbf s_t,\mathbf t_t).
\]
So the left factors of Claim~\ref{cl:weights} are, as numbers, sums of products over the blocks of such values: the coaction stays in the serif family.
\end{claim}

\emph{Proof.} On the torus the complement of the letters is a principal open set, so a regular function there is a Laurent polynomial over a product of powers of the letters, and its monomials give the families. With $\alpha^{\mathbf x}$ inserted, $x_{\cE}=0$, the integral of the form is holomorphic near $\mathbf x=0$. On each sector cube it is the sum of the cube integrals of the terms (Lemma~\hb{8.6}{55}), which are the sector forms of their periods and continue them (Thm.~\hb{8.7}{55}(ii),(iii)). Along $\varepsilon\mathbf v$ no term lies on a pole hyperplane, and the constant term is linear. $\square$

Each continued period is a value of the serif family: it is the residue on the balance wall of the interest, a Gramma over the Gammas of the exponents (Thm.~\hb{8.3}{54}, Thm.~\hb{8.10}{56}), and where the masses and the walls are positive it is a stacked Gamma times a Trailmoment of the joint law of the letters (Prop.~\hb{8.11}{56}). The constant term along a balanced direction is the operation of the pseudointegral (Def.~\hb{16.8}{96}(iii)), read on another ray.

\subsection{The blocks sit on the tree}\status{derived}

Descend the tree from the root. A block lands in exactly one of three places:
\begin{itemize}
\item inside a bunch: a \emph{bunch piece};
\item meeting both the bunch and the room of a node, or two childcules of a split: a \emph{node piece};
\item inside a leaf: a \emph{leaf piece}.
\end{itemize}

\begin{claim}\label{cl:children14}
Let $\cB$ be a table with bunch $B$ and room $R$, and $B=X\sqcup Y$. The minor $\cB\setminus X/Y$ on $R$ is the child at some pattern of $B$ (Lemma~\hb{12.2}{70}).
\end{claim}

\emph{Proof.} Contracting $Y$ is contracting a maximal independent $Y'\subseteq Y$ and deleting the rest. A mark $b\in B\setminus Y'$ outside the closure of $R\cup Y'$ is a coloop of the one table on $R+b$ of Prop.~\hb{12.5}{70}, so deleting it and contracting it agree. Grow $Y'$ this way to $Y''$ with $R\cup Y''$ spanning. Then $Y''$ is the trace on $B$ of a basis, a pattern. $\square$

So the letters of a block attached at a node are Foreman polynomials of minors of the tables at that node. The peel carries everything the room's blocks need: going down the tree loses nothing the coaction can see.

\subsection{The node at any width}\status{derived; the levels checked at every width on $1340$ letters, the discriminants at width two computed}

The width of a set is the width of its window (Def.~\hb{7.10}{49}, Thm.~\hb{7.11}{49}): every basis meets $B$ in between $d(B)$ and $d(B)+c(B)$ marks, every number in between occurs, and
\[
c(B)=r(B)+r(R)-r .
\]
Nothing in Claims~\ref{cl:oneform}--\ref{cl:children14} uses it, and nothing in the tree requires it. The Foreman bunch is a class with the fewest marks (Def.~\hb{13.6}{78}), and its definition places no condition on the width; \S14.6 has a Foreman bunch of width two. Width enters only in how the bunch and the room meet at a node.

Minors do not raise the width. For a minor $\cB'=\cB\setminus X/Y$ on a block $Z$ meeting both $B$ and $R$, $X\sqcup Y=\cE\setminus Z$,
\[
r_{\cB'}(Z\cap B)+r_{\cB'}(Z\cap R)-r_{\cB'}(Z)\ \le\ r_{\cB}(B)+r_{\cB}(R)-r_{\cB} .
\]
So a node piece at a node of width $c$ has width at most $c$ at its own split. The lists $\binom J1$ have width at most one at every split.

\begin{claim}\label{cl:hinge14}
At a node of width $c$, read the bunch side of a node piece $Z$ in scale and shape, $\alpha_{Z\cap B}=\rho\,u$ (Lemma~\hb{2.7}{14}), and the room side on its simplex, $\alpha_{Z\cap R}=u'$. Every letter is $\rho^{d}$ times its letter on the scale (Def.~\hb{18.17}{122}),
\begin{gather*}
P(\rho;u,u')=\sum_{k=d}^{d+c}\rho^{\,k-d}K^{(k)}(u,u')=K^{(d+c)}(u,u')\prod_{i=1}^{c}\big(\rho+\xi_i(u,u')\big),\\
K^{(k)}=\sum_{|p|=k}u^{p}\,K_{(p)}(u'),
\end{gather*}
whose coefficients are the level letters (Def.~\hb{12.15}{73}), each a sum over the patterns of its level of a monomial of the bunch against a child on the room, the two extreme ones products of a bunch letter and a room letter (Thm.~\hb{12.17}{73}), and whose roots are the negatives of the Skeleton values (Def.~\hb{12.9}{71}). Every period of a form on the piece is
\[
\iint\sum_tA_t(u)\,B_t(u')\ \CT_{\kappa=a_t}\Gst2{|J_t|}{\kappa,\ |\mathbf w_t|-\kappa}{\mathbf w_t}\,\sB^{\{J_t\}}_{\mathbf w_t,\,|\mathbf w_t|-\kappa;\,I_{J_t}}\big(\boldsymbol\xi(u,u')\big)\,d\omega(u)\,d\omega(u'),
\]
with $A_t$ and $B_t$ rational on the bunch and on the room: the scale integral of each term is one Breta of the joint spectrum (Prop.~\hb{15.1}{87}, Thm.~\hb{15.2}{87}), its letter exponents those of the letters on their Skeleton values, read at its finite part (Thm.~\hb{1.6}{12}). The bunch and the room meet only through the joint spectrum. At a root this describes the period of every left factor of the coaction of $P_F^{\mathfrak m}$ (Claim~\ref{cl:oneform} with the block $\cE$).
\end{claim}

\emph{Proof.} Expand the form over the torus and read it in scale and shape. Each letter contributes $\rho^{d}$, its top level, which is separable, and its roots. The scale integral of the sum converges at almost every point and is its own finite part (Prop.~\hb{1.5}{12}); finite parts are linear (Def.~\hb{1.4}{11}), and the finite part of each term is the constant term of the value of Prop.~\hb{15.1}{87} at its index (Thm.~\hb{1.6}{12}). $\square$

At integer data the reads are finite algebra (Ch.~\hb{37}{319}): the scale exchange reduces every read, with coefficients rational in the level letters, to the $D$ masters of the letters on the scale, $D$ the number of distinct points of the joint spectrum, without finding a root (Thm.~\hb{37.7}{324}), and at a corner the $3D$ corner numbers of the masters determine every convergent read (Thm.~\hb{37.10}{325}).

For the period of the family itself the claim is Thm.~\hb{20.13}{168} with its scale integral closed into one Breta (Thm.~\hb{15.2}{87}). Opened on one share per term of a level factorization, it is the gluing across a bunch of any width (Thm.~\hb{21.13}{178}): a fibre integral of stacked Gammas against the periods of the bunch word and the room word.

\emph{Width one: the bunch is one mark to its room} (Thm.~\hb{18.6}{116}). There are two levels, the extreme ones, so each is one station, a product of a bunch letter and a room letter (Thm.~\hb{12.17}{73}):
\[
K^{(d)}=K_{\cB/R}\,K_{\cB|R},\qquad K^{(d+1)}=K_{\cB|B}\,K_{\cB/B}.
\]
Each letter has one Skeleton value, the room's hinge value over the bunch's (Cor.~\hb{15.3}{88}, Cor.~\hb{18.7}{117}):
\[
\xi=\frac{K^{(d)}}{K^{(d+1)}}=\frac{\gamma_R(u')}{\gamma_B(u)},\qquad \gamma_B=\frac{K_{\cB|B}}{K_{\cB/R}},\quad \gamma_R=\frac{K_{\cB|R}}{K_{\cB/B}}.
\]
So the coefficients of the Bretas separate: on their share side (Cor.~\hb{15.4}{88}) they enter through the powers $\gamma_R(u')^{-b}\gamma_B(u)^{\,b}$, and wherever that side converges the integrals over the bunch and over the room separate under the share integral. For the family's own period this is the gluing across a bunch of width one (Thm.~\hb{18.13}{119}); for one list,
\[
P_F=\frac1{2\pi i}\int_{\re\sigma^d=\sigma^d_\ast}\Gst42{a^d,\ b^d,\ a^r,\ b^r}{s,\ t}\,P_{F^{\mathrm{bunch}}}(\mvec|_B)\,P_{F^{\mathrm{room}}}(\mvec|_R)\,d\sigma^d,
\]
with the even share entering through $(\varrho_B/\varrho_R)^{-\sigma^d}$ (Lemma~\hb{18.15}{121}). Claim~\ref{cl:separation} is this read before the share is opened: its kernel $J$ is the Breta of the two Skeleton values. When all the active letters have one and the same Skeleton value, the share is pinned and $P_F$ is the Beta corner $\Gst21{\kappa_B,\ 2\sigma-\kappa_B}{2\sigma}$ times the two periods (Cor.~\hb{18.14}{121}).

\emph{Width two and more.} The Skeleton values of a letter are, up to sign, the roots of its letter on the scale, a polynomial of degree $c$ whose extreme coefficients are separable and whose middle coefficients sum over the stations of their level (Def.~\hb{20.12}{168}). They need not be rational. On $\binom{\{1,2,3,4\}}2$ with $B=\{1,2\}$, $c=2$ and $d=0$:
\[
K(\rho\alpha_1,\rho\alpha_2,\alpha_3,\alpha_4)=\alpha_3\alpha_4+\rho\,(\alpha_1+\alpha_2)(\alpha_3+\alpha_4)+\rho^2\alpha_1\alpha_2 .
\]
Its discriminant is $12$ at the gauge, so it is not a square, and the two Skeleton values are not rational functions of the point. On the lists of \texttt{note\_lemmas\_check.py}, all $560$ letters of width two have a discriminant that is not a square (computed). For a list with a representation over $\mathbb R$ or $\mathbb C$ whose nonzero maximal minors all have absolute value one, Thm.~\hb{20.15}{169} writes
\[
K(\rho u,u')=\rho^{d}K_{\cB/R}(u)\,K_{\cB/B}(u')\det\big(\rho\,S_B(u)+S_R(u')\big),
\]
with positive definite $c\times c$ matrices, and the Skeleton values are the eigenvalues of $S_B^{-1/2}S_RS_B^{-1/2}$. At width $c$ the bunch and the room meet through one $c\times c$ matrix on each side; at width one the matrices are $\gamma_B$ and $\gamma_R$.

\emph{The two corrections to Claim~\ref{cl:nodeclosure}.}
\begin{enumerate}[label=(\arabic*)]
\item Its list of boundary strata omitted the faces $\alpha_S\to0$ with $S$ meeting both $B$ and $R$, other than $B$ and $R$. These are node pieces of the same node, of width at most its width, and Claim~\ref{cl:hinge14} describes them.
\item Its step ``a class of the relative cohomology is a form on the open part together with forms on the boundary strata'' needs the affine complement, hence the $\binom J1$. Once those are adjoined, the boundary forms can be removed altogether (Claim~\ref{cl:oneform}).
\end{enumerate}
With both corrections, the conclusion of \S12.6 stands, at every width, in the precise form of Claims~\ref{cl:weights}, \ref{cl:serif14} and~\ref{cl:hinge14}.

\subsection{Pinned cells, at any width; which pieces split off}\status{derived; the one-hinge criterion checked on $1340$ letters; the factorization checked as exact identities of forms on two substitutions with two cells each, one of them at a full lean with cells that are not dual pairs; read alike against dual pairs on $263$ pieces}

The first form of this entry factored $P_F^{\mathfrak m}$ at a bunch $B$ of width one with $M/B=(M|R)^\ast$: the period went to the cell of the bunch $M^+_B$ (Thm.~\hb{18.6}{116}) and to the room, and the whole room was read through one hinge value. That is where width one came from, and it came from the single hinge value, not from the cell.

\begin{claim}\label{cl:onehinge14}
Let $\cB$ be a list on $S\sqcup S'$. Its letter $K$ reads $S$ through one hinge value, $K=A(\alpha_S)\,P\big(\alpha_{S'},\gamma(\alpha_S)\big)$ with $A,\gamma$ homogeneous, $\gamma$ of degree one and $P$ a polynomial, exactly when $c(S)\le1$. The same holds for $\Psi$.
\end{claim}

\emph{Proof.} At $c\le1$ the two extreme levels give it, with $\gamma=K_{\cB|S}/K_{\cB/S'}$ the hinge value (Thm.~\hb{18.6}{116}). Conversely, scaling $\alpha_S$ by $\tau$ gives every level as $K^{(k)}=A\gamma^{k-a}P_{k-a}$, a product of a polynomial on $S$ and one on $S'$. The level is a sum over its patterns of a monomial against a child (Thm.~\hb{12.17}{73}), with $0/1$ coefficients, so all its children agree and $K^{(k)}=f_k(\alpha_S)C_k(\alpha_{S'})$, $f_k$ the sum of the monomials of the patterns of size $k$. Then $f_k$ is proportional to $A\gamma^{k-a}$, and at $c\ge2$ this gives $f_{d+1}^2\propto f_df_{d+2}$. That is impossible: $\alpha_S^{2p}$ occurs in $f_{d+1}^2$ for a pattern $p$ of size $d+1$, and in $f_df_{d+2}$ it would need $q+q'=2p$ with $|q|=d$ and $|q'|=d+2$, both sets, so $q=q'=p$. $\square$

So a set of width $c\ge2$ is read through no single function of its marks, and its Skeleton values are needed (Claim~\ref{cl:hinge14}). The remedy is to read the room through several hinge values, one for each cell. A table built by substitution, $\cB=\mathcal H\langle\mathcal A_1,\dots,\mathcal A_q\rangle$ with the marks of $L$ left in place (Def.~\hb{20.1}{163}), has each $R_i$ of width one, whose cell of the bunch is $\mathcal A_i$ (Prop.~\hb{20.3}{164}(i)), and its letters read the core at the hinge values of the cells (Thm.~\hb{20.4}{164}):
\[
K_{\cB}=\prod_iK_{\mathcal A_i/g_i}\ K_{\mathcal H}\big(\alpha_L,(\gamma^K_{\mathcal A_i})_i\big),\qquad \Psi_{\cB}=\prod_i\Psi_{\mathcal A_i\setminus g_i}\ \Psi_{\mathcal H}\big(\alpha_L,(\gamma^\Psi_{\mathcal A_i})_i\big).
\]
A bunch $B\subseteq L$ has in $\cB$ the width it has in the core, since the traces of the bases on $L$ are those of $\mathcal H$. That width is at most $|L\setminus B|+q$, and every width occurs.

\begin{claim}\label{cl:pinned14}
Let $F$ be a family of lists substituted table by table, with integer data, every proper wall positive, and the active letters reading every cell $R_i$ through one hinge value $\gamma_i$. Let $\kappa_i$ be the wall of the cell (Thm.~\hb{20.7}{165}) and $2\sigma=\sum_j(s_j+t_j)$. Then
\[
P_F^{\mathfrak m}=P^{\mathfrak m}_{F_{\mathcal H}}\prod_iP^{\mathfrak m}_{G_i},\qquad P^{\mathfrak m}_{F_i}=\Gst21{\kappa_i,\ 2\sigma-\kappa_i}{2\sigma}P^{\mathfrak m}_{G_i},
\]
every family convergent. Here:
\begin{itemize}
\item the core family $F_{\mathcal H}$ carries the masses $\mvec|_L$ and $\kappa_i$ on the hinge $h_i$;
\item the pinned family $G_i$ on $R_i$, with the form $\gamma_i^{-\kappa_i}\alpha^{\mvec|_{R_i}}\prod_jK_{\mathcal A_{ij}/g_i}^{-s_j}\Psi_{\mathcal A_{ij}\setminus g_i}^{-t_j}$, is the load family of Thm.~\hb{21.4}{172} at integer loads adding up to $\kappa_i$, none on an inactive letter, and its period is the pinned shape integral of Thm.~\hb{21.5}{173}; its form, period and motivic period do not depend on the loads;
\item $F_i$ is the cell family on $R_i\sqcup\{g_i\}$, with the mass $2\sigma-\kappa_i$ on the hinge.
\end{itemize}
The Galois conjugates of $P_F^{\mathfrak m}$ lie in the span of products of theirs. Nothing is assumed about the width of a bunch in $L$.
\end{claim}

\emph{Proof.} The map $[\alpha]\mapsto\big([\alpha_L:\gamma_1:\dots:\gamma_q],[\alpha_{R_1}],\dots,[\alpha_{R_q}]\big)$ is birational, a shear in logarithmic coordinates, and carries $\omega_F$ to the product of the forms: the scale of $R_i$ comes out at the power $\kappa_i$ (Thm.~\hb{20.7}{165}). Positivity and Tonelli give convergence, and the birational invariance of motivic periods (Lemma~8.1 of the note) gives the identity. For the cell family, read the cell as $\binom{\{g,h\}}1\langle g,\mathcal A\rangle$, with the core a Beta form on two marks whose motive is $H^1(\mathbb A^1,\{0,\infty\})\cong\mathbb Q(0)$. $\square$

As numbers this is Thm.~\hb{21.5}{173} and Cor.~\hb{21.6}{173}. The two instances:
\begin{itemize}
\item \emph{Dual pairs,} $\mathcal A_i/g_i=(\mathcal A_i\setminus g_i)^\ast$. Then $\gamma^K_{\mathcal A_i}=\gamma^\Psi_{\mathcal A_i}=K_{\cB|R_i}/\Psi_{\cB|R_i}$ whatever the exponents, and $G_i=(\cB|R_i;\ \kappa_i,\ s+t-\kappa_i)$ is one list, with $\kappa_i=m_{R_i}-(s+t)\big(r(R_i)-1\big)$.
\item \emph{One list at a full lean,} $(s,0)$ or $(0,t)$: one letter is active, every cell of every substitution is pinned, and along a hinge tree $P_F^{\mathfrak m}$ is a positive rational times the product of the motivic periods of the prime nodes, each at the masses of its marks and, on each marker, the wall of $F$ at the marks beyond its edge. This is Cor.~\hb{21.7}{174}, Thm.~\hb{25.6}{208} and Prop.~\hb{25.10}{209} made motivic; a list without prime nodes has a rational motivic period at its full leans.
\end{itemize}

On five marks (\S12.4) the list is $\binom{\{1,2,h\}}1\langle1,2,\binom{\{3,4,5,g\}}2\rangle$, one dual-pair cell, and both factors are the three-mark list at its centre, so $(3L)^2$ is an identity of motivic periods. On eight marks, $\binom{\{b_1,b_2,h_1,h_2\}}2\langle b_1,b_2,\mathcal A_1,\mathcal A_2\rangle$ with $\mathcal A_i=\binom{\{x_i,y_i,z_i,g_i\}}2$ has the bunch $\{b_1,b_2\}$ of width two, with irrational Skeleton values, and at the centre
\[
P^{\mathfrak m}=P^{\mathfrak m}_{\binom42}\cdot\big(P^{\mathfrak m}_{\binom32}\big)^2,\qquad P=P_{\binom42}\cdot\Big(3\,\Tx{2}{\big(\Ilat{1+3\mathbb Z}-\Ilat{2+3\mathbb Z}\big)}{0}\Big)^2,
\]
the bunch of width two going whole into the core $\binom42$, which is prime; its centre has the motivic period of $\binom32$ (\S14.6), so $P^{\mathfrak m}=\big(P^{\mathfrak m}_{\binom32}\big)^3$.

The cell of the bunch is not confined to width one: at the width $c$ it is the level cell $\cB^\natural_B$ on $B\sqcup G$, $|G|=c$, with the bases $p\sqcup F$, $p$ a pattern of excess $c-|F|$ (Lemma~\hb{27.23}{237}). On eight marks the core $\binom42$ is the level cell of $\{b_1,b_2\}$. What is confined to width one is a single hinge: a bunch is stored on one mark exactly at width one (Thm.~\hb{27.22}{237}), and Claim~\ref{cl:onehinge14} is the letter form of that theorem.

\emph{Which pieces split off} (derived). Read a bunch $S$ as a piece with the stations of its room, with its scale $\gamma^\rho_S$, seam vector $\boldsymbol\varrho^\rho_S$ and law $(\boldsymbol\varrho^\rho_S)_\ast\phi^\rho_S$ (Def.~\hb{27.1}{224}, Def.~\hb{27.5}{227}). Everything glued on it sees it through its wall and its law (Thm.~\hb{27.9}{230}, Thm.~\hb{27.18}{235}), and the law of one mark is a point mass.

\begin{claim}\label{cl:reach14}
For a bunch $S$ of a convergent family of lists with $s_j,t_j\ge0$, the law of $S$ is a point mass exactly when every table with an active letter has width at most one at $S$ and the active letters of those of width one read $S$ through one hinge value. Then $S$ acts on everything glued to it as one mark, and with width one and $|S|\ge2$ it is a pinned cell (Claim~\ref{cl:pinned14}). For one list with both letters active and $\cB|S$, $\cB/S'$ connected, this is width one with $\cB/S'=(\cB|S)^\ast$.
\end{claim}

\emph{Proof.} The scale form has a positive density, so the law is a point mass exactly when the seam vector is constant. If it is, every active letter is its extreme excess sum times a polynomial in $\gamma^\rho_S$ with room coefficients (Prop.~\hb{27.2}{225}(iii)), so it reads $S$ through $\gamma^\rho_S$, and Claim~\ref{cl:onehinge14} gives width at most one. At width one the coordinates are the hinge values over $\gamma^\rho_S$ (Prop.~\hb{27.10}{231}(ii)), constant exactly when the hinge values are proportional, and proportional letters of content one are equal. For one list, $K_{\cB|S}\Psi_{\cB|S}=K_{\cB/S'}\Psi_{\cB/S'}$; a connected list on two marks or more has an irreducible letter (a factorization in disjoint variables would split its bases), so unique factorization and the degrees $r\neq r-1$ give $\cB/S'=(\cB|S)^\ast$. $\square$

The converse, that a piece acting as one mark on every tinker of its room has the law of one mark, is the retention theorem (Thm.~\hb{27.18}{235}) against a piece on one mark with the seam vector $\mathbf x_0$, whose seam function is entire.

So the width-one condition belongs to the factors, and only to them. A bunch of width two or more never splits off; it stays in the core and enters through its law. With both letters active, the dual-pair condition is not a convenience: with connected hinge minors it is exactly the condition for a piece to split off (computed on the test lists: of $263$ pieces of width one, $20$ are read alike, exactly the $20$ dual pairs). In a prime table (Thm.~\hb{20.9}{166}) every set with two marks on each side has width at least two, so the factorization along the hinge tree ends at the prime nodes, by this claim. Their only cells are the complements of single marks; at a full lean these split off and trade the node for a Beta corner and the two-table family of its deletion and contraction on one fewer mark, $((\cB\setminus e;\,s-m_e,0),(\cB/e;\,m_e,0))$. On four marks the pair $\{3,4\}$ of $\binom41$ has $\gamma^K=\alpha_3+\alpha_4$ and $\gamma^\Psi=\alpha_3\alpha_4/(\alpha_3+\alpha_4)$, whose ratio $u(1-u)$ is not constant: it does not split off, and $7\,\Tx3{\IZ}0$ is read across it as a glue (Thm.~\hb{27.9}{230}).

The general object is the core read at the seam points of its pieces, at any width (Thm.~\hb{28.1}{239}). The scale--shape map at the pieces is birational, so Lemma~8.1 of the note makes it an identity of motivic periods. A pinned piece is one over which the stand-in core is constant after the change of scale by its hinge value, and every left factor at every node has the node form of Claim~\ref{cl:hinge14}.

\subsection{The note made self-contained; four small lists}\status{derived; the room-station factorization checked at $796$ sets, the cells of $362$ sets of width one as minors, the lattice identities in ball arithmetic}

The note now carries, with proofs, everything it used from the handbook: the levels and Skeleton values, sections and scale--shape, the walls of complementary sets, the stacked Gamma, Gramma, Zamma, Breta of one bunch and Trailmoment on a line, the period as a stacked Gamma times a Trailmoment of the joint law of the letters (Prop.~\hb{8.11}{56}), finite parts as constant terms, the scale integral as a stacked Gamma times a Breta with its share side (Prop.~\hb{15.1}{87}, Cor.~\hb{15.4}{88}), the continuation by cube integrals, substitution (a matroid by a direct exchange argument; its letters, Thm.~\hb{20.4}{164}; its cells, Prop.~\hb{20.3}{164}), the cell of a set of width one as a minor $M/Y_0\setminus(R'\setminus(Y_0+r_0))$, and pieces with stations, their laws and the glue (Def.~\hb{27.4}{227}, Def.~\hb{27.5}{227}, Thm.~\hb{27.9}{230}). The handbook is cited only for provenance.

One step changed. The converse in Claim~\ref{cl:reach14} used the retention theorem (Thm.~\hb{27.18}{235}) with its condition at the native split. The partners on one mark already suffice, with no condition.

\begin{claim}\label{cl:probe14}
Let a piece $S$ with the stations $\mathcal D$ have $0<\kappa_S<w=\sum_j(s_j+t_j)c_j$. For $\mathbf p$ with positive coordinates, one per active letter and station, let $b_{\mathbf p}$ be the partner on one mark $b$ with the room letters $p_{jD}\alpha_b^{c_j-i(D)}$ and $p^\ast_{jD}\alpha_b^{i(D)}$. If $P_{S\sqcup b_{\mathbf p}}=c\,I_S(\mathbf x_0\circ\mathbf p)$ for every $\mathbf p$, with
\[
I_S(\mathbf y)=\int_0^\infty v^{\kappa_S-1}\prod_j\Big(\sum_Dy_{jD}v^{i(D)}\Big)^{-s_j}\Big(\sum_Dy^\ast_{jD}v^{c_j-i(D)}\Big)^{-t_j}dv
\]
the period of two one-mark pieces, then the law of $S$ is $c\,\delta_{\mathbf x_0}$.
\end{claim}

\emph{Proof.} By the glue, $P_{S\sqcup b_{\mathbf p}}=\int I_S(\mathbf x\circ\mathbf p)\,d\mu_S(\mathbf x)$. On shares $(\mathbf a,\mathbf b)$ with positive coordinates on the pins of the seam, $p^{(\mathbf a,\mathbf b)}I_S(\mathbf x\circ\mathbf p)$ is invariant under the letter scalings and the scale, and its integral over a complement of these directions is $c_W\prod_j\Gst{|\mathcal D_j|}1{\mathbf a_j}{s_j}\Gst{|\mathcal D_j|}1{\mathbf b_j}{t_j}\,\mathbf x^{-(\mathbf a,\mathbf b)}$: the $v$-integral unfolds the scale, and what is left is a Dirichlet integral per letter. So the hypothesis gives $\int\mathbf x^{-(\mathbf a,\mathbf b)}d\mu_S=c\,\mathbf x_0^{-(\mathbf a,\mathbf b)}$ on an open set of the pins. That is a moment generating function on the directions of the pins, finite near a point, so the image of $\mu_S$ in the dual of those directions is $c$ times a point mass; and $\mathbf x\mapsto\log\mathbf x$ modulo the pin normals is injective on the seam torus, the relations at the excesses $0$ and $c_j$ fixing the letter scalings and the product relation the scale. $\square$

So Theorem~D of the note reads: (i) the law is a point mass $\Leftrightarrow$ (ii) pinned $\Leftrightarrow$ (iii) $S$ acts as one mark on every partner $\Leftrightarrow$ (iv) on every one-mark partner.

\emph{Four lists at their centres} (derived; $A_r=\Tx{2}{\Ilat{r+3\mathbb Z}}{0}$).
\begin{itemize}
\item $\binom32$: in the chart $\alpha_3=1$, with $\alpha_{\{1,2\}}=\rho(u,1-u)$ and $Q=u(1-u)$, the letters are $1+\rho$ and $\rho(1+\rho Q)$, so
\begin{align*}
P&=\int_0^1\frac{-\log Q}{1-Q}\,du=-2\int_0^1\frac{\log u}{1-u+u^2}\,du\\
&=2\,\Tx{2}{\big(\Ilat{1+6\mathbb Z}+\Ilat{2+6\mathbb Z}-\Ilat{4+6\mathbb Z}-\Ilat{5+6\mathbb Z}\big)}{0}=3\,(A_1-A_2),
\end{align*}
the last by the dilations $2+6\mathbb Z=2(1+3\mathbb Z)$, $4+6\mathbb Z=2(2+3\mathbb Z)$ (alias $3L(\chi_{-3},2)$).
\item $\binom42$: the complement $R$ of a mark $e$ is a dual-pair cell, $\binom42|R=\binom32$ and $\binom42/e=\binom31=\binom32^\ast$, with $\kappa_R=1$; the core is $\binom{\{e,h\}}1$ at the masses $(1,1)$, whose motivic period is the Beta corner $\Gst21{1,\ 1}2=1$. So $P^{\mathfrak m}_{\binom42}=P^{\mathfrak m}_{\binom32}$ and $P=3(A_1-A_2)$. Its pairs have width two and none splits off.
\item $\binom41$: at the pair $S=\{3,4\}$ the seam vector is $(1,q^{-1},q,1)$ with $q=(u(1-u))^{1/2}$, the scale form $q^{-1}du$, and the room the same with the excesses reversed. Both absorbed,
\begin{gather*}
P=\int_0^1\!\!\int_0^1\frac{\log(Q'/Q)}{Q'-Q}\,du\,du'=\int_0^\infty F(t)^2\,dt,\\
F(t)=\int_0^1\frac{du}{t+u(1-u)}=\frac2r\log\frac{r+1}{r-1},\qquad r=(1+4t)^{1/2},
\end{gather*}
and then $P=4\int_0^1\log^2z\,dz/(1-z^2)=8\,\Tx3{\Ilat{1+2\mathbb Z}}0=7\,\Tx3{\IZ}0$: the glue of the laws of two pairs, which does not close to one evaluation.
\item The eight marks of \S14.5: $P^{\mathfrak m}=\big(P^{\mathfrak m}_{\binom32}\big)^3$ and $P=27(A_1-A_2)^3$; the bunch of width two never splits off, but the core splits further at the complement of a mark.
\end{itemize}

\subsection{What the proof uses of the Foreman tree}\status{derived}

\begin{itemize}
\item Claims~\ref{cl:weights}, \ref{cl:serif14}, \ref{cl:children14} and~\ref{cl:hinge14} use only that the tree is a tree of bunches, at any width, as do (i), (ii), (iv), (v) of Thm.~\hb{19.57}{159}.
\item The separation of the coefficients of the Bretas of Claim~\ref{cl:hinge14}, one Skeleton value per letter, the room's hinge value over the bunch's, uses width one at the node.
\item Claim~\ref{cl:pinned14} places no condition on the width of a bunch: width one is what a cell is (Claim~\ref{cl:onehinge14}), and the bunch sits in the core at whatever width it has.
\end{itemize}

\emph{A Foreman bunch of width two} (derived; the numbers computed). Take the list of rank three on six marks whose bases are all triples except $\{1,2,3\}$. The permutations of each triple are automorphisms, so marks of one triple have one reading, and
\[
\Pr(e\in T)=\tfrac9{19}\ \ (e\in\{1,2,3\}),\qquad \Pr(e\in T)=\tfrac{10}{19}\ \ (e\in\{4,5,6\}).
\]
At a lean $\delta$ the odd view of a mark with no partners has first derivative $\delta\big(2\Pr(e\in T)-1\big)$, which is $-\delta/19$ on the first triple and $\delta/19$ on the second. So the classes at a lean are the two triples (Def.~\hb{10.1}{64}), and the Foreman bunch there is a triple. Each triple has width $2+3-3=2$. Along $\{1,2,3\}$ the levels are single products,
\[
K(\rho u,\alpha_R)=\alpha_4\alpha_5\alpha_6+\rho\,e_1(u)\,e_2(\alpha_R)+\rho^2e_2(u)\,e_1(\alpha_R),
\]
yet the discriminant is $45$ at the gauge, so the Skeleton values are not rational. At the centre the exchange of the two triples carries $K$ to $\Psi$, and the list is an atom. So the three Foreman trees of this list (Thm.~\hb{14.6}{82}) are one atom and two trees whose root bunch has width two. At those roots Claims~\ref{cl:weights}--\ref{cl:hinge14} hold as stated, with two Skeleton values per letter. The list is prime: its sets with two marks on each side have widths two and three (computed). So its only cells are the complements of single marks, and Claim~\ref{cl:pinned14} cannot split it; at the centre it does not pin even those, the two hinge values of every such cell differing (computed).

The claims hold along every tree of bunches, and along the Foreman tree what they attach to each vertex is fixed by the family. What the tree provides is proved in the reference:
\begin{itemize}
\item it makes no choice: with the tables fixed, one tree per ray of the strengths and ray of the lean, so three for one list, possibly coinciding (Thm.~\hb{14.6}{82}; for $\binom{[4]}1$ all three are one atom);
\item it follows the symmetries: a bijection mapping bases onto bases maps each tree onto itself, and one mapping bases onto cobases maps the tree at a lean onto the tree at the opposite lean (Thm.~\hb{14.8}{82});
\item its native split hands a centre down as centres (Prop.~\hb{12.24}{75}), so that along the tree of a centre both letters stay active at every vertex;
\item on a molecule with a connected active table it is a chain whose bunches never shrink, the leaf the largest (Cor.~\hb{14.23}{85});
\item it suffices: the period is the word of the Foreman tree, as it is the word of every tree of bunches (Cor.~\hb{19.9}{137}, Thm.~\hb{19.57}{159});
\item the tree of the centre is the tree of the pseudointegral: the procedure draws it once from the tables and follows it at every lean, in the Foreman's order (Def.~\hb{16.12}{97}, Def.~\hb{16.26}{101}), and computes the pseudointegral with no finite part, chart, order or closing rule (Thm.~\hb{16.13}{98}); the pseudointegral equals the integral wherever that converges and every regular iterated finite part off the pole hyperplanes (Cor.~\hb{16.14}{98}, Thm.~\hb{16.27}{102}, Cor.~\hb{16.33}{105}), while an iterated integral in another order computes a quasi-integral, which on a pole hyperplane can depend on the order (Thm.~\hb{16.21}{100}).
\end{itemize}
So the program's three words about the Foreman bunch (deterministic; working wherever some bunch does; a pseudo- among generalized-) are statements of the reference, the third for the tree of the centre. (Revised in \S14.8: the first wording here undersold the tree, and the second called these a contention.)

\emph{Checks} (\texttt{note\_lemmas\_check.py}):
\begin{itemize}
\item face factorization on $1592$ faces;
\item the children lemma on $8198$ minors;
\item the width bound on $305564$ block tables;
\item the levels at every width on $1340$ letters ($680$ of width one, $560$ of width two, $100$ of width three), each level against the sum over its patterns of the children, the extreme levels against the products;
\item at width one with the $\binom J1$ adjoined, $14432$ letters;
\item at width two, a non-square discriminant for all $560$ letters;
\item the dual-pair identities on four dual pairs, exactly, and the multi-hinge identities of the eight-mark list;
\item Claim~\ref{cl:onehinge14} through its criterion on the levels, on the same $1340$ letters at all their sets: it holds exactly at width at most one; of the $660$ letters of width two or three, $220$ fail at a level with two children and $440$ only at $f_{d+1}^2$ against $f_df_{d+2}$;
\item Claim~\ref{cl:pinned14} on the eight-mark list at $(1,1)$ and at $(2,1)$, and on $\binom42\langle b_1,b_2,\binom{R_1+g_1}2,\binom{R_2+g_2}3\rangle$ at the full lean $(4,0)$, whose cells are not dual pairs: the walls and balances of every family, the identity of the degree-zero parts of the forms and the cell families against the Beta form, exactly, and the logarithmic Jacobian $\pm1$ at ten random rational points;
\item Claim~\ref{cl:reach14} for one list at the centre on the test lists: no piece of width two ($344$) or three ($90$) has the law of one mark; of the $263$ pieces of width one, the $20$ read alike are exactly the $20$ dual pairs; proportional hinge values always equal;
\item the six-mark list above: the two probabilities, the widths, the discriminant $45$, primality; at the centre no cell $\cE\setminus\{e\}$ is read through one hinge value, and at $(2,0)$ with unit masses each is pinned, the form the Beta form times that of the deletion and the contraction, exactly.
\end{itemize}
The note had one independent referee pass, which found no remaining gap after the corrections recorded here.

Alias, once:
\begin{itemize}
\item Claim~\ref{cl:weights} is Brown's small graphs principle in its affine form (Thm.~8.11 of his \emph{Feynman amplitudes, coaction principle, and cosmic Galois group}), here for families of lists.
\item The removed $\binom J1$ are his affine model.
\item Whether they can be dropped is his Remark~8.4.
\end{itemize}


\subsection{The note made native}\status{derived; the new lemmas checked: duality on the $21$ test lists, nested windows along $5846$ maximal chains, children against widths in $30336$ inequalities, the exchange of the six-mark list; four referee passes, three against the handbook}

Rule of the pass: if a step can be done from the reference, it is. The note now argues from the exchange property and the serif family alone; from outside it keeps only Brown's motivic periods, coaction and affine model, with resolution of singularities and the functoriality of relative cohomology on the motivic side. Classical names (matroid, Kirchhoff, Symanzik, connectivity, $U_{k,n}$, Stieltjes, Mellin, $\zeta$, $L$, Feynman) appear only as parenthetical aliases.

\emph{The foundation, rebuilt from the exchange property} (derived):
\begin{itemize}
\item \emph{Duality.} For bases $T,T'$ and $x\in T'\setminus T$ some $y\in T\setminus T'$ has $T-y+x$ a basis, by induction on $|T\setminus T'|$: exchange a different mark of $T'\setminus T$ first. So the cobases form a list.
\item \emph{The window} (Thm.~\hb{7.11}{49}). The traces of a set fill $[d_N(S),r_N(S)]$ along an exchange path. Along a chain $S_1\subseteq\cdots\subseteq S_\ell$ some basis is at the top of every $S_i$: exchanging out a mark outside $S_\ell$ can only bring in a mark outside $S_{\ell-1}$, since the traces below are at their tops. This replaces the greedy algorithm, and it is Lemma~\hb{8.5}{55}.
\item \emph{Extreme levels are products} (Thm.~\hb{12.17}{73}). If two bases are at the top of $B$, their room sides can be swapped: exchanging out a room mark cannot bring in a mark of $B$. Restriction and contraction are defined this way, as top traces and the room sides of top bases, and $(N|B)^\ast=N^\ast/R$.
\item \emph{Children} (Lemma~\hb{12.2}{70}). Children are lists, $(N_{(p)})^\ast=(N^\ast)_{(B\setminus p)}$, children of children are children, and children do not widen: the traces of $A$ in $N_{(p)}$ are those of $A\cup W$ in $N$ shifted by $|p\cap W|$. This replaces submodularity.
\item \emph{Bottom levels} replace the initial forms. The bottom level of a chain is the product of the block lists $N_{(T_0\setminus Z_i)}$, children at the traces of a basis at the bottom of every $S_j$; one block can be exchanged at a time by the extreme-level swap at $S_{i-1}$ and at $S_i$. Along a maximal chain it is one monomial with coefficient $1$, which is all the charts need.
\end{itemize}

\emph{Single steps replaced} (derived):
\begin{itemize}
\item \emph{Substitution.} The augmentation step is the exchange property at the hinge: from $Y_1+h$ toward $Y_2\in\mathcal H\setminus h$ at $h$ itself.
\item \emph{The cell of a set of width one.} An adjacent pattern $Y_0+b$ of a bottom pattern $Y_0$ of $R'$ is reached by exchanges, and the cell is the child of $M$ at $Y_0$ on $R'-b$ (Prop.~\hb{12.5}{70}). Checked at all $362$ sets of width one, at every bottom pattern.
\item \emph{Gauss's lemma.} Proportional hinge values are equal because the leading term of a product of $0/1$ polynomials has coefficient $1$.
\item \emph{Mellin inversion} is Splitting a coefficient (Cor.~\hb{4.4}{26}), proved by Closing (Thm.~\hb{4.5}{26}). Shift the line left across the poles of $\Gamma(a)$; the residues are the expansion lattice of $\sG_{w;(1,1)}$ (Prop.~\hb{2.10}{15}). The shifted line dies because consecutive terms of $(|w|)_Jz^J/\Gamma(J+\frac12)$ have ratio tending to $z<1$, and the bounds come from the shift relation and the shape integral. No Stirling.
\item \emph{Moment generating functions.} The station transform (Def.~\hb{27.20}{236}) gives the seam function on an open set of the pins. Then, with $\nu$ the law reweighted at one point of the pins,
\[
\int\Big((\mathbf x/\mathbf x_0)^{\boldsymbol\delta/2}-(\mathbf x/\mathbf x_0)^{-\boldsymbol\delta/2}\Big)^2d\nu=C-2C+C=0
\]
forces a point mass.
\item \emph{The glue} is a Trailmoment of the two laws, its kernel a Zamma (Thm.~\hb{27.9}{230}):
\[
P_{S\sqcup Z}=\Gamma(\kappa_S)\big(\sT_{(\kappa_S,\mathbf s,\mathbf t);t_{\mathcal D}}[\mathbf x\circ\mathbf x']\,(\mu_S\otimes\mu_Z)\big)(\mathbb T\times\mathbb T').
\]
\end{itemize}

\emph{Line one.} The introduction opens with the centre and both letters. Cremona,
\[
P_{((N_j;s_j,t_j))_j}(\mathbf m)=P_{((N_j;t_j,s_j))_j}(2\sigma\one-\mathbf m),
\]
swaps the full leans and fixes the centre. $\binom{[4]}1$ has rational periods at both full leans: at $(s,0)$ the period is the shape integral $\Gst41{\mathbf m}{m_{\cE}}$. At its centre it has $7\,\Tx3{\IZ}0$, and no pair splits off.

\emph{The Foreman tree} is presented as above (\S14.7). Its peel is read with its record, children indexed by table and pattern. Merging equal children can cancel a letter when exponents of both signs merge. A referee found such a family: $((\binom{[5]}2;1,0),(\binom{[5]}3;-1,0),(\binom{[5]}4;4,0))$ at the masses $3\cdot\one$, peeled at $\{1\}$, loses $e_2$ on the room. With the record, the native split hands a full lean down full table by table.

\emph{Corrections from the referee passes} (derived):
\begin{itemize}
\item ``A set of width two stays inside a factor'' was false. In the eight-mark list, $\{b_1,x_1\}$ has width two and Theorem~C cuts it. The true statements are: it never splits off, and a bunch $B\subseteq L$ stays in the core with its width.
\item The width-one separation of Theorem~B(iv) is now stated on contours of shares where the two simplex integrals converge absolutely; the third pass below fixes where it holds at the node form.
\item The Foreman bunch at a lean is chosen first by the class of the centre containing it (Def.~\hb{13.6}{78}).
\item $\Ilat{b+q\mathbb Z}$ is the lattice input on the coset, and $\IZ$ the counting input.
\end{itemize}

\emph{Third and fourth referee passes} (derived; the fourth checked every handbook number cited in the note):
\begin{itemize}
\item \emph{Theorem~B(iv).} The contour conditions name the two integrands,
\[
\int_{\Delta_{Z_B}}|A_t|\prod_g(\gamma^g_B)^{\re b_g}\,d\omega<\infty,\qquad \int_{\Delta_{Z_R}}|B_t|\prod_g(\gamma^g_R)^{-\re b_g}\,d\omega<\infty,
\]
and the proof bounds the share kernels for Fubini, by Splitting a coefficient and the shape integral. There is one share for each Skeleton value of the term but one. If $0<a_t<|\mathbf w_t|$ the separation holds at $\kappa=a_t$; otherwise it holds on the strip, with $\mathrm{Sc}_t$ pointwise the constant term of the continuation. A constant term does not pass through the shape integrals: $\int_0^1u^{\kappa-1}du=1/\kappa$ has constant term $0$ at $\kappa=0$, while $u^{-1}$ is not integrable.
\item \emph{The abstract and the introduction} now carry the hypotheses of the theorems: convergence and nonnegative exponents for splitting off, partners glued balanced, a pinned set for the motivic split, one list for the full-lean product, at most one Skeleton value, constant terms of stacked Gammas times Bretas, every left factor only at the root, and the Gramma for the interest at positive exponents with $\kappa_E>0$. ``Splits off'' is defined once, as Theorem~D(iii).
\item \emph{The Foreman tree.} The order of classes is ``the order'' (Def.~\hb{13.5}{78}, Thm.~\hb{13.7}{79}); ``the Foreman's order'' is the integration order (Def.~\hb{16.12}{97}), along the tree of the centre at every lean (Def.~\hb{16.26}{101}). The classes are the level sets of the reading (Def.~\hb{10.1}{64}), and for one list the orbits of the bijections mapping bases onto bases or onto cobases (Cor.~\hb{11.17}{69}). The tree is defined for nonnegative exponents (Def.~\hb{9.1}{61}). The word is integrated over shares of the exponents among patterns, its contours through the native splits at the native gauge of the point, which is the table's gauge only at the centre with masses $\sigma\one$.
\item Corollary~5.1 gives constant terms of continued periods; its title no longer calls them values of the serif family. A block family whose walls are positive at its masses has nonnegative exponents and positive masses, so its period is one.
\end{itemize}

\subsection{The note substantiated by the handbook}\status{derived; the new evaluations checked in ball arithmetic; every handbook citation of the note checked item by item against the reading copy; two referee passes}

Rule of the pass: what is not born of the handbook distracts from the note's self-containedness and suggests that the import is needed. So every step the handbook takes natively is taken its way, and every definition and statement on the side of the periods names the item it specializes, by printed number and page. The proofs stay complete in the note; the remarks and examples that go further quote the handbook where they do.

\emph{The concordance} (derived; $106$ citations checked against the reading copy):
{\sloppy
\begin{itemize}
\item Lists, windows, levels, children: Def.~\hb{7.1}{48}, Prop.~\hb{7.2}{48}, Def.~\hb{7.3}{48}, Thm.~\hb{7.11}{49}, Lemma~\hb{8.5}{55}, Def.~\hb{12.1}{70}, Lemma~\hb{12.2}{70}, Thm.~\hb{12.17}{73}, Prop.~\hb{14.5}{82}; Skeleton and hinge values: Def.~\hb{12.9}{71}, Def.~\hb{12.13}{73}, Def.~\hb{18.4}{115}, Cor.~\hb{18.7}{117}.
\item Families, walls, Cremona: Def.~\hb{8.1}{53}, Thm.~\hb{8.10}{56}, Def.~\hb{9.1}{61}, Thms.~\hb{8.12}{57}, \hb{8.13}{57}, Prop.~\hb{8.15}{57}; charts and continuation: Thm.~\hb{8.7}{55}, Lemma~\hb{8.6}{55}.
\item The serif family: Defs.~\hb{1.2}{11}, \hb{2.1}{13}, \hb{2.12}{15}, \hb{3.1}{18}, \hb{5.1}{28}, \hb{5.2}{28}; Thm.~\hb{3.9}{21}; Prop.~\hb{8.11}{56}; finite parts: Def.~\hb{1.4}{11}, Thm.~\hb{1.6}{12}; the share side: Cor.~\hb{4.4}{26}, Thm.~\hb{4.5}{26}, Prop.~\hb{15.1}{87}, Cor.~\hb{15.4}{88}; the Laplace side: Thm.~\hb{5.11}{31}, now Lemma~2.23 of the note.
\item Cells and substitution: Def.~\hb{20.1}{163}, Prop.~\hb{20.3}{164}, Thms.~\hb{20.4}{164}, \hb{20.7}{165}, Prop.~\hb{12.5}{70}, Thm.~\hb{18.6}{116}; pieces, laws, the glue: Defs.~\hb{27.1}{224}, \hb{27.4}{227}, \hb{27.5}{227}, \hb{27.20}{236}, Props.~\hb{27.2}{225}, \hb{27.10}{231}, Lemma~\hb{27.8}{229}, Thms.~\hb{27.9}{230}, \hb{27.18}{235}.
\item The theorems at the level of periods: B is Thm.~\hb{15.2}{87} with Cor.~\hb{15.7}{89}, and B(iv) is Thm.~\hb{18.13}{119}; C is Thm.~\hb{21.5}{173} and Cor.~\hb{21.6}{173}, for one cell Thm.~\hb{18.10}{118} and Prop.~\hb{18.12}{119}; D is Thm.~\hb{27.18}{235}, at width one Thm.~\hb{22.7}{185} and Cor.~\hb{22.8}{186}; Corollary~5.5 is Thm.~\hb{20.9}{166}, Cor.~\hb{21.7}{174}, Thm.~\hb{25.6}{208}, Prop.~\hb{25.10}{209}.
\item What the note adds: the motive (Theorems A, B(v), C(ii), Corollary~5.5); Theorem~B for every form on every node piece; one hinge value of any degree forces width at most one (the converse of Thm.~\hb{18.6}{116} reads it off two terms); so a point-mass law forces width at most one; and two connected hinge minors with one even letter are a dual pair (with Thms.~\hb{7.9}{49}, \hb{10.6}{65}).
\end{itemize}
}

\emph{$\binom{[4]}1$ at its centre, natively} (derived; the identity, the residues and the value checked at $256$ bits). Theorem~B(iv) at the root bunch $\{3,4\}$ is the one-table display of Thm.~\hb{18.13}{119}, its free share the even share:
\[
P_{\binom{[4]}1}=\frac1{2\pi i}\int_{(1/2)}\Gst40{a,\ a,\ 1-a,\ 1-a}{}\ \Gst21{1-a,\ 1-a}{2-2a}\ \Gst21{a,\ a}{2a}\,da .
\]
Reflection and shift (Prop.~\hb{1.3}{11}) give
\[
\Gst21{a,\ a}{2a}\Gst21{1-a,\ 1-a}{2-2a}=\frac{2\pi\cot\pi a}{1-2a}=-\frac{\partial_a\log\Gst40{a,\ a,\ 1-a,\ 1-a}{}}{1-2a},
\]
so the centre is the critical-line word $\operatorname{crit}_1[1]$ of Def.~\hb{33.3}{298}: \S9.7 at $k=1$, $R_1=1$. Closing on the triple poles at the positive integers, residues $8(1-2m)^{-3}$, gives (Thm.~\hb{33.6}{298}, Thm.~\hb{5.12}{32}(i))
\[
P_{\binom{[4]}1}=\sum_{m\ge1}\big(m-\tfrac12\big)^{-3}=\Tx{3}{\Ilat{\frac12+\mathbb Z}}{0}=7\,\Tx{3}{\IZ}{0}.
\]
This replaces the roots of $u^2-u-t$, two substitutions, a series from $\int_0^1z^{a-1}\log^2z\,dz$, and the Stieltjes alias. At the full leans one letter is active, with one Skeleton value, and no share is free (Cor.~\hb{18.14}{121}); at the centre the share runs over a line.

\emph{$\binom{[3]}2$ at its centre, natively} (derived; checked). At the bunch $\{1,2\}$, room one mark, the share integral $\frac1{2\pi i}\int_{(1/2)}\Gst40{a,\ a,\ 1-a,\ 1-a}{}Q^{-a}da$, $Q=u(1-u)$, closes to the left to $\log Q/(Q-1)$ (Lemma~\hb{5.16}{34}(ii)). The reflection $u\mapsto1-u$, the coordinate $y=\log(1/u)$ and the Laplace side (Thm.~\hb{5.11}{31}) give
\[
P_{\binom{[3]}2}=2\,\Tx{2}{\big(\Ilat{1+6\mathbb Z}+\Ilat{2+6\mathbb Z}-\Ilat{4+6\mathbb Z}-\Ilat{5+6\mathbb Z}\big)}{0}=3\,\Tx{2}{\big(\Ilat{1+3\mathbb Z}-\Ilat{2+3\mathbb Z}\big)}{0},
\]
the Laplace transform of the period-six input being $(e^{-y}+e^{-2y})/(1+e^{-3y})$, a sum of geometric series. This replaces the expansion of $1/(1-u+u^2)$ and $\int_0^1u^{a-1}\log u\,du$. The share-side closing on three marks stays open (Not yet computed, (5)). And $\binom{[4]}2$ splits off a mark by one Beta corner at the level of periods by Cor.~\hb{22.8}{186}(iii), its bases being closed under complement.

\emph{Proofs} (derived). Proportional hinge values are equal: along a maximal chain the bottom level of each letter is one monomial with coefficient $1$, and the bottom level of a product of polynomials with nonnegative coefficients is the product of theirs. This replaces the monomial order of \S14.8. In the one-list criterion the factors of a homogeneous basis polynomial are homogeneous, so a factorization would make the variables of one factor a separator; irreducibility is Thm.~\hb{7.9}{49}, and the factor-by-factor reading Thm.~\hb{10.6}{65}.

\emph{Corrections from the two referee passes} (derived; one counterexample computed):
\begin{itemize}
\item The period separates on the shares at every width only as the handbook states it. One share per station needs the bunch word and the room word to have Newton sums of full dimension (Thm.~\hb{24.12}{204}, Lemma~\hb{28.8}{244}); otherwise every direction a side does not see costs one more pin, with reduced periods (Thm.~\hb{27.11}{232}). Counterexample to the unconditional reading (computed): the spanning trees of a $4$-cycle with one chord, at the centre, $S$ three edges of width two. The two room sides of excess one have different room children, so every station of room sides is one room side, the room word is a monomial, and its seam function is $+\infty$ at every share.
\item Width two: ``the discriminant is not a square'' is now ``not a constant times a square'', tested on the primitive part; the value $12$ at $\alpha=\one$ excluded only rational squares. On the line $\alpha=(1,1,1,\lambda)$ the discriminant of $\binom{[4]}2$ along a pair is $4(\lambda^2+\lambda+1)$, with simple roots. All $560$ letters of width two pass the stronger test (computed).
\item Three citations corrected: the condition of Thm.~\hb{27.22}{237} when both letters are active; Thm.~\hb{22.7}{185} with Def.~\hb{22.6}{185} for the width-one converse in Theorem~D; Prop.~\hb{18.12}{119} for one cell in Theorem~C(ii). The coordinate $y=\log(1/u)$ is no longer attributed to Def.~\hb{1.1}{10}(ii).
\end{itemize}

\subsection{One letter against two on $\binom{\cE}1$}\status{derived; the chart against Cremona checked on $495$ integer data, the values in ball arithmetic; the Bessel values and the wheel recalled}

Raised by the pitch: on $\binom{[4]}1$ every convergent period of one letter is rational, at the data of a family (integer exponents and masses, Def.~2.11 of the note). The reason is Corollary~5.5(b) of the note: the list has one parallel node and no prime node, and one letter alone reaches only what the prime nodes carry. On a prime list one letter already gives an odd zeta value: the wheel with three spokes has a prime list, and its Feynman period, the full lean $(0,2)$ at unit masses, is $6\,\Tx3\IZ0$ (recalled).

\emph{The full leans} (derived; checked). $K=e_1=|\alpha|$ and $\Psi=e_3=\alpha^EK(\alpha^{-1})$, so the shape integral (Thm.~\hb{3.9}{21}) and Cremona (Thm.~\hb{8.13}{57}) give
\begin{gather*}
P_{(\binom{[4]}1;\,s,0)}(\mvec)=\Gst41{m_1,\ m_2,\ m_3,\ m_4}{m_E},\\
P_{(\binom{[4]}1;\,0,t)}(\mvec)=\Gst41{t-m_1,\ t-m_2,\ t-m_3,\ t-m_4}{t},
\end{gather*}
with $s=m_E$, resp.\ $m_E=3t$. They converge exactly when every $m_e\ge1$, resp.\ every $m_e\le t-1$, which needs $t\ge4$; the smallest are $\int\Omega/K^4=\int(\alpha_1\alpha_2\alpha_3\alpha_4)^2\,\Omega/\Psi^4=\frac16$. In the frame of Example~6.3 of the note, one letter makes the $\rho$-integral a monomial in the two ratios of hinge values times a Beta corner,
\[
\int_0^\infty\rho^{A}\,(\varrho_S\rho+\varrho_{S'})^{-t}\,d\rho=\varrho_S^{-A-1}\varrho_{S'}^{A+1-t}\,\Gst21{A+1,\ t-A-1}{t},
\]
and the shape integrals over $\Delta_S$ and $\Delta_{S'}$ are then Beta corners. Both letters make it $(\log\varrho_{S'}-\log\varrho_S)/(\varrho_{S'}-\varrho_S)$, a function of the ratio, and one share stays free. Off the data of a family a full lean is a stacked Gamma at fractional arguments: $\Psi$ alone at unit masses balances only at $t=\frac43$, where
\[
P=\Gst41{\frac13,\ \frac13,\ \frac13,\ \frac13}{\frac43}=3\,\Gamma(\tfrac13)^3 .
\]

\emph{The centre by a second route} (derived). The identity $\frac{\log b-\log a}{b-a}=\int_0^\infty\frac{d\tau}{(\tau+a)(\tau+b)}$ separates the bunch from the room:
\[
P_{\binom{[4]}1}=\int_0^\infty\Big(\int_0^1\frac{du}{\tau+u(1-u)}\Big)^2d\tau
=2\int_0^1\log^2\frac{1+w}{1-w}\,\frac{dw}{w}
=4\int_0^\infty y^2\,\frac{e^{-y}}{1-e^{-2y}}\,dy ,
\]
by $\tau+\frac14=\frac1{4w^2}$ and then $\frac{1+w}{1-w}=e^{y}$. By the Laplace side (Thm.~\hb{5.11}{31}) the last is $8\,\Tx{3}{\Ilat{1+2\mathbb Z}}{0}$, and the dilation by $2$ (Thm.~\hb{5.9}{31}) and the lattice relation (Thm.~\hb{5.12}{32}(i)) give $\Tx{3}{\Ilat{\frac12+\mathbb Z}}{0}=7\,\Tx3\IZ0$, the value of \S14.9.

\emph{Every $n$} (derived; the aliases recalled). The centre of $\binom{\cE}1$ at unit masses is \S6.1 at $\sigma=1$, the Gramma power
\[
P_n=\int_{\sigma_{\cE}}\frac{\Omega_{\cE}}{e_1\,e_{n-1}}=\int_0^\infty\sG_{0;(1,-1)}(1,y)^n\,dy .
\]
Read $y$ first, as in \S6.2: it gives $e_n(u)/e_{n-1}(u)$ against $\prod_eu_e^{-1}e^{-|u|}$, and scale and shape leave $\int_{\Delta_{\cE}}d\omega/e_{n-1}$. (Alias: $\sG_{0;(1,-1)}(1,y)=2K_0(2\sqrt y)$, so $P_n=2^{n-1}\int_0^\infty t\,K_0(t)^n\,dt$, the equal-mass vacuum banana with $n$ propagators in two dimensions, a Bessel moment.) So $P_2=1$; $P_3=3\,\Tx{2}{(\Ilat{1+3\mathbb Z}-\Ilat{2+3\mathbb Z})}{0}$, Example~6.1 of the note read through duality; and $P_4=7\,\Tx3\IZ0$. These agree with the recalled Bessel moments $\frac12$, $\frac34L(\chi_{-3},2)$, $\frac78\zeta(3)$. At $n=5$ the recalled value is the modular $L$-value of Not yet computed~(8). Every full lean of every $\binom{\cE}1$ is rational at integer data, by one parallel node and no prime node. So the whole sequence of centres is made by the second letter.

\emph{For the pitch.} The sentence becomes ``every convergent period of $K$ alone or of $\Psi$ alone, with integer exponents, is rational''. The centre is better named as the four-propagator banana before the reader names it.

\subsection{The centre of $\binom{\cE}1$ on the share side: the room is a shorter centre ray}\status{derived; the five-mark closing not yet computed; the content graded by tier proposed}

Take the bunch $S$ a pair and the room $S'$ the other $n-2$ marks. On $\Delta_S\times\Delta_{S'}$, $K(\rho u,u')=\rho+1$ and
\[
\Psi(\rho u,u')=\rho\,\big(\varrho_S(u)\,e_{n-3}(u')\,\rho+e_{n-2}(u')\big),
\]
since a set of $n-1$ marks misses one mark, of $S'$ or of $S$. The share side of Theorem~B(iv) of the note, as in its Example~6.3, gives
\[
P_n=\frac1{2\pi i}\int_{(1/2)}\Gst40{a,\ a,\ 1-a,\ 1-a}{}\ \Gst21{1-a,\ 1-a}{2-2a}\ R_{n-2}(a)\,da,
\]
with the room's factor
\[
R_{n-2}(a)=\int_{\Delta_{S'}}e_{n-3}^{-a}\,e_{n-2}^{\,a-1}\,d\omega=P_{(\binom{S'}1;\,a,\,a)}(a\one),
\]
the centre ray of the $(n-2)$-mark list read at $\sigma=a$ (balance: $(n-2)a=a+(n-3)a$). By \S6.1 both factors are Mellin transforms of Gramma powers,
\begin{multline*}
\Gst40{a,\ a,\ 1-a,\ 1-a}{}\Gst21{1-a,\ 1-a}{2-2a}\,R_{n-2}(a)\\
=\int_0^\infty y^{-a}\,\sG_{0;(1,-1)}(1,y)^2\,dy\cdot\int_0^\infty y^{a-1}\,\sG_{0;(1,-1)}(1,y)^{n-2}\,dy,
\end{multline*}
so the share pairs $\sG^2$ against $\sG^{n-2}$, and the display returns $\int_0^\infty\sG_{0;(1,-1)}(1,y)^n\,dy$, \S6.1 at $\sigma=1$ (a check). The collapses:
\begin{itemize}
\item $n=3$: the room is one mark and $R_1=1$; the closing to the left gives the frequency $\frac13$ (\S8.5; Example~6.1 of the note, by duality).
\item $n=4$: $R_2(a)=\Gst21{a,\ a}{2a}$, of tier one; the integrand is a stacked Gamma, the critical-line word, and the closing gives the coset $\frac12$ (\S8.6; Example~6.3 of the note).
\item $n=5$: $R_3$ is the three-mark centre ray in the share, of tier two (\S9.5): at integer $\sigma=k$ its values are $A_kP_3-B_k$ with rational $A_k,B_k$, and they satisfy the recurrence of symbol $9\lambda^2-10\lambda+1$ (\S8.7, \S13.8). This is the first $n$ at which the share integrand is not a stacked Gamma, and the closing must read the Laurent data of a function of tier two at its poles.
\end{itemize}
The five-mark value is in hand: it is the display above at $n=5$, equivalently $\int_0^\infty\sG_{0;(1,-1)}(1,y)^5\,dy$. It converges, it is certified to any precision, and it carries the relations (Thm.~\hb{8.3}{54}), the Cremona symmetry (Thm.~\hb{8.13}{57}) and the ray recurrence (Thm.~\hb{35.19}{313}) of the family. What the closing adds is its reading inside the family (not yet computed):
\begin{itemize}
\item \emph{Its input.} The closing against a room of tier two gives a Trailmoment of a live input of tier two. Its row, clearing and frame along the ray (\S8.7) are the value's computational properties.
\item \emph{Its floor} (proposed). Grade the content by tier: tier one is the lattice inputs of $n=3,4$, and tier two adds the Trailmoments of live inputs of tier two. \S9.5 then reads: the tier of the room, which the tree reads before any value, says on which floor of the content a value lives.
\item \emph{Its primes.} The critical values $25,9,1$ carry $3$ and $5$ (\S13.6). Whether they enter the live input's clearing, or stay in the frame as $2$ did at four marks (\S13.8).
\item \S13.5: the first value on which the cutting rule at tier two acts.
\end{itemize}
The alias at the end, and only there: the value is recalled as a modular $L$-value of level $15$ (the five-Bessel moment). It names the number for an outside reader; it is not the criterion.

\emph{Correction.} \S14.10 first gave $\sG_{0;(1,-1)}(1,\cdot)$ a second name; it now reads the centre as the Gramma power of \S6.1.

\subsection{What an answer is: members, relations, independence}\status{derived; refined after discussion; the classical readings recalled}

Every period of a family is a member of the serif family with its slots filled. On five marks $\int_\sigma\Omega/(K\Psi)$ is $\int_0^\infty\sG_{0;(1,-1)}(1,y)^5\,dy$, the Gramma power at the slot $n=5$ (\S6.1). Computability does not tell one member from another. What makes a value an answer is that its place among the other members is known, and comparison runs through the family's relations. It has three parts.
\begin{itemize}
\item \emph{Equality} is proved by the relations. On four marks the closing and the lattice relation (Thm.~\hb{5.12}{32}) tie the Gramma power at $n=4$ to the Trailmoment of the counting input,
\[
\int_0^\infty\sG_{0;(1,-1)}(1,y)^4\,dy=7\,\Tx3\IZ0,
\]
and along the whole ray $\langle1\rangle_k=A_k\cdot7\,\Tx3\IZ0-B_k$ with $A_k,B_k\in\mathbb Q$ (\S13.7). Which member of a tied class is used as its representative is a convention; the tie is the content.
\item \emph{Inequality} of two given members is certified by evaluation in ball arithmetic.
\item \emph{Independence}, the absence of any rational relation, is what relations cannot give: a relation proves a tie, and no number of relations proves that there is none. Every vocabulary meets this (recalled: for the values of the Gamma function at rationals, that reflection, shift and multiplication are all the relations is a conjecture; for the zeta values, independence is proved on the motivic side). In the family this is the coaction's part: on the motivic side it separates members that no relation ties.
\end{itemize}
At five marks the task is to find the ties of $\int_0^\infty\sG_{0;(1,-1)}(1,y)^5\,dy$ to the members already in use, by the closing of \S14.11 through a Trailmoment of a live input of tier two, or to certify by the coaction that there are none, in which case it is a new representative. Either outcome is an answer inside the family. An outside name, if there is one, is an alias.

\subsection{The five-mark centre is a trail: Corollary 19.40 worked}\status{derived from the handbook; the values computed in ball arithmetic, no quadrature}

The handbook already answers the question of \S14.11--\S14.12. By Cor.~\hb{17.20}{113} the centre of the circuit on five marks at unit masses is $\mathcal M_5(1)$, where $\mathcal M_q(z)=\int_0^\infty y^{z-1}\sG_{0;(1,-1)}(1,y)^q\,dy$, and Cor.~\hb{19.40}{153} writes it as a nested lattice of Gamma data. Worked out:
\begin{itemize}
\item \emph{One mark} (Lemma~\hb{19.27}{148}(i) at $\nu=0$): $\sG_{0;(1,-1)}(1,y)=\sum_{m\ge0}\frac{y^m}{m!^2}\big(2(\log\Gamma)'(1+m)-\log y\big)$.
\item \emph{Two marks} (Thm.~\hb{19.39}{153}(i)): $\mathcal M_2(z)=\Gst41{z,\ z,\ z,\ z}{2z}$.
\item \emph{One mark against the rest} (Thm.~\hb{19.39}{153}(ii), $D=-\partial_z$):
\[
D^i\mathcal M_q(z)=\sum_{m\ge0}\frac1{m!^2}\big(2(\log\Gamma)'(1+m)\,D^i\mathcal M_{q-1}(z+m)+D^{i+1}\mathcal M_{q-1}(z+m)\big).
\]
\end{itemize}
Three times, summed from the inside out at the root rates $\frac14,\frac19,\frac1{16}$ (Thm.~\hb{19.39}{153}(iii)):
\begin{multline*}
\int_0^\infty\sG_{0;(1,-1)}(1,y)^5\,dy\\
=\sum_{a\ge0}\frac1{a!^2}\sum_{b\ge0}\frac1{b!^2}\sum_{c\ge0}\frac1{c!^2}\prod_{x\in\{a,b,c\}}\big(2(\log\Gamma)'(1+x)+D\big)\,\Gst41{z,\ z,\ z,\ z}{2z}\bigg|_{z=1+a+b+c}.
\end{multline*}
\emph{Its slots at integer data.} $2(\log\Gamma)'(1+x)=-2\,\Tx1\IZ{x}$, regularized (Prop.~\hb{5.14}{33}). With $N=1+a+b+c$, the jets of the pair's word are its value times a polynomial in the tails of Prop.~\hb{5.13}{32},
\[
\log\frac{\mathcal M_2(N-\delta)}{\mathcal M_2(N)}=\sum_{i\ge1}\frac{\delta^i}i\Big(4\,\Tx{i}{\IZ}{N-1}-2^i\,\Tx{i}{\IZ}{2N-1}\Big),
\]
and for $i\ge2$ the cosets (Thm.~\hb{5.12}{32}(i)) split the second tail, $2^i\,\Tx{i}{\IZ}{2N-1}=\Tx{i}{\IZ}{N-1}+\Tx{i}{\Ilat{\frac12+\mathbb Z}}{N-\frac12}$; at $i=1$ the regularized values carry $\log 2$ as well. So every point of the nest carries tails of $\mathbb Z$ and of $\frac12+\mathbb Z$ cut at the running points $a,b,c,N-1,N-\frac12$, against the live densities $1/x!^2$ on the increments and $\Gst41{z,\ z,\ z,\ z}{2z}$ at the end, the Gramma input of Thm.~\hb{6.10}{46} with the recurrence $2(2z+1)f(z+1)=z^3f(z)$. A nest of reads is one Trailmoment on the joint source (Thm.~\hb{5.21}{36}(ii)); opening each tail as a trail variable makes it a trail of lattice inputs on $\mathbb Z$ and $\frac12+\mathbb Z$ (Def.~\hb{5.23}{37}, Thm.~\hb{5.24}{37}(iv)), with live densities. The five-mark centre \emph{is} a trail. The coset $\frac12$ is the one of four marks, and it comes from the $\Gamma(2z)$ of the pair's word.

\emph{The order of the reads is part of the trail.} Summed flat over the three increments, the terms on the shell $a+b+c=N$ grow like $(9/4)^N$. Inside out each level converges, but it cancels: its largest term exceeds its value by about $9^z$ for three marks read from two and $4^z$ for four read from three (computed). So each level is summed to $10^{-(\text{loss}+70)}$ of its largest term, at $1100$ bits.

\emph{Values} (computed with the nest alone, \texttt{five\_mark\_trail.py}, in $2.6$ seconds):
\begin{align*}
\mathcal M_5(1)&=39.94558801199624569894402849842399589184\ldots,\\
\mathcal M_5(2)&=0.29469610741827513139030320865599973007\ldots,\\
\mathcal M_5(3)&=0.01893533579656526200673898655484419383\ldots,\\
\mathcal M_5(4)&=0.00364887571944556684733283235472259553\ldots.
\end{align*}
Checks, all to $10^{-70}$ or better: the same nest on three marks gives $3\,\Tx{2}{(\Ilat{1+3\mathbb Z}-\Ilat{2+3\mathbb Z})}{0}$, and on four marks $7\,\Tx3\IZ0$; the pair-first nest (a pair against a triple, root rate $\frac49$) gives the same $\mathcal M_5(1)$; the corner relation $5\mathcal M_5(1)-285\mathcal M_5(2)+225\mathcal M_5(3)=120$ and the ray law $\mathcal M_5(1)-201\mathcal M_5(2)+1062\mathcal M_5(3)-225\mathcal M_5(4)=0$ of the five-mark centre ray (from the parallel session) hold.

\emph{What is open, in these terms.} On four marks the live trail collapses, by the pair closing of Example~6.3 of the note, to the single lattice Trailmoment $\Tx{3}{\Ilat{\frac12+\mathbb Z}}{0}$. On five marks the same construction is a trail with live densities on $\mathbb Z$ and $\frac12+\mathbb Z$. Whether it collapses to trails with rational densities is the reduction question of \S14.12. Either way the value is a trail. \emph{Correction}: \S14.11 and \S14.12 discussed this value without Cor.~\hb{19.40}{153}, which already writes it as a nested lattice of Gamma data.

\subsection{Five marks shortened as far as the notation goes}\status{derived; the structure checked in exact rational arithmetic through $N=12$; the offset jet off the edge not yet computed}

\emph{Is it a serif closed form?} Yes. Write the trail of jet tails of the parallel session,
\[
\big(\sT[K]\,w_1\mathbb I^{(q_1)}_{\mathbb N}\cdots w_d\mathbb I^{(q_d)}_{\mathbb N}\big)(N):=g_1(N),\qquad g_{d+1}=K,\quad g_i(y)=\sum_{l\in\mathbb N}\partial_\varepsilon^{\,q_i}\big[w_i(l-\varepsilon)\,g_{i+1}(y+l-\varepsilon)\big]_{\varepsilon=0},
\]
with the words $w_1(x)=\Gamma(1+x)^{-2}$ and $w_2(x)=\Gst14{1+2x}{1+x,\ 1+x,\ 1+x,\ 1+x}$ and the kernel $M_2=\mathcal M_2$. Then $\int_0^\infty\sG_{0;(1,-1)}(1,y)^5\,dy=\big(\sT[M_2]\,w_2\mathbb I^{(2)}_{\mathbb N}\,w_1\mathbb I^{(1)}_{\mathbb N}\big)(1)$, a Trailmoment with its slots filled. The question is whether it \emph{shortens}, as four marks shortens to $\Tx3{\Ilat{\frac12+\mathbb Z}}0$.

\emph{One density on the quadrant} (the shift of Prop.~\hb{1.3}{11}):
\begin{gather*}
w_2(x)\,w_1(y)\,M_2(1+x+y)=\Phi(x,y)=\frac{\binom{x+y}x^4}{(1+2x+2y)\binom{2x+2y}{2x}\binom{2y}y},\\
\int_0^\infty\sG_{0;(1,-1)}(1,y)^5\,dy=-\sum_{l\ge0}\sum_{k\ge0}(\partial_x^2\partial_y\Phi)(l,k),
\end{gather*}
the inner sum over $k$ first. On the edge, $\Phi(x,0)=1/(1+2x)$, the density of four marks: five marks is the four-mark trail with its density pushed off the edge into the quadrant.

\emph{The pair collapses in every slice.} At an integer $y=k$, $\Phi(x,k)=k!^{-2}R_{k+1}(x)$ with the rational word
\[
R_N(x)=w_2(x)M_2(N+x)=\frac{(1+x)_{N-1}^4}{(1+2x)_{2N-1}}.
\]
Its integer poles cancel against the fourth powers; it has simple poles at $x=-j-\frac12$, $0\le j\le N-1$, with residues $c_{N,j}=2^{7-8N}\binom{2j}j^3\binom{2m}m^3j!^2m!^2$, $m=N-1-j$; it is odd under $x\mapsto-N-x$; and its polynomial part $P_N$ has degree $2N-3$. The pair's second jet moves every pole to index $3$ (the jet law), the oddness kills the even columns (Thm.~\hb{6.6}{44}(iv)), the column sum is read by Thm.~\hb{6.4}{43}, and $P_N$ is paid on the far line (the parallel session's Theorem~B). So the four-mark ray at every integer offset is
\begin{gather*}
M_4(N)=A_N\,\Tx3{\Ilat{\frac12+\mathbb Z}}0-B_N,\\
A_N=2\sum_jc_{N,j}=\frac{(N-1)!^2}{2^{8N-8}}\sum_{j+m=N-1}\frac{\binom{2j}j^3\binom{2m}m^3}{\binom{j+m}j^2},\\
B_N=2\sum_{j=0}^{N-1}c_{N,j}\,\Th{3}{\Ilat{\frac12+\mathbb Z}}{0}{j}-\sum_{i\ge1}(-1)^i\mu_i\,\frac{P_N^{(i)}(-1)}{i!},
\end{gather*}
with $\mu_i=i\,\nu_{i-1}$ and $\nu_k=(-1)^{k-1}k\,\Tx{1-k}{\IZ}{0}$ for $k\ge1$, $\nu_0=1$: the far line is paid in counting Trailmoments at nonpositive index (Prop.~\hb{5.5}{29}). Exact checks (\texttt{four\_mark\_word.py}): oddness, poles and the formula for $A_N$ through $N=12$; $A_2,\dots,A_5=\frac1{16},\frac7{256},\frac9{256},\frac{6111}{65536}$ and $B_2,\dots,B_5=\frac38,\frac{27}{128},\frac{37}{128},\frac{25555}{32768}$, the parallel session's table. The two indices $j$ and $m$ count the poles from the two ends of the odd word.

\emph{The order-dependence, located.} The slice $k$ has a polynomial part of degree $2k-1$. On the edge, $k=0$, it is zero, and the pair collapses with no far line; that is all of four marks. Off the edge every slice collapses only with its far line. Summed mark-first, nothing diverges and nothing collapses; summed pair-first, every slice collapses and every slice off the edge pays a far line.

\emph{The mark does not collapse.} With the mark outermost (Thm.~\hb{19.39}{153}(ii)),
\[
\int_0^\infty\sG_{0;(1,-1)}(1,y)^5\,dy=\sum_{k\ge0}\frac1{k!^2}\Big(2(\log\Gamma)'(1+k)\,M_4(1+k)-M_4'(1+k)\Big).
\]
The mark's word $k!^{-2}$ cancels the $(N-1)!^2$ of $A_N$ exactly, so the first half puts on $\Tx3{\Ilat{\frac12+\mathbb Z}}0$ the coordinate
\[
-2\sum_{j,m\ge0}\frac{\binom{2j}j^3\binom{2m}m^3}{2^{8(j+m)}\binom{j+m}j^2}\,\Tx1\IZ{j+m},
\]
a Trailmoment on $\mathbb N^2$ whose variables are the pole columns of the two pairs, the mark between them entering as $\binom{j+m}j^{-2}$ and as the tail at $j+m$; the regularized index-one values cancel against the second half by balance. The second half is the offset jet,
\[
\partial_NR_N(x)=R_N(x)\big(4\psi(N+x)-2\psi(2N+2x)\big)=R_N(x)\big(3\psi(1+x)-\psi(\tfrac12+x)-2\log2+\rho_N(x)\big),
\]
by duplication, with $\rho_N$ rational and odd under $x\mapsto-N-x$: $R_N\rho_N$ is even, so it lands on index $4$ of $\Ilat{\frac12+\mathbb Z}$, and the digamma terms make trails of depth two on $\mathbb Z$ and $\frac12+\mathbb Z$. For $N\ge2$ the far line of $P_N$ against the digamma terms carries logarithms; that piece is not yet computed.

So the five-mark centre is the content of four marks with live coordinates: double Trailmoments of explicit Gamma-ratio densities on $\mathbb N^2$. On the edge the coordinates are rational because there is one slice. Whether the trail shortens is now a question about these coordinates, not about an integral.

\subsection{What five marks is, without collapsing}\status{derived; the edge identity checked against the nest; the anatomy computed}

The trail does not shorten, and comparison does not need it to: both centres are members of one family with known relations, and five marks is built out of four.

\emph{Five marks is four marks read through one more mark} (Thm.~\hb{19.39}{153}(ii), the mark outermost, $D=-\partial_\lambda$):
\[
\int_0^\infty\sG_{0;(1,-1)}(1,y)^5\,dy=\sum_{k\ge0}\frac1{k!^2}\Big[\big(2(\log\Gamma)'(1+k)+D\big)M_4\Big](1+k).
\]
The fifth mark's lattice walks along the four-mark ray at the offsets $1,2,3,\dots$, with the weights $k!^{-2}$ and one jet; along that ray $M_4(N)=A_N\,\Tx3{\Ilat{\frac12+\mathbb Z}}0-B_N$ (\S14.14).

\emph{The edge carries the value.} The step $k=0$ reads the four-mark centre itself. Inside the strip the pair level holds without a far line, so with
\[
h(x)=3\big((\log\Gamma)'(1+x)-(\log\Gamma)'(1)\big)-\big((\log\Gamma)'(\tfrac32+x)-(\log\Gamma)'(\tfrac12)\big),
\]
whose values at the integers are the heads $3\,\Th1\IZ0l-\Th1{\Ilat{\frac12+\mathbb Z}}0{l+\frac12}$ (the regularized index-one values and $\log2$ cancel by balance and duplication),
\[
\text{edge}=-\sum_{l\ge0}\partial_x^2\Big[\frac{h(x)}{1+2x}\Big]_{x=l}
=\sum_{l\ge0}\Big[\frac{2\,\tau_3(l)}{1+2l}+\frac{4\,\tau_2(l)}{(1+2l)^2}-\frac{8\,h(l)}{(1+2l)^3}\Big],
\]
with $\tau_s(l)=3\,\Tx{s}{\IZ}{l}-\Tx{s}{\Ilat{\frac12+\mathbb Z}}{l+\frac12}$. This is a trail of depth two on $\mathbb Z$ and $\frac12+\mathbb Z$ with rational coefficients, of weight four: the four-mark word $1/(1+2x)$, of weight three, carrying the mark's lattice coefficient $h$. In the slots of Def.~\hb{5.1}{28}, with $x=l+\frac12$, set against four marks:
\begin{align*}
M_4(1)&=\Tx3{\Ilat{\frac12+\mathbb Z}}0,\\
\text{edge}&=\big(\sT_{3;(1)}[x]\,(-h(x-\tfrac12))\,\Ilat{\frac12+\mathbb Z}\big)((0,\infty))\\
&\quad+\big(\sT_{2;(1)}[x]\,\tau_2(x-\tfrac12)\,\Ilat{\frac12+\mathbb Z}\big)((0,\infty))+\big(\sT_{1;(1)}[x]\,\tau_3(x-\tfrac12)\,\Ilat{\frac12+\mathbb Z}\big)((0,\infty)).
\end{align*}
Kernel, lattice and cut are those of four marks. The fifth mark enters the input slot as a density, its own Trailmoment at the running point, and the pair's second jet is shared between the kernel and that density: the indices split as $3+1$, $2+2$, $1+3$, always four. Its value is $39.38158212\ldots$ (the nest; the explicit sum to $l=40000$ agrees to the size of its tail).

\emph{The live part is a small tail.} The steps $k\ge1$ read the four-mark ray beyond its centre and add $0.5475$, $0.0159$, $0.0006$, $2.5\cdot10^{-5},\dots$, each at most $\frac1{16}$ of the one before; together $0.5640\ldots$. That is where the live coordinates of \S14.14 sit. So the difference between four and five marks is located: on the edge, five marks is four marks one weight up; off the edge the mark's weights $k!^{-2}$ read the rest of the four-mark ray, and that is the new part.

\emph{Its place on its own ray.} The four-mark centre ray is $\mathbb Q+\mathbb Q\,\Tx3{\Ilat{\frac12+\mathbb Z}}0$: one master. The five-mark ray law has order three, and the corner relation $5\mathcal M_5(1)-285\mathcal M_5(2)+225\mathcal M_5(3)=120$ leaves two masters, $\mathcal M_5(1)$ and $\mathcal M_5(2)$ (the parallel session; checked here to $10^{-70}$). The five-mark centre generates its ray as $7\,\Tx3\IZ0$ generates the four-mark one.

\subsection{What the five-mark trail implies}\status{derived; the ray law re-derived here from the Gramma's relations}

Read off $M_5(s)=\big(\sT[M_2]\,w_2\mathbb I^{(2)}_{\mathbb N}\,w_1\mathbb I^{(1)}_{\mathbb N}\big)(s)$ without evaluating or collapsing it.
\begin{enumerate}[label=(\roman*)]
\item \emph{The corner.} Only the first term of the trail, $l=k=0$, has a pole at $s=0$: the kernel's $M_2(u)=2u^{-3}+\cdots$, raised by the jets of orders $2$ and $1$, so $M_5(s)=2\cdot3\cdot4\cdot5\,s^{-6}+\cdots=5!\,s^{-6}+\cdots$. With jets of total order $n-2$, $M_n(s)=n!\,s^{-n-1}+\cdots$.
\item \emph{The law along the ray.} The kernel and the words are the share side of $g=\sG_{0;(1,-1)}(1,\cdot)$, and the ladder and exchange of Thm.~\hb{2.3}{13} give $\vartheta^2g=yg$, $\vartheta=y\partial_y$. The fifth power of a solution obeys
\[
\big[\vartheta^6-y(35\vartheta^4+70\vartheta^3+63\vartheta^2+28\vartheta+5)+y^2(259\vartheta^2+518\vartheta+285)-225y^3\big]g^5=0,
\]
and in the index
\begin{multline*}
s^6M_5(s)-(35s^4+70s^3+63s^2+28s+5)M_5(s+1)\\
+(259s^2+518s+285)M_5(s+2)-225M_5(s+3)=0
\end{multline*}
(derived here; it agrees with the parallel session).
\item \emph{The relation at the corner.} At $s\to0$ the law and (i) give $5M_5(1)-285M_5(2)+225M_5(3)=5!$. Four marks, the same way: $4M_4(1)-64M_4(2)=4!$, which $A_2=\frac1{16}$, $B_2=\frac38$ of \S14.14 satisfy exactly.
\item \emph{Masters.} The law has order three and its last coefficient is $-225\neq0$, so every value at an integer index is a rational combination of $1$, $M_5(1)$, $M_5(2)$: two masters. The four-mark ray has one. The next coefficient of the corner, $-240\,\Tx1\IZ0$ at $s^{-5}$ (regularized), gives the jets $M_5'$ the same treatment, with that one new constant.
\end{enumerate}
From four marks to five, read off the trails alone: the corner goes from $4!\,s^{-5}$ to $5!\,s^{-6}$, the jets from total order $2$ to $3$, and the masters of the ray from one to two.

\emph{The masters and the levels} (the two counts derived; that they agree is a count agreement, not a mechanism: see \S14.17). On $n$ marks the law along the centre ray has order $\lceil n/2\rceil$, one for each nonzero critical value $(n-2j)^2$ (the degree in $y$ of the $n$-th power of $\vartheta^2-y$), and the corner gives exactly one relation among the values at integer indices; the finite parts at the other poles are new constants, not relations. So over $\mathbb Q$ the ray is spanned by $1$ and $\lceil n/2\rceil-1$ masters. The shortest trail of $n$ marks has $\lceil n/2\rceil-1$ levels: $(2,\dots,2,1)$ for odd $n$, all pairs for even $n$. The two counts agree for every $n$: four marks, one level and one master; five marks, two levels and two masters. With the parallel session's values: $M_4(1)-16M_4(2)=8.41439832211715999780-16\cdot0.15089989513232249986=6$, while $5M_5(1)-285M_5(2)+225M_5(3)=120$ ties three values and leaves $M_5(1)$ and $M_5(2)$ unrelated by the law over $\mathbb Q$; what the law cannot see, and what is left once more is admitted, is \S14.17.

\subsection{What the ray law cannot see: the jet slot of the outer pair}\status{derived; the five-mark settings, their block and the minors computed in ball arithmetic, no quadrature (\texttt{jet\_slot.py}); the jet identities of the pair's word checked exactly through $m=39$; the parallel session's relation identified as a minor; its Trailmoment on $\mathbb Z^2$ recalled}

The question: why the five-mark ray needs two numbers and the four-mark ray one. The answer below is a statement about what the law along the ray can see, read inside the trail.

\emph{A level before it meets its kernel.} For a word $w$ and a jet order $q$ put
\[
(w\,\mathbb I^{(q)}_{\mathbb N})(y):=\sum_{m\in\mathbb N}\partial_\varepsilon^{\,q}\big[w(m-\varepsilon)\,y^{m-\varepsilon}\big]_{\varepsilon=0}.
\]
A level of a trail is this function closed against its kernel,
\[
\big(\sT[K]\,w\,\mathbb I^{(q)}_{\mathbb N}\big)(s)=\int_0^\infty y^{s-1}(w\,\mathbb I^{(q)}_{\mathbb N})(y)\,k(y)\,dy,
\]
when $K$ is the index reading of $k$ and the closing holds (Lemma~\hb{19.29}{148}). The mark is the word $w_1$ at jet one: $(w_1\mathbb I^{(1)}_{\mathbb N})(y)=\sG_{0;(1,-1)}(1,y)=g(y)$ (Lemma~\hb{19.27}{148}(i) at $l=0$). At jet zero it is $(w_1\mathbb I^{(0)}_{\mathbb N})(y)=\sum_my^m\Gamma(1+m)^{-2}$, of reach $1$ (Def.~\hb{19.26}{148}) and at most $\big(\sum_my^{m/2}/m!\big)^2=e^{2\sqrt y}$.

\emph{The law does not see the jet} (derived). Since $w_1(x)/w_1(x-1)=x^{-2}$,
\[
(\vartheta^2-y)\sum_{m\in\mathbb N}w_1(m-\varepsilon)\,y^{m-\varepsilon}=\frac{\varepsilon^2\,y^{-\varepsilon}}{\Gamma(1-\varepsilon)^2},
\]
so the mark at jets $0$ and $1$ both obey $\vartheta^2f=yf$, the law of Thm.~\hb{2.3}{13}. The pair's word is the squared law read in the index: $w_2(x)/w_2(x-1)=2(2x-1)/x^3$, and $\vartheta^3-2y(2\vartheta+1)$ annihilates every product of two solutions of $\vartheta^2f=yf$, so
\[
\big(\vartheta^3-2y(2\vartheta+1)\big)\sum_{m\in\mathbb N}w_2(m-\varepsilon)\,y^{m-\varepsilon}=-\varepsilon^3\,w_2(-\varepsilon)\,y^{-\varepsilon}.
\]
Jets $0,1,2$ are therefore solutions; a solution is fixed by its coefficients of $y^0(-\log y)^j$, $j\le2$; matching them,
\[
(w_2\mathbb I^{(2)}_{\mathbb N})=g^2,\qquad (w_2\mathbb I^{(1)}_{\mathbb N})=g\cdot(w_1\mathbb I^{(0)}_{\mathbb N}),\qquad (w_2\mathbb I^{(0)}_{\mathbb N})=(w_1\mathbb I^{(0)}_{\mathbb N})^2 .
\]
Coefficientwise these are identities of harmonic numbers, $\sum_aH_a\binom ma^2=\binom{2m}m(2H_m-H_{2m})$ and its second-order companion, checked exactly through $m=39$. The jet order of a level counts its marks at jet one. One and two marks are the laws of degree one in $y$, which is what lets a level's word be a single stacked Gamma.

\emph{The settings of the outer pair.} At five marks lower the jet of the outer pair, and write $\sT^{(q)}$ for the trail with that slot at $q$:
\[
\sT^{(q)}(s):=\big(\sT[M_2]\,w_2\mathbb I^{(q)}_{\mathbb N}\,w_1\mathbb I^{(1)}_{\mathbb N}\big)(s)=\int_0^\infty y^{s-1}\,g(y)^3\,(w_2\mathbb I^{(q)}_{\mathbb N})(y)\,dy,\qquad q=2,1,0,
\]
with integrands $g^5$, $g^4\,(w_1\mathbb I^{(0)}_{\mathbb N})$, $g^3\,(w_1\mathbb I^{(0)}_{\mathbb N})^2$; $\sT^{(2)}=M_5$. Each integrand is a product of five solutions of $\vartheta^2f=yf$, so each obeys the fifth-power law of \S14.16(ii), and, decaying at infinity, the law along the ray. The law is blind to $q$.

\emph{The corner sees only the full jet} (derived). At $s=0$ the kernel $M_2(s)=2s^{-3}+\cdots$ is raised one order by each jet, so $\sT^{(q)}$ has its pole of order $4+q$, and $s^6\sT^{(q)}(s)$ tends to $5!$, $0$, $0$ for $q=2,1,0$. The law at $s=0$ reads that limit:
\[
5\,\sT^{(q)}(1)-285\,\sT^{(q)}(2)+225\,\sT^{(q)}(3)=\begin{cases}5!,&q=2,\\0,&q=1,0.\end{cases}
\]

\emph{Which settings close} (derived). The pair's lattice has reach $(1+1)^2=4$ at every jet (Lemma~\hb{19.28}{148}); $k$ marks beneath it are at most $Cy^ae^{-2\sqrt{k^2y}}$ (Lemma~\hb{19.27}{148}(ii)). Lemma~\hb{19.29}{148} closes the pairing when $4<k^2$, its tail at most $C(1+N)^D(4/k^2)^N$: the condition $w_{S_1}<w_{S_2}$ of Thm.~\hb{19.39}{153}(ii), now at every jet.
\begin{itemize}
\item Five marks, $k=3$: $4<9$. Every jet closes, at the root rate $4/9$: three settings on the ray.
\item Four marks, $k=2$: $4=4$, the case of the commentary to Thm.~\hb{19.38}{152} where the points meet and domination decides nothing. The jet decides. At $s=1$, $w_2(x)M_2(1+x)=1/(1+2x)$, and
\[
\big(\sT[M_2]\,w_2\mathbb I^{(q)}_{\mathbb N}\big)(1)=\sum_{m\ge0}\frac{q!\,2^q}{(1+2m)^{q+1}}=\frac{q!}2\,\Tx{q+1}{\Ilat{\frac12+\mathbb Z}}0 :
\]
one Trailmoment on $\frac12+\mathbb Z$ with its index slot at $3,2,1$. Index $3$ is $M_4(1)$; index $2$ is $\frac12\Tx2{\Ilat{\frac12+\mathbb Z}}0=\frac14\Gamma(\frac12)^4$; index $1$ diverges. Two settings on the ray.
\end{itemize}
In general the setting with $j$ marks at jet zero is at most $Cy^ae^{-2(n-2j)\sqrt y}$ and closes when $j<n/2$; at four marks $j=n/2$ is the divergent index one.

\emph{The count} (derived; independence computed at five marks). The closing settings are independent solutions of the law: at four marks their $2\times2$ minor at $s=1,2$ is $\frac3{32}\Gamma(\frac12)^4$ (the two corners and the value $\frac14\Gamma(\frac12)^4$), at five marks their $3\times3$ minor at $s=1,2,3$ is not zero (below). There are as many of them as the order of the law, so they span its solutions on the integer ray, and the corner, $n!$ on the centre and $0$ on every lowered setting, fixes the centre up to adding lowered settings. Over $\mathbb Q$:
\[
\#\text{masters}=\#\{j:1\le j<n/2\}=\lceil n/2\rceil-1 .
\]
From four marks to five the new master is the pair with both marks at jet zero: on two marks it balances (index one on $\frac12+\mathbb Z$); on three it closes.

\emph{Computed} (\texttt{jet\_slot.py}, $1400$ bits, every sum cut at $10^{-75}$ of its largest term, $5.3$\,s).
\begin{itemize}
\item $\sT^{(2)}(s)-M_5(s)$, $s=1,\dots,4$, against the mark-first nest: below $10^{-73}$. $\sT^{(1)}(1)$ against the other nesting $\sum_kM_4(1+k)/k!^2$: $10^{-74}$.
\item The law at $s=1$ for $q=2,1,0$: below $10^{-71}$. The corner: $120$, $0$, $0$, to $10^{-71}$.
\item Four marks, by the other nesting: $\sum_kM_3(1+k)/k!^2-\frac14\Gamma(\frac12)^4$ and the corner of that setting, below $10^{-72}$.
\end{itemize}

\emph{The lowered settings close on stacked Gammas at fifteenths} (computed; the values at $s=1$ are the parallel session's, re-verified here). With
\[
C=\frac1{240\sqrt5}\,\Gst44{\frac1{15},\ \frac2{15},\ \frac4{15},\ \frac8{15}}{\frac12,\ \frac12,\ \frac12,\ \frac12},
\]
\begin{align*}
\sT^{(1)}(1)&=8\,\Gamma(\tfrac12)^4\,C, & \sT^{(1)}(2)&=\tfrac{104}{225}\,\Gamma(\tfrac12)^4\,C-\tfrac4{1125}\,\Gamma(\tfrac12)^4\,C^{-1},\\
\sT^{(0)}(1)&=2\sqrt{15}\,\Gamma(\tfrac12)^2\,C, & \sT^{(0)}(2)&=\tfrac{26\sqrt{15}}{225}\,\Gamma(\tfrac12)^2\,C+\tfrac{\sqrt{15}}{1125}\,\Gamma(\tfrac12)^2\,C^{-1},
\end{align*}
each below $10^{-72}$. The block $\sT^{(1)}(1)\,\sT^{(0)}(2)-\sT^{(1)}(2)\,\sT^{(0)}(1)=\frac{16\sqrt{15}}{1125}\Gamma(\frac12)^6$, and the corners make the $3\times3$ minor of the three settings at $s=1,2,3$ equal to $\frac8{15}$ of it (derived; checked to $10^{-72}$). The settings at $s\ge3$ follow by the law with rational coefficients.

\emph{The two masters against the lowered settings.} Over the scalars of the lowered settings, $(M_5(1),M_5(2))$ is equivalent to the two minors
\[
\Delta_q=M_5(1)\,\sT^{(q)}(2)-M_5(2)\,\sT^{(q)}(1),\qquad q=1,0,
\]
the change of coordinates having the block above as its determinant. The parallel session's relation is the minor at $q=0$, exactly:
\[
\Big(13C+\frac1{10C}\Big)M_5(1)-225C\,M_5(2)=\frac{15\sqrt{15}}{2\,\Gamma(\frac12)^2}\,\Delta_0
\]
(computed: its coefficient ratio is $\sT^{(0)}(2)/\sT^{(0)}(1)$ to $10^{-74}$, and resolved on the two minors its coefficient on $\Delta_1$ is below $10^{-72}$). Its value $\frac{675\sqrt{15}}{\Gamma(\frac12)^6}L(4)$, with $L$ the parallel session's Trailmoment on $\mathbb Z^2$ through the norm forms of discriminant $-15$, is recalled, not recomputed here; so $\Delta_0=\frac{90}{\Gamma(\frac12)^4}L(4)$ (recalled). The other minor,
\[
\Delta_1=4.340002931535662034070630216892006273160\ldots,
\]
is not identified, and not yet compared against anything. At four marks the one minor of the centre against its lowered setting is fixed by the corners, $M_4(1)\cdot\frac1{64}\Gamma(\frac12)^4-M_4(2)\cdot\frac14\Gamma(\frac12)^4=\frac3{32}\Gamma(\frac12)^4$, so there the lowered setting takes no master away.

\emph{Relative to what} (derived from the above; the $\mathbb Z^2$ entry recalled). What the law, its corner and the identified members leave open, by the scalars admitted. These are counts of open directions, not statements of transcendence.
\begin{center}
\begin{tabular}{lcc}
scalars admitted & five marks & four marks\\\hline
$\mathbb Q$, the law's own & two: $M_5(1)$, $M_5(2)$ & one: $M_4(1)$\\
$+$ the lowered settings & two: $\Delta_1$, $\Delta_0$ & one: $M_4(1)$\\
$+$ Trailmoments on $\frac12+\mathbb Z$, $\mathbb Z^2$ & one: $\Delta_1$ & none\\
the serif family & none & none
\end{tabular}
\end{center}
At four marks the third row closes because $M_4(1)=\Tx3{\Ilat{\frac12+\mathbb Z}}0$; in the last row every value on either ray is its trail with the slots filled, $M_5(s)=\sT^{(2)}(s)$.

``Two masters'' is the first row. The law and its corner fix the centre up to the settings they cannot see; the trail leaves nothing open. What the trail does not yet shorten is one number, the minor $\Delta_1$ of the centre against the setting with one mark of the outer pair at jet zero.

\emph{On \S14.16.} ``The masters count the levels'' is a count agreement, not a mechanism: at five marks both lowered settings sit in the jet slot of the one outer pair.

\subsection{The machine in general: the span of the mark's exchange}\status{derived on one scale; the counts for $(1,-1)$ derived; for $(2,-1)$ the mark's law derived and checked in the index, the counts of masters not yet computed}

Two recurrences are in play: the exchange relation of one mark, in the mark's own index $\nu$ (Thm.~\hb{2.3}{13}(ii)), and the law along the centre's ray, in the index $s$. The first causes the second, in five moves.
\begin{enumerate}[label=(\roman*)]
\item \emph{The mark's span.} On one scale with integer exponents $t=(t_1,\dots,t_N)$, the exchange relation $\nu\,\sG_{\nu;t}=\sum_kt_k\xi_k\,\sG_{\nu+t_k;t}$ links the indices from $\nu+\min(0,\min t)$ to $\nu+\max(0,\max t)$. Its span is
\[
w(t)=\max(0,\max_kt_k)-\min(0,\min_kt_k).
\]
\item \emph{The mark's law.} Let the shared scale $y$ be the coefficient of the exponent $-1$, as in a circuit. The ladder (Thm.~\hb{2.3}{13}(i)) makes each step down in the index one derivative in $y$, and the exchange relation becomes a linear law in $y$ of order $w(t)$. The mark has $w(t)$ independent solutions: the Gramma and $w(t)-1$ partners.
\item \emph{$n$ marks.} Every product of $n$ solutions obeys the $n$-th power of the mark's law, which never reads which solution sits in which slot.
\item \emph{The ray.} The index reading $\int_0^\infty y^{s-1}(\cdot)\,dy$ sends $\vartheta$ to $-s$ and $y$ to the shift $s\mapsto s+1$. Every product whose reading converges obeys the resulting law in $s$.
\item \emph{The corner.} At $s=0$ the leading coefficient of that law vanishes, and the law reads the top pole. Only the Gramma's own power carries a $-\log y$ from every mark, so the corner is one condition. The masters over $\mathbb Q$ are the solutions of the law in $s$ beyond the centre, one fewer than the law's order.
\end{enumerate}

\emph{One exponent.} $w=1$: $\sG_{\nu;(1)}(\xi)=\Gamma(\nu)\,\xi^{-\nu}$ is a bare power with no partner, and a scale integral of a product of such marks diverges at one end. A centre needs a mark with exponents of both signs, and both signs already give $w\ge2$.

\emph{Two exponents, $t=(1,-1)$: the circuit} (Cor.~\hb{17.20}{113}). Here $w=2$, with the solutions $g$ and $h=\sum_my^m\Gamma(1+m)^{-2}$ of growth $e^{\mp2\sqrt y}$. The product $g^{n-j}h^j$ has growth $e^{-2(n-2j)\sqrt y}$ and converges exactly when $j<n/2$; these products are exactly the solutions of the law in $s$ (\S14.17).
\begin{center}
\begin{tabular}{ccc}
$n$ & products that converge & masters over $\mathbb Q$\\\hline
$1$ & $g$ & $0$\\
$2$ & $g^2$ & $0$\\
$3$ & $g^3,\ g^2h$ & $1$\\
$4$ & $g^4,\ g^3h$ & $1$\\
$5$ & $g^5,\ g^4h,\ g^3h^2$ & $2$\\
$6$ & $g^6,\ g^5h,\ g^4h^2$ & $2$
\end{tabular}
\end{center}
The rows without masters are the collapses $M_1(s)=\Gamma(s)^2$ and $M_2(s)=\Gst41{s,\ s,\ s,\ s}{2s}$ (Thm.~\hb{19.39}{153}(i)). A master is added at each odd $n$: from even $n$ to $n+1$ the tie $g^{n/2}h^{n/2}$, which does not decay, becomes $g^{n/2+1}h^{n/2}$, which does.

\emph{Span three, $t=(2,-1)$} (derived). With $G_\mu=\sG_{\mu;(2,-1)}(1,y)$ the exchange relation is $\mu\,G_\mu=2G_{\mu+2}-y\,G_{\mu-1}$, spanning $\mu-1,\dots,\mu+2$. At $\mu=\nu-2$, with the ladder, it gives for $f=G_\nu$
\[
y\,f'''-(\nu-2)\,f''+2f=0,
\]
of order three, whose index reading is $2F(s+2)=s(s+1)(s+\nu)\,F(s)$ with $F(s)=\frac12\Gamma(s)\Gamma(\frac{s+\nu}2)$ (checked). At large $y$ the leading balance $y\,\varphi'^3=-2$ gives one decaying solution and two growing ones. Two things change against $(1,-1)$: there are two partners, and the step of two splits the integer ray into two classes, each with its own inputs. At $\nu=0$ one mark already reads $\frac12\Gamma(\frac12)$ at $s=1$ and $\frac12$ at $s=2$. The counts of masters for $n$ such marks are not yet computed. For a table that is not a circuit, a mark sits in its letters with other exponents and on more scales, with one exchange relation per scale; not worked out here. The configuration $(2,-1)$ is a toy: no mark of a circuit has it (\S14.19).

\subsection{What a configuration carries}\status{derived; the counts of masters for the configuration $(2,1,-1,-2)$ not yet computed}

\emph{Signs are letters.} In a circuit (Cor.~\hb{17.20}{113}) a mark's variable enters the cobasis letter $\Psi=\ip\theta\alpha$ linearly, giving the exponent $+1$ and the kernel, and the basis letter $K=(\prod\alpha)\sum_e\theta_e/\alpha_e$ inversely, giving $-1$, the place where the shared scale sits. The reflection $t\mapsto-t$ (Prop.~\hb{2.4}{14}) is the inversion $\alpha\mapsto1/\alpha$, which exchanges the two letters. A single table's letters are multilinear, so a mark of a circuit always has the configuration $(1,-1)$: magnitudes beyond one would need a mark entering a letter to a higher power.

\emph{The index reading of one mark} (derived from Def.~\hb{2.1}{13}). If the shared scale $y$ is the coefficient of the exponent $t_y$ and $t'$ is the configuration of the other exponents, then
\[
\int_0^\infty y^{s-1}\,\sG_{\nu;t}(\dots,y,\dots)\,dy=\Gamma(s)\,\sG_{\nu-t_ys;\,t'}(\dots),
\]
since $\int_0^\infty y^{s-1}e^{-yx^{t_y}}dy=\Gamma(s)\,x^{-t_ys}$. For the circuit's mark this is $\Gamma(s)\,\sG_{\nu+s;(1)}(1)=\Gamma(s)\Gamma(s+\nu)$, the share side of Cor.~\hb{17.20}{113}. The exchange relation of $t'$ is the law of one mark along the ray.

\emph{Three readings of a configuration.}
\begin{enumerate}[label=(\roman*)]
\item \emph{Its span} $w(t)$ counts the mark's solutions (\S14.18).
\item \emph{What is left, $t'$}, decides whether the mark has a lattice at $y=0$. If $t'$ has one sign, $\sG_{\cdot;t'}$ has poles (Thm.~\hb{2.2}{13}), the mark has a lattice expansion, and where those poles meet the poles of $\Gamma(s)$ it carries $-\log y$. That log is what the corner reads. If $t'$ has both signs, $\sG_{\cdot;t'}$ is entire, $\Gamma(s)$ gives the only poles, and $y=0$ is an ordinary point of the mark's law. For the circuit's mark at the centre, $t'=(1)$ and $\nu=0$, so the poles meet everywhere: $\Gamma(s)^2$.
\item \emph{The far side} decides which solutions decay. The shared scale's exponent balances against the exponents of the other sign, and the balance fixes how fast each solution grows or decays, hence which products converge.
\end{enumerate}

\emph{The configuration $(2,1,-1,-2)$, shared scale on $-1$} (derived). With the other coefficients $1$, the exchange relation $\mu\,G_\mu=2G_{\mu+2}+G_{\mu+1}-y\,G_{\mu-1}-2G_{\mu-2}$ at $\mu=\nu-2$ and the ladder give
\[
2f''''-y\,f'''+(\nu-2)\,f''+f'-2f=0 ,
\]
of order four, and its index reading reduces exactly to the exchange relation of $t'=(2,1,-2)$ for $F(s)=\Gamma(s)\,\sG_{\nu+s;(2,1,-2)}(1,1,1)$. The leading coefficient is the constant $2$: $y=0$ is an ordinary point, the four solutions are power series there, and there is no $-\log y$ anywhere, as (ii) predicts, since $t'$ has both signs. At large $y$ the balance $y\,\varphi'^3=-2$ gives one decaying solution (the Gramma) and two growing ones of order $y^{2/3}$ in the exponent, and the balance $2\varphi'=y$ gives a fourth growing as $e^{y^2/4}$. So the circuit's mechanism, in which the corner sees only the full jet, is absent: every product has a value at $y=0$, and the corner reads all of them. The masters for $n$ such marks are not yet computed. Since the configuration has two exponents of each sign, removing any one leaves both signs, wherever the shared scale sits.

\emph{Where larger spans come from in the families.} A series bunch of $k$ marks is read by its room as one mark (Ch.~18; Prop.~\hb{19.34}{150}), its transform the product of $k$ Grammas of configuration $(1,-1)$. It has $k+1$ solutions, $g^{k-j}h^j$, of growth $e^{-2(k-2j)\sqrt y}$: a ladder of growths, not the star of $y^{2/3}$ branches of $(2,-1)$. The trail's levels are such bunches, of one and two marks.

\subsection{The ray read from the list}\status{the reading is Thm.~\hb{35.12}{310}; the dictionary derived; the sign values of the uniform lists computed; their tiers and any native gauges off the signs not yet computed}

\S\S14.13--14.19 read the neighbours of the centre through the Gramma of one list, $\binom{\cE}1$. The handbook builds them from any family. The centre ray at the strengths $\boldsymbol\sigma$ (Def.~\hb{35.9}{309}) has the values
\[
\langle1\rangle_k=\int_{y_{e_0}=0}f^k\,dy=P_F(k\sigma\one;k\boldsymbol\sigma,k\boldsymbol\sigma),\qquad f=\prod_j\big(K_j(e^y)K_j(e^{-y})\big)^{-\sigma_j},
\]
and for $\binom{\cE}1$ at $\sigma=1$ these are $\langle1\rangle_k=M_n(k)/\Gamma(k)^2$ (Cor.~\hb{17.20}{113}): the same numbers. What decides the law of the ray is read from the list: its native gauges include the signs $\eta\in\{\pm1\}^{\cE}$, $\eta_{e_0}=1$, with $K_j(\eta)\neq0$, of native value $\prod_jK_j(\eta)^{-2\sigma_j}$; distinct native gauges give independent solutions, and their native values are roots of the symbol of every recurrence of the ray (Thm.~\hb{35.12}{310} and the commentary before Def.~\hb{35.9}{309}).

\emph{Dictionary for $\binom{\cE}1$} (derived). $K(\eta)=\sum_e\eta_e=n-2j$ at a sign with $j$ marks at $-1$. The real gauge, all marks at $+1$, is the centre itself, $f(\one)=n^{-2}$; the partner $h$ of \S14.17 is a mark sent to $-1$; the product $g^{n-j}h^j$ is the sign with $j$ marks at $-1$, of native value $(n-2j)^{-2}$, the root found in \S14.16; the tie $g^{n/2}h^{n/2}$ at even $n$ is the sign with $K(\eta)=0$, which is not a native gauge. So the order $\lceil n/2\rceil$ of the banana's ray law is the number of distinct nonzero $|K(\eta)|$, and the two masters at five marks are the three sizes $5,3,1$ less the corner.

\emph{Other lists on the same marks} (computed). For $\binom{[n]}k$, $K(\eta)=\sum_i(-1)^i\binom ji\binom{n-j}{k-i}$ depends only on $j$:
\begin{center}
\begin{tabular}{lll}
list & $K(\eta)$, $j=0,1,\dots,n$ & distinct nonzero $|K(\eta)|$\\\hline
$\binom{[4]}1$ & $4,2,0,-2,-4$ & $4,2$\\
$\binom{[4]}2$ & $6,0,-2,0,6$ & $6,2$\\
$\binom{[5]}1$ & $5,3,1,-1,-3,-5$ & $5,3,1$\\
$\binom{[5]}2$ & $10,2,-2,-2,2,10$ & $10,2$\\
$\binom{[6]}1$ & $6,4,2,0,-2,-4,-6$ & $6,4,2$\\
$\binom{[6]}2$ & $15,5,-1,-3,-1,5,15$ & $15,5,3,1$\\
$\binom{[6]}3$ & $20,0,-4,0,4,0,-20$ & $20,4$
\end{tabular}
\end{center}
On five marks the banana's signs force three solutions and those of $\binom{[5]}2$ force two; on six marks the counts are three, four and two. The count belongs to the list, not to the number of marks, and the banana is the sequence of lists in which $K(\eta)$ is the progression $n-2j$. The signs give a lower bound: native gauges off the signs come in pairs (Thm.~\hb{35.12}{310}(i)) and could add solutions. The tiers of these rays, and their corners, are not yet computed.

\section{The two frameworks as one superset}

Brown's results hold on a subclass of the families, and nothing here contradicts them. What separates the two frameworks is a choice of letters, not a disagreement. This entry records:
\begin{itemize}
\item where his class sits inside ours;
\item what each side carries that the other does not;
\item what the union gave at once (\S14);
\item the gaps that the union points at.
\end{itemize}

\subsection{Brown's class inside the families}\status{derived; his definitions recalled}

A graph period without kinematics is the family of one list, the cycle list of the graph, with the Symanzik letter alone:
\[
\int_{\sigma_{\cE}}\frac{\Omega_{\cE}}{\Psi^2}\qquad\leftrightarrow\qquad(\cB;\ s=0,\ t=2),\quad \mvec=\one .
\]
With massless kinematics, his second Symanzik polynomial is the Symanzik letter of a weighted table: the truncation of the cycle list, whose bases are the spanning $2$-forests, weighted by the squared momentum through the cut,
\[
\Phi(q)=\sum_{\text{2-forests }F}q(F)^2\,\alpha^{\cE\setminus F}=\Psi_{(\mathrm T\cB,\ \theta_q)} .
\]
Its Kirchhoff counterpart $K_{(\mathrm T\cB,\theta_q)}=\sum_Fq(F)^2\alpha^F$ is the letter his presentations never write. With masses his $\Xi$ is not multilinear and is not a letter of a table.

\subsection{What each side carries}\status{recalled from his papers; the reference by printed number}

\emph{His, not in the reference:}
\begin{itemize}
\item the motive of the period and its relative cohomology, the coaction, weights, and face maps;
\item the affine model;
\item downstream of these: Galois-theoretic constraints on periods, single-valued periods, and the decomposition of the content into its letters of weight (alias: the $f$-alphabet).
\end{itemize}

\emph{Ours, not in his:}
\begin{itemize}
\item the serif family, in which every period is a final answer: the stacked Gamma, the Gramma and the Zamma, the Breta and the Trailmoment (Def.~\hb{1.1}{10}(vi)), closed under products and nested reads (Thm.~\hb{5.21}{36}); the interest of a family is a Gramma (Thm.~\hb{8.3}{54}), its period a stacked Gamma times a Trailmoment of the joint law of its letters (Prop.~\hb{8.11}{56}) and a fibre integral of stacked Gammas of shares (Cor.~\hb{17.15}{111}, Cor.~\hb{17.16}{112});
\item both letters, with independent exponents: the lean, the even letter $\mathbf E=K\Psi$ and the centre (Def.~\hb{8.17}{57}); Cremona exchanges the letters together with their exponents (Thm.~\hb{8.13}{57});
\item tables with weights and general lists, not only graphs;
\item a canonical decomposition: the Foreman bunch and tree (Def.~\hb{13.6}{78}, Def.~\hb{14.4}{81});
\item the levels and Skeleton values of a bunch at any width, its scale integral, what crosses it and the gluing across it (Def.~\hb{12.9}{71}, Thm.~\hb{12.17}{73}, Thm.~\hb{15.2}{87}, Thm.~\hb{20.13}{168}, Thm.~\hb{21.13}{178}), and at width one the cells, their hinge values and ratio, and the gluing across a bunch of width one (Thm.~\hb{18.6}{116}, Def.~\hb{18.4}{115}, Thm.~\hb{18.13}{119}). This is structure inside the open part, which his faces do not see.
\item continuation in the exponents and masses, with poles on the walls (Thm.~\hb{8.7}{55});
\item the letter and mark relations (Thm.~\hb{8.3}{54});
\item finite parts as constant terms (Thm.~\hb{1.6}{12}, Cor.~\hb{1.7}{12});
\item residues at walls as products of face periods (Thm.~\hb{16.23}{100});
\item explicit evaluation, by shares, closings and the one Trailmoment, and the arithmetic of rows and clearings.
\end{itemize}

\emph{The union, at once} (\S14): Claims~\ref{cl:weights}, \ref{cl:serif14}, \ref{cl:hinge14} and~\ref{cl:pinned14}. His coaction acts on our values, and the left factors it produces are again values of the serif family, attached to the tree. Neither side alone can state them.

\subsection{The gap his own text names}\status{not yet computed}

After his weight relations he writes that his small-graphs conjecture needs ``a notion of regularisation on the level of motivic periods, which is beyond the scope of the notes.'' The reference regularizes in the data, and the wall residues of Thm.~\hb{16.23}{100} are his faces read through the continuation:
\begin{gather*}
\operatorname*{Res}_{\kappa_S=0}P_F=P_{F^S}\,P_{F^R}\qquad\text{(Thm.~\hb{16.23}{100})},\\
\Reg\!\int_0^\infty f\,dx=\CT_{\varepsilon=0}\int_0^\infty x^{\varepsilon}f(x)\,dx\qquad\text{(Thm.~\hb{1.6}{12})}.
\end{gather*}
The letter and mark relations shift the data by integers. They write the continued period near a wall as a combination, rational in the data, of convergent periods at shifted data. So every Laurent coefficient at a wall, the finite part among them, is a rational combination of convergent periods and of their derivatives in the exponents. These are integrals with logarithms of letters inserted, and each is a period of a motive.

If the relations of Thm.~\hb{8.3}{54} hold between motivic periods, this is a regularization on the level of motivic periods. Those relations are integrations by parts, so what has to be checked is that their boundary terms vanish on the model. It would let the faces of Claim~\ref{cl:weights} be read through the walls rather than through the adjoined $\binom J1$, which is the form of his conjecture for families of lists.

\subsection{Two smaller gaps}\status{not yet computed}

\begin{enumerate}[label=(\arabic*)]
\item \emph{The coaction on the content.} The content $\mathcal C$ needs a coaction of its own. Its trails carry one by the cutting rule (\S13.3). To make $\mathfrak W/\mathfrak I$ (Def.~\hb{36.1}{315}) a comodule, the moves of Thm.~\hb{18.35}{133} and the relations of Def.~\hb{36.1}{315} must be shown to be relations between motivic periods, as Lemma~8.1 of the note does for one change of variables.
\item \emph{Width two and more: settled} (\S14.5). The width-one restriction came from reading the whole room through one hinge value (Claim~\ref{cl:onehinge14}); with one hinge value per cell the bunch has any width (Claim~\ref{cl:pinned14}), and a piece splits off exactly when it is pinned (Claim~\ref{cl:reach14}).
\end{enumerate}

\section{Not yet computed}

\begin{enumerate}[label=(\arabic*)]
\item The level read off the pencil (\S13.6--\S13.8): the equality is refuted at four marks; Claim~\ref{cl:criticalprime} proves that every critical prime is a prime of degeneration of the letters. Open: whether every prime of degeneration of the letters is a critical prime of $L$ (\S13.8(b)).
\item The coaction on the content (\S9.3) is written: \S13, with the verb of \S9.3 in \S13.2 and the deconcatenation in \S13.3; at tier two, \S13.5.
\item The interchange of \S9.4 in general: the top column sum of a hinge closing as a port at the wall. The four-mark constant-term test passes and is proved (\S10.2).
\item \S9.5: the tier hypothesis, first on the five-mark list, where the room along the hinge is the three-mark list and its values carry $\Tx{2}{(\Ilat{1+3\mathbb Z}-\Ilat{2+3\mathbb Z})}{0}$.
\item The two congruences of \S9.7 are proved (\S10.3). Open: the $2$-adic sharpness $64^{k-1}$ for $B_k$, and the same derivation on three marks, where the live input $(j!)^2/(2j+1)!$ of \S8.5 should be the half-integer pole structure of a critical-line word at the frequency $\frac13$.
\item The full power of the frame in the clearing: \S8.7 derives $|c|^{-k/h}$ and computes $|c|^{-k}$ on two rays. Needs: the denominators of the solutions of a recurrence with $\operatorname{lc}(p_h)=c_h$, i.e.\ that the divisions by $p_h(k)$ do not cancel in the power, only in the Gamma part (the $(k!)^3\to d_k^3$ of Thm.~\hb{34.9}{302} is such a cancellation).
\item The two recurrences of \S8.7 from the ray relations (Thm.~\hb{35.3}{308}) by hand, as relations in $\mathbb C[X_F]\otimes\mathbb C(k)$ modulo $\mathcal L$ (Def.~\hb{35.4}{308}).
\item The frequency $\frac13$ of \S8.5 from the family: the gauges $(1,\omega,\omega^2)$ at which $K$ vanishes, or the reflection of the live input's Laplace side (Thms.~\hb{5.18}{35}, \hb{5.19}{35}); the general statement for $\binom{\cE}1$ on $n$ marks, where at $n=5$ the recalled value is a modular $L$-value and no Dirichlet frequency can be the whole story.
\item The law of the free shares of a general atom on its wall contour (Cor.~\hb{17.15}{111}) with its Cremona-even and Cremona-odd parts, of which Def.~\hb{33.1}{297} and Thm.~\hb{33.2}{297} are the two-mark corner.
\item The gap on a (family, direction) pair as \S8.8 states it, with the full power supplied by item (6); then the search: rays of the two-mark family and of the lists on three and four marks, with insertions (Beta ports, Thm.~\hb{34.13}{303}) chosen to make the native value a unit, the native value read from the stride before any closing.
\item \S7.5: the invariants of $S_6$ on the weight-$5$ columns of the centre closing of $\binom{\cE}1$ on six marks.
\item A general reason, or a counterexample, for \S5.1: are the atom-formula pins of every list unimodular?
\item Whether the lists $\binom J1$ can be dropped from the pieces of \S14: for the model without them, are the left factors of weight $\le k$ still generated by the faces? This is the question Brown leaves open for graphs.
\item (Settled in \S14.5: a piece splits off exactly when it is pinned, Claim~\ref{cl:reach14}; a bunch of width $c\ge2$ sits in the core, Claim~\ref{cl:pinned14}.)
\item Whether the letter and mark relations (Thm.~\hb{8.3}{54}) hold between motivic periods, making the continuation in the data a regularization of motivic periods (\S15.3).
\end{enumerate}

\end{document}
